\xchapter{The Westminster Confession of Faith, A.D. 1647}

[The English text is taken from the second edition which appeared under the title, '\emph{The Humble} | \emph{Advice} | \emph{of the} | \emph{Assembly} | \emph{of} | \emph{Divines,} | \emph{now by Authority of Parliament} | \emph{sitting at Westminster,} | \emph{concerning }| \emph{a Confession of Faith:} | \emph{with the Quotations and Texts of} | \emph{Scripture annexed. Presented by them lately to both Houses of Parliament.} | \emph{Printed at London;} | \emph{and} | \emph{reprinted at Edinburgh by Evan Tyler, Printer to} | \emph{the Kings most Excellent Majestie.} 1647.' The spelling and punctuation are conformed to modern usage.

The changes of the American revision, which occur chiefly in Ch. XXIII., relating to the Civil Magistrate, and in Ch. XXXI., relating to Synods and Councils, are inserted in their proper places, and marked by italics. Minor changes are indicated in footnotes.

The Latin translation of the Westminster Confession and Catechisms by G. D. (see Preface) appeared first at Cambridge, 1656 (also 1659); at Edinburgh, 1694, etc.; and at Glasgow, 1660), under the title, 'C\textsc{onfessio} F\textsc{idei} | \emph{in Conventu theologorum authoritate | Parliamenti Anglicani indicto | Elaborata; | eidem Parliamento postmodum | Exhibita; | Quin et ab eodem, deindeque ab Ecclesia Scoticana | Cognita et Approbata; | und cum} | C\textsc{atechismo} | \emph{duplici}, M\textsc{ajori}, M\textsc{inorique}; | \emph{E Sermone Anglicano summa cum fide | in Litinum versa. Cantabrigiæ: excudebat Johannes Field, celeberrimæ Academiæ typographus}.']

Confession of Faith.

Confessio Fidei.

\xsection{Chapter I. Of the Holy Scripture}

I. Although the light of nature, and the works of creation and providence, do so far manifest the goodness, wisdom, and power of God, as to leave men inexcusable;\footnote{Rom. ii. 14, 15; i. 19, 20;  Psa. xix. 1–3; Rom. i. 32;  ii. 1.} yet are they not sufficient to give that knowledge of God, and of his will, which is necessary unto salvation;\footnote{1 Cor. i. 21; ii. 13, 14.} therefore it pleased the Lord, at sundry times, and in divers manners, to reveal himself, and to declare that his will unto his Church;\footnote{Hebrews 1:1.} and afterwards, for the better preserving and propagating of the truth, and for the more sure establishment and comfort of the Church against the corruption of the flesh, and the malice

I. \emph{Quanquam naturæ lumen, operaque Dei cum Creationis tum Providentiæ, bonitatem ejus, sapientiam, potentiamque eo usque manifestant, ut homines vel inde reddantur inexcusabiles:}\footnote{Rom. ii. 14, 15;  i. 19, 20; Psa. xix. 1–3;  Rom. i. 32; ii. 1.}\emph{ eam tamen Dei, voluntatisque divinæ cognitionem, quæ porro est ad salutem necessaria, nequeunt nobis ingenerare.}\footnote{1 Cor. i. 21; ii. 13, 14.}\emph{ Quocirca Domino complacitum est, variis quidem modis vicibusque Ecclesiæ suæ semetipsum revelare, suamque hanc voluntatem patefacere;}\footnote{Hebrews 1:1.}\emph{ sed et eandem omnem postea literis consignare, quo et veritati suæ tam conservandæ quam propagandæ melius consuleret, nec Ecclesia sua contra carnis corruptelam, contra malitiam mundi}

of Satan and of the world, to commit the same wholly unto writing;\footnote{Prov. xxii. 19–21; Luke i. 3, 4;  Rom. xv. 4; Matt. iv. 4, 7, 10; Isa. viii. 19, 20.} which maketh the holy Scripture to be most necessary;\footnote{2 Tim. iii. 15; 2 Pet. i. 19.} those former ways of God's revealing his will unto his people being now ceased.\footnote{Heb. i. 1, 2.}

\emph{Satanæque, præsidio foret ac solatio destituta.}\footnote{Prov. xxii. 19–21; Luke i. 3, 4;  Rom. xv. 4; Matt. iv. 4, 7, 10; Isa. viii. 19, 20.}\emph{ Unde factum est, ut, postquam pristini illi modi, quibus olim populo suo Deus voluntatem suam revelabat, jam desiverint,}\footnote{2 Tim. iii. 15; 2 Pet. i. 19.}\emph{ Scriptura Sacra sit maxime necessaria.}\footnote{Heb. i. 1, 2.}

II. Under the name of holy Scripture, or the Word of God written, are now contained all the Books of the Old and New Testament, which are these:

II. \emph{Sacræ Scripturæ, nomine, seu Verbi Dei scripti continentur hodie omnes illi libri tam Veteris quam Novi Instrumenti,}\footnote{[So the Cambridge eds. of 1656 and 1659. The Edinb. ed. reads \emph{Testamenti}.]} \emph{nempe quorum inferius subsequuntur nomina.}

\emph{Of the Old Testament.}

Veteris Testamenti.

Genesis.

Ecclesiastes.

\emph{Genesis.}

\emph{Ecclesiastes.}

Exodus.

The Song of Songs.

\emph{Exodus.}

\emph{Canticum Canticorum.}

Leviticus.

Isaiah.

\emph{Leviticus.}

\emph{Isaias.}

Numbers.

Jeremiah.

\emph{Numeri.}

\emph{Jeremias.}

Deuteronomy.

Lamentations.

\emph{Deuteronomium.}

\emph{Lamentationes.}

Joshua.

Ezekiel.

\emph{Josua.}

\emph{Ezechiel.}

Judges.

Daniel.

\emph{Judices.}

\emph{Daniel.}

Ruth.

Hosea.

\emph{Ruth.}

\emph{Hosea.}

I. Samuel.

Joel.

\emph{Samuelis} 1.

\emph{Joel.}

II. Samuel.

Amos.

\emph{Samuelis} 2.

Amos.

I. Kings.

Obadiah.

\emph{Regum} 1.

\emph{Obadias.}

II. Kings.

Jonah.

\emph{Regum} 2.

\emph{Jonas.}

I. Chronicles.

Micah.

\emph{Chronicorum} 1.

\emph{Micheas.}

II. Chronicles.

Nahum.

\emph{Chronicorum} 2.

\emph{Nahum.}

Ezra.

Habakkuk.

\emph{Ezra.}

\emph{Habucuc.}

Nehemiah.

Zephaniah.

\emph{Nehemias.}

\emph{Zephanias.}

Esther.

Haggai.

\emph{Esther.}

\emph{Haggæus.}

Job.

Zechariah.

\emph{Job.}

\emph{Zacharias.}

Psalms.

Malachi.

\emph{Psalmi.}

\emph{Malachias.}

Proverbs.

\emph{Proverbia.}

\emph{Of the New Testament.}

Novi autem.

The Gospels according to

\emph{Evangelium secundem}

Matthew,

Luke,

\emph{Matthæum},

\emph{Lucam},

Mark,

John.

\emph{Marcum},

\emph{Johannem}.

The Acts of the

To Timothy II.

\emph{Acta apostolorum.}

Apostles.

To Titus.

\emph{Titum.}

Paul's Epistles to the

To Philemon.

\emph{Pauli espistolæ ad}

\emph{Philemonem.}

Romans.

The Epistle to the

\emph{Romanos.}

\emph{Epist. ad Hebræos.}

Corinthians I.

Hebrews.

\emph{Corinthios I. II.}

Corinthians II.

The Epistle of James.

\emph{Jacobi Epistola.}

Galatians.

The First and Second

\emph{Galatas.}

\emph{Petri Epist. I. II.}

Ephesians.

Epistles of Peter.

\emph{Ephesios.}

Philippians.

The First, Second, and

\emph{Philippenses.}

\emph{Johan. Epist. I. II.}

Colossians.

Third Epistles of

\emph{Collossenses.}

\emph{III.}

Thessalonians I.

John.

\emph{Thessalonicens I. II.}

Thessalonians II.

The Epistle of Jude.

\emph{Judæ Epistola.}

To Timothy I.

The Revelation.

\emph{Timotheum I. II.}

\emph{Apocalypsis.}

All which are given by inspiration of God, to be the rule of faith and life.\footnote{Luke xvi. 29, 31; Eph. ii. 20;  Rev. xxii. 18, 19; 2 Tim. iii. 16.}

\emph{Qui omnes divina inspiratione dati sunt in Fidei vitæque regulam.}\footnote{Luke xvi. 29, 31; Eph. ii. 20;  Rev. xxii. 18, 19; 2 Tim. iii. 16.}

III. The books commonly called Apocrypha, not being of divine inspiration, are no part of the Canon of the Scripture; and therefore are of no authority in the Church of God, nor to be any otherwise approved, or made use of, than other human writings.\footnote{Luke xxiv. 27, 44; Rom. iii. 2; 2 Pet. i. 21.}

III. \emph{Libri Apocryphi, vulgo dicti, quum non fuerint divinitus inspirati, Canonem Scripturæ nullatenus constituunt; proindeque nullam aliam authoritatem obtinere debent in Ecclesia Dei, nec aliter quam alia humana scripta, sunt aut approbandi aut adhibendi}.\footnote{Luke xxiv. 27, 44; Rom. iii. 2; 2 Pet. i. 21.}

IV. The authority of the holy Scripture, for which it ought to be believed and obeyed, dependeth not upon the testimony of any man or church, but wholly upon God (who is truth itself), the Author thereof; and therefore it is to be received, because it is the Word of God.\footnote{2 Pet. i. 19, 21; 2 Tim. iii. 16; 1 John v. 9;  1 Thess. ii. 13.}

IV. \emph{Authoritas Scripturæ sacræ propter quam ei debetur fides et observantia, non ab ullius aut hominis aut Ecclesiæ pendet testimonio, sed a solo ejus authore Deo, qui est ipsa veritas: eoque est a nobis recipienda, quoniam est Verbum Dei}.\footnote{2 Pet. i. 19, 21; 2 Tim. iii. 16; 1 John v. 9;  1 Thess. ii. 13.}

V. We may be moved and induced by the testimony of the

V. \emph{Testimonium Ecclesiæ efficere quidem potest ut de Scriptura sacra }

Church to an high and reverent esteem of\footnote{[Am. ed. \emph{ for.}]} the holy Scripture;\footnote{1 Tim. iii. 15.} and the heavenliness of the matter, the efficacy of the doctrine, the majesty of the style, the consent of all the parts, the scope of the whole (which is to give all glory to God), the full discovery it makes of the only way of man's salvation, the many other incomparable excellencies, and the entire perfection thereof, are arguments whereby it doth abundantly evidence itself to be the Word of God; yet, notwithstanding, our full persuasion and assurance of the infallible truth, and divine authority thereof, is from the inward work of the Holy Spirit, bearing witness by and with the Word in our hearts.\footnote{1 John ii. 20, 27;  John xvi. 13, 14; 1 Cor. ii. 10–12; Isa. lix. 21.}

\emph{quam honorifice sentiamus;}\footnote{1 Tim. iii. 15.}\emph{ materies insuper ejus cœlestis, doctrinæ vis et efficacia, styli majestas, partium omnium consensus, totiusque scopus} (\emph{ut Deo nempe omnis gloria tribuatur}), \emph{plena denique quam exhibet unicæ ad salutem viæ commonstratio, præter alias ejus virtutes incomparabiles, et perfectionem summam, argumenta sunt quibus abunde se Verbum Dei et luculenter probat; nihilominus tamen plena persuasio et certitudo de ejus tam infallibili veritate, quam authoritate divina non aliunde nascitur quam ab interna operatione Spiritus Sancti, per verbum et cum verbo ipso in cordibus nostris testificantis}.\footnote{1 John ii. 20, 27; John xvi. 13, 14; 1 Cor. ii. 10–12;  Isa. lix. 21.}

VI. The whole counsel of God, concerning all things necessary for his own glory, man's salvation, faith, and life, is either expressly set down in Scripture, or by good and necessary consequence may be deduced from Scripture: unto which nothing at any time is to be added, whether by new revelations of the Spirit, or traditions of men.\footnote{2 Tim. iii. 15–17; Gal. i. 8, 9;  2 Thess. ii. 2.} Nevertheless we acknowledge the inward illumination of the Spirit of God to be necessary for the saving

VI. \emph{Consilium Dei universum de omnibus quæ ad suam ipsius gloriam, quæque ad hominum salutem, fidem, vitamque sunt necessaria, aut expresse in Scriptura continetur, aut consequentia bona et necessaria derivari potest a Scriptura; cui nihil deinceps addendum est, seu novis a spiritu revelationibus, sive traditionibus hominum.}\footnote{2 Tim. iii. 15–17; Gal. i. 8, 9;  2 Thess. ii. 2.}\emph{ Internam nihilominus illuminationem Spiritus Dei ad salutarem eorum perceptionem, quæ in Verbo Dei }

understanding of such things as are revealed in the Word;\footnote{John vi. 45;  1 Cor. ii. 9, 10, 12.} and that there are some circumstances concerning the worship of God, and government of the Church, common to human actions and societies, which are to be ordered by the light of nature and Christian prudence, according to the general rules of the Word, which are always to be observed.\footnote{1 Cor. xi. 13, 14; xiv. 26, 40.}

\emph{revelantur, agnoscimus esse necessariam:}\footnote{John vi. 45; 1 Cor. ii. 9, 10, 12.} quin etiam nonnullas esse circumstantias cultum Dei spectantes et Ecclesiæ regimen, iis cum humanis actionibus et societatibus communes, quæ naturali lumine ac prudentia Christiana secundum generales verbi regulas (perpetuo quidem illas observandas) sunt regulandæ.\footnote{1 Cor. xi. 13, 14; xiv. 26, 40.}

VII. All things in Scripture are not alike plain in themselves, nor alike clear unto all;\footnote{2 Pet. iii. 16.} yet those things which are necessary to be known, believed, and observed, for salvation, are so clearly propounded and opened in some place of Scripture or other, that not only the learned, but the unlearned, in a due use of the ordinary means, may attain unto a sufficient understanding of them.\footnote{Psa. cxix. 105, 130.}

VII. \emph{Quæ in Scriptura continentur non sunt omnia æque aut in se perspicua, aut omnibus hominibus evidentia,}\footnote{2 Pet. iii. 16.}\emph{ ea tamen omnia quæ ad salutem necessaria sunt cognitu, creditu, observatu, adeo perspicue, alicubi saltem in Scriptura, proponuntur et explicantur, ut eorum non docti solum, verum indocti etiam ordinariorum debito usu mediorum, sufficientem assequi possint intelligentiam.}\footnote{Psa. cxix. 105, 130.}

VIII. The Old Testament in Hebrew (which was the native language of the people of God of old), and the New Testament in Greek (which at the time of the writing of it was most generally known to the nations), being immediately inspired by God, and by his singular care and providence kept pure in all ages, are therefore authentical;\footnote{Matt. v. 18.} so as in all controversies

VIII. \emph{Instrumentum Vetus Hebræa lingua (antiqua Dei populo nempe vernacula) Novum autem Græca} (\emph{ut quæ apud Gentes maxime omnium tunc temporis, quum scriberetur illud, obtinuerat}), \emph{immediate a Deo inspirata, ejusque cura et Providentia singulari per omnia huc usque secula pura et intaminata custodita, ea propter sunt authentica.}\footnote{Matt. v. 18.} \emph{Adeo sane ut ad }

of religion the Church is finally to appeal unto them.\footnote{Isa. viii. 20; Acts xv. 15;  John v. 39, 46.} But because these original tongues are not known to all the people of God who have right unto, and interest in the Scriptures, and are commanded, in the fear of God, to read and search them,\footnote{John v. 39.} therefore they are to be translated into the vulgar language of every nation unto which they come,\footnote{1 Cor. xiv. 6, 9,11, 12, 24, 27, 28.} that the Word of God dwelling plentifully in all, they may worship him in an acceptable manner,\footnote{Col. iii. 16.} and, through patience and comfort of the Scriptures, may have hope.\footnote{Rom. xv. 4.}

\emph{illa ultimo in omnibus de religione controversiis Ecclesia debeat appellare.}\footnote{Isa. viii. 20;  Acts xv. 15; John v. 39, 46.} \emph{Quoniam autem Originales istæ linguæ non sunt toti Dei populo intellectæ} (\emph{Quorum tamen et jus est ut scripturas habeant, et interest plurimum, quique eas in timore Dei legere jubentur et perscrutari})\footnote{John v. 39.}\emph{ proinde sunt in vulgarem cujusque Gentis, ad quam pervenerint linguam transferendæ,}\footnote{1 Cor. xiv. 6, 9,11, 12, 24, 27, 28.}\emph{ ut omnes, verbo Dei opulenter in ipsis habitante, Deum grato acceptoque modo colant,}\footnote{Col. iii. 16.}\emph{ et per patientiam ac consolationem Scripturarum spem habeant.}\footnote{Rom. xv.4.}

IX. The infallible rule of interpretation of Scripture is the Scripture itself; and therefore, when there is a question about the true and full sense of any Scripture (which is not manifold, but one), it must\footnote{[Am. ed. may.]} be searched and known by other places that speak more clearly.\footnote{2 Pet. i. 20, 21; Acts xv. 15; [Am. John v. 46.]}

IX. \emph{Infallibilis Scripturam interpretandi regula est Scriptura ipsa. Quoties igitur cunque oritur quæstio de. vero plenoque Scripturæ cujusvis sensu} (\emph{unicus ille est non multiplex}), \emph{ex aliis locis, qui apertius loquuntur, est indagandus et cognoscendus.}\footnote{2 Pet. i. 20, 21; Acts xv. 15; [Am.  John v. 46.]}

X. The Supreme Judge, by which all controversies of religion are to be determined, and all decrees of councils, opinions of ancient writers, doctrines of men, and private spirits, are to be examined, and in whose sentence we are to rest, can

X. \emph{Supremus judex, a quo omnes de religione controversiæ sunt determinandæ, omnia Conciliorum decreta, opiniones Scriptorum Veterum, doctrinæ denique hominum, et privati quicunque Spiritus sunt examinandi, cujusque sententia tenemur}

be no other but the Holy Spirit speaking in the Scripture.\footnote{Matt. xxii. 29, 31;  Eph. ii. 20; Acts xxviii. 25.}

\emph{nemur acquiescere, nullus alius esse potest, præter Spiritum Sanctum in Scriptura pronunciantem.}\footnote{Matt. xxii. 29, 31; Eph. ii. 20; Acts xxviii. 25.}

\xsection{Chapter II. Of God, and of the Holy Trinity}

I. There is but one only\footnote{Deut. vi. 4;  1 Cor. viii. 4, 6.} living and true God,\footnote{1 Thess. i. 9; Jer. x. 10.} who is infinite in being and perfection,\footnote{Job xi. 7, 8, 9; xxvi. 14.} a most pure spirit,\footnote{John iv. 24.} invisible,\footnote{1 Tim. i. 17.} without body, parts,\footnote{Deut. iv. 15, 16; John iv. 24; Luke xxiv. 39.} or passions,\footnote{Acts xiv. 11, 15.} immutable,\footnote{James i. 17;  Mal. iii. 6.} immense,\footnote{1 Kings viii. 27;  Jer. xxiii. 23, 24.} eternal,\footnote{Psa. xc. 2;  1 Tim. i. 17.} incomprehensible,\footnote{Psa. cxlv. 3.} almighty,\footnote{Gen. xvii. 1;  Rev. iv. 8.} most wise,\footnote{Rom. xvi. 27.} most holy,\footnote{Isa. vi. 3; Rev. iv. 8.} most free,\footnote{Psa. cxv. 3.} most absolute,\footnote{Exod. iii. 14.} working all things according to the counsel of his own immutable and most righteous will,\footnote{Eph. i. 11.} for his own glory;\footnote{Prov. xvi. 4; Rom. xi. 36; [Am. ed. Rev. iv. 11] .} most loving,\footnote{1 John iv. 8, 16.} gracious, merciful, longsuffering, abundant in goodness and truth, forgiving iniquity, transgression, and sin;\footnote{Exod. xxxiv. 6, 7.} the rewarder of them that diligently seek him;\footnote{Heb. xi. 6.} and withal most just and terrible in his judgments;\footnote{Neh. ix. 32, 33.}

I. \emph{Unus est unicusque},\footnote{Deut. vi. 4; 1 Cor. viii. 4, 6.}\emph{ vivens ille et verus Deus:}\footnote{1 Thess. i. 9;  Jer. x. 10.}\emph{ qui idem est essentia et perfectione infinitus,}\footnote{Job xi. 7, 8, 9;  xxvi. 14.}\emph{ Spiritus purissimus,}\footnote{John iv. 24.}\emph{ invisibilis,}\footnote{1 Tim. i. 17.}\emph{ sine corpore, sine partibus,}\footnote{Deut. iv. 15, 16; John iv. 24; Luke xxiv. 39.}\emph{ sine passionibus,}\footnote{Acts xiv. 11, 15.}\emph{ immutabilis,}\footnote{James i. 17; Mal. iii. 6.}\emph{ immensus,}\footnote{1 Kings viii. 27; Jer. xxiii. 23, 24.} æternus,\footnote{Psa. xc. 2;  1 Tim. i. 17.}\emph{ incomprehensibilis,}\footnote{Psa. cxlv. 3.}\emph{ omnipotens,}\footnote{Gen. xvii. 1;  Rev. iv. 8.}\emph{ summe sapiens,}\footnote{Rom. xvi. 27.}\emph{ summe sanctus,}\footnote{Isa. vi. 3;  Rev. iv. 8.}\emph{ liberrimus,}\footnote{Psa. cxv. 3.}\emph{ maxime absolutus;}\footnote{Exod. iii. 14.}\emph{ operans omnia secundum consilium immutabilis suæ ac justissimæ voluntatis,}\footnote{Eph. i. 11.}\emph{ ad suam ipsius gloriam;}\footnote{Prov. xvi. 4; Rom. xi. 36; [Am. ed. Rev. iv. 11].}\emph{ idemque summa benignitate,}\footnote{1 John iv. 8, 16.} \emph{gratia, misericordia, et longanimitate; bonitate abundans et veritate; condonans iniquitatem, transgressionem et peccatum;}\footnote{Exod. xxxiv. 6, 7.}\emph{ studiose quærentium ipsum remunerator;}\footnote{Heb. xi. 6.}\emph{ sed et in judiciis suis justissimus idem ac tremendus maxime;}\footnote{Neh. ix. 32, 33.}

hating all sin,\footnote{Psa. v. 5, 6.} and who will by no means clear the guilty. .\footnote{Nahum i. 2, 3; Exod. xxxiv. 7.}

\emph{peccatum omne perosus,}\footnote{Psa. v. 5, 6.}\emph{ et qui sontem nullo unquam absolvet modo.}\footnote{Nahum i. 2, 3; Exod. xxxiv. 7.}

II. God hath all life,\footnote{John v. 26.} glory,\footnote{Acts vii. 2.} goodness,\footnote{Psa. cxix. 68.} blessedness,\footnote{1 Tim. vi. 15;  Rom. ix. 5.} in and of himself; and is alone in and unto himself allsufficient, not standing in need of any creatures which he hath made,\footnote{Acts xvii. 24, 25.} nor deriving any glory from them,\footnote{Job xxii. 2, 23.} but only manifesting his own glory in, by, unto, and upon them: he is the alone foundation of all being, of whom, through whom, and to whom are all things;\footnote{Rom. xi. 36.} and hath most sovereign dominion over them, to do by them, for them, or upon them whatsoever himself pleaseth.\footnote{Rev. iv. 11;  1 Tim. vi. 15; Dan. iv. 25, 35.} In his sight all things are open and manifest;\footnote{Heb. iv. 13.} his knowledge is infinite, infallible, and independent upon the creature;\footnote{Rom. xi. 33, 34; Psa. cxlvii. 5.} so as nothing is to him contingent or uncertain.\footnote{Acts 15:18; Ezeck. xi. 5.} He is most holy in all his counsels, in all his works, and in all his commands.\footnote{Psa. cxlv. 17;  Rom. vii. 12.} To him is due from angels and men, and every other creature, whatsoever worship, service, or obedience, he is pleased to require of them.\footnote{Rev. v. 12–14.}

II. \emph{Omnem vitam,}\footnote{John v. 26.}\emph{ omnem gloriam,}\footnote{Acts vii. 2.}\emph{ bonitatem,}\footnote{Psa. cxix. 68.}\emph{ beatitudinemque}\footnote{1 Tim. vi. 15;  Rom. ix. 5.}\emph{ omnem in sese habet et a seipso Deus; qui solus in se sibique est ad omnia sufficiens; creaturarum, quas ipse condidit, nullius egens,}\footnote{Acts xvii. 24, 25.}\emph{ nec gloriam ab eis derivans ullam,}\footnote{Job xxii. 2, 23.}\emph{ verum in iis, per eas, iis ipsis, ac super eas propriam ipsius gloriam tantummodo manifestans. Is omnis entitatis fans est unicus, a quo, per quem et ad quem omnia;}\footnote{Rom. xi. 36.}\emph{ summumque in ea dominium habet, ac per illa, pro illis, in illa pro suo arbitrio quidlibet agendi potestatem.}\footnote{Rev. iv. 11;  1 Tim. vi. 15; Dan. iv. 25, 35.}\emph{ In conspectu ejus aperta sunt omnia ac manifesta;}\footnote{Heb. iv. 13.}\emph{ scientia ejus infinita est, infallibilis, atque a creatura independens,}\footnote{Rom. xi. 33, 34; Psa. cxlvii. 5.}\emph{ adeo ut illi contingens incertumve nihil sit;}\footnote{Acts 15:18; Ezeck. xi. 5.}\emph{ in omnibus ejus consiliis, operibus et mandatis est sanctissimus.}\footnote{Psa. cxlv. 17;  Rom. vii. 12.}\emph{ Quicquid cultus, quicquid officii, quicquid obsequii ab Angelis illi, ab hominibus, aut a quavis creatura exigere placet, id illi omne jure optimo debetur.}\footnote{Rev. v. 12–14.}

III. In the unity of the Godhead

III. In Deitatis unitate personæ

head there be three persons, of one substance, power, and eternity: God the Father, God the Son, and God the Holy Ghost.\footnote{1 John v. 7; Matt. iii. 16, 17;  xxviii. 19; 2 Cor. xiii. 14.} The Father is of none, neither begotten nor proceeding; the Son is eternally begotten of the Father;\footnote{John i. 14, 18.} the Holy Ghost eternally proceeding from the Father and the Son.\footnote{John xv. 26; Gal. iv. 6.}

\emph{tres sunt unius ejusdemque essentiæ, potential ac æternitatis; Deus Pater, Deus Filius, ac Deus Spiritus Sanctus.}\footnote{1 John v. 7; Matt. iii. 16, 17;  xxviii. 19; 2 Cor. xiii. 14.}\emph{ Pater quidem a nullo est, nec genitus nempe nec procedens: Filius autem a Patre est æterne genitus:}\footnote{John i. 14, 18.}\emph{ Spiritus autem Sanctus æterne procedens a Patre Filioque.}\footnote{John xv. 26; Gal. iv. 6.}

\xsection{Chapter III. Of God's Eternal Decree}\footnote{[Am. ed. \emph{ decrees}.]}

I. God from all eternity did, by the most wise and holy counsel of his own will, freely and unchangeably ordain whatsoever comes to pass;\footnote{Eph. i. 11; Rom. xi. 33;  Heb. vi. 17; Rom. ix. 15,18.} yet so as thereby neither is God the author of sin,\footnote{James i. 13,17; 1 John i. 5; [Am. ed.  Eccl. vii. 29].} nor is violence offered to the will of the creatures, nor is the liberty or contingency of second causes taken away, but rather established.\footnote{Acts ii. 23;  Matt. xvii. 12; Acts iv. 27, 28;  John xix. 11; Prov. xvi. 33.}

I. \emph{Deus, e sapientissimo sanctissimoque consilio voluntatis suæ, libere ac immutabiliter, quicquid unquam evenit, ab omni æterno ordinavit;}\footnote{Eph. i. 11;  Rom. xi. 33; Heb. vi. 17;  Rom. ix. 15,18.}\emph{ ita tamen, id inde nec author peccati evadat Deus,}\footnote{James i. 13,17;  1 John i. 5; [Am. ed. Eccl. vii. 29].}\emph{ nec voluntati creaturarum sit vis illata, neque libertas aut contingentia causarum secundarum ablata sit, verum potius stabilita.}\footnote{Acts ii. 23; Matt. xvii. 12;  Acts iv. 27, 28; John xix. 11;  Prov. xvi. 33.}

II. Although God knows whatsoever may or can come to pass upon all supposed conditions,\footnote{Acts xv. 18;  1 Sam. xxiii. 11, 12; Matt. xi. 21, 23.} yet hath he not decreed any thing because he foresaw it as future, or as that which would come to pass upon such conditions.\footnote{Rom. ix. 11,13,16,18.}

II. \emph{Quamvis omnia cognoscat Deus, quæ suppositis quibusvis conditionibus sunt eventu possibilia;}\footnote{Acts xv. 18; 1 Sam. xxiii. 11, 12; Matt. xi. 21, 23.}\emph{ non tamen ideo quicquam decrevit quoniam illud præviderat aut futurum, aut positis talibus conditionibus eventurum.}\footnote{Rom. ix. 11,13,16,18.}

III. By the decree of God, for the manifestation of his glory,

III. \emph{Deus, quo gloriam suam manifestaret, nonnullos hominum}

some men and angels\footnote{1 Tim. v. 21;  Matt. xxv. 41.} are predestinated unto everlasting life, and others foreordained to everlasting death.\footnote{Rom. ix. 22, 23; Eph. i. 5, 6;  Prov. xvi. 4.}

\emph{ac Angelorum}\footnote{1 Tim. v. 21; Matt. xxv. 41.}\emph{ decreto suo ad æternam vitam prædestinavit, alios autem ad mortem æternam præordinavit.}\footnote{Rom. ix. 22, 23; Eph. i. 5, 6;  Prov. xvi. 4.}

IV. These angels and men, thus predestinated and foreordained, are particularly and unchangeably designed; and their number is so certain and definite that it can not be either increased or diminished.\footnote{2 Tim. ii. 19; John xiii. 18.}

IV. \emph{prædestinati illi et præordinati homines Angelique, particulariter sunt ac immutabiliter designati, certusque illorum est ac definitus numerus, adeo ut nec augeri possit nec imminui.}\footnote{2 Tim. ii. 19; John xiii. 18.}

V. Those of mankind that are predestinated unto life, God, before the foundation of the world was laid, according to his eternal and immutable purpose, and the secret counsel and good pleasure of his will, hath chosen in Christ, unto everlasting glory,\footnote{Eph. i. 4, 9, 11; Rom. viii. 30; 2 Tim. i. 9; 1 Thess. v. 9.} out of his mere free grace and love, without any foresight of faith or good works, or perseverance in either of them, or any other thing in the creature, as conditions, or causes moving him thereunto;\footnote{Rom. ix. 11, 13, 16; Eph. i, 4, 9.} and all to the praise of his glorious grace.\footnote{Eph. i. 6, 12.}

V. \emph{Qui ex humano genere sunt ad vitam prædestinati, illos Deus ante jacta mundi fundamenta, secundum æternum suum ac immutabile propositum, secretumque voluntatis suæ consilium et beneplacitum, elegit in Christo ad æternam gloriam,}\footnote{Eph. i. 4, 9, 11; Rom. viii. 30; 2 Tim. i. 9; 1 Thess. v. 9.}\emph{ idque ex amore suo et gratia mere gratuita; nec fide, nec bonis operibus, nec in his illave perseverantia, sed neque ulla alia re in creatura, prævisis, ipsum tanquam causis aut conditionibus ad id moventibus;}\footnote{Rom. ix. 11, 13, 16; Eph. i, 4, 9.}\emph{ quo totum nempe in laudem cederet gloriosæ suæ gratiæ.}\footnote{Eph. i. 6, 12.}

VI. As God hath appointed the elect unto glory, so hath he, by the eternal and most free purpose of his will, foreordained all the means thereunto.\footnote{1 Pet. i. 2;  Eph. i. 4, 5; ii. 10;  2 Thess. ii. 13.} Wherefore they who are elected, being fallen in

VI. \emph{Quemadmodum autem Deus electos ad gloriam destinavit, sic omnia etiam quibus illam consequantur media præordinavit, voluntatis suæ proposito æterno simul et liberrimo.}\footnote{1 Pet. i. 2; Eph. i. 4, 5;  ii. 10; 2 Thess. ii. 13.}\emph{ Quapropter electi, postquam }

Adam, are redeemed by Christ,\footnote{1 Thess. v. 9, 10; Tit. ii. 14.} are effectually called unto faith in Christ by his Spirit working in due season; are justified, adopted, sanctified,\footnote{Rom. viii. 30; Eph. i. 5; 2 Thess. ii. 13.} and kept by his power through faith unto salvation.\footnote{1 Pet. i. 5.} Neither are any other redeemed by Christ, effectually called, justified, adopted, sanctified, and saved, but the elect only.\footnote{John xvii. 9; Rom. viii. 28 to the end; John vi. 64, 65; viii. 47; x. 26; 1 John ii. 19.}

\emph{lapsi essent in Adamo, a Christo sunt redempti;}\footnote{1 Thess. v. 9, 10; Tit. ii. 14.}\emph{ per Spiritum ejus opportuno tempore operantem, ad fidem in Christum vocantur efficaciter; justificantur, sanctificantur,}\footnote{Rom. viii. 30; Eph. i. 5; 2 Thess. ii. 13.}\emph{ et potentia ipsius per fidem custodiuntur ad salutem.}\footnote{1 Pet. i. 5.}\emph{ Nec alii quivis a Christo redimuntur, vocantur efficaciter justificantur, adoptantur, sanctificantur et salvantur, præter electos solos.}\footnote{John xvii. 9; Rom. viii. 28 to the end; John vi. 64, 65; viii. 47; x. 26; 1 John ii. 19.}

VII. The rest of mankind God was pleased, according to the unsearchable counsel of his own will, whereby he extendeth or withholdeth mercy as he pleaseth, for the glory of his sovereign power over his creatures, to pass by, and to ordain them to dishonor and wrath for their sin, to the praise of his glorious justice.\footnote{Matt. xi. 25, 26; Rom. ix. 17, 18, 21, 22; 2 Tim. ii. 19, 20; Jude 4; 1 Pet. ii. 8.}

VII. \emph{Reliquos humani generis Deo placuit secundum consilium voluntatis suæ inscrutabile} \{\emph{quo misericordiam pro libitu exhibet abstinetve}) \emph{in gloriam supremæ suæ in creaturas potestatis, præterire; eosque ordinare ad ignominiam et iram pro peccatis suis, ad laudem justitiæ suæ gloriosæ.}\footnote{Matt. xi. 25, 26; Rom. ix. 17, 18, 21, 22; 2 Tim. ii. 19, 20; Jude 4; 1 Pet. ii. 8.}

VIII. The doctrine of this high mystery of predestination is to be handled with special prudence and care,\footnote{Rom. ix. 20; xi. 33; Deut. xxix. 29.} that men attending the will of God revealed in his Word, and yielding obedience thereunto, may, from the certainty of their effectual vocation, be assured of their eternal election.\footnote{2 Pet. i. 10.} So shall this doctrine afford matter of praise, reverence, and admiration of God;\footnote{Eph. i. 6; Rom. xi. 33.} and of humility, diligence, and abundant consolation

VIII. \emph{Doctrina de sublimi hoc prædestinationis mysterio non sine summa cura et prudentia tractari debet,}\footnote{Rom. ix. 20; xi. 33; Deut. xxix. 29.}\emph{ quo nimirum homines, dum voluntati Dei in verbo ejus revelatæ advertant animos, eique debitam exhibeant obedientiam, de efficaci sua vocatione certiores facti, ad æternæ suæ electionis assurgere possint certitudinem.}\footnote{2 Pet. i. 10.}\emph{ Ita demum suppeditabit hæc doctrina laudandi, reverendi, admirandique Deum argumentum,}\footnote{Eph. i. 6; Rom. xi. 33.} \emph{quin etiam humilitatis, }

to all that sincerely obey the gospel.\footnote{Rom. xi. 5, 6, 20; 2 Pet. i. 10; Rom. viii. 33; Luke x. 20.}

\emph{diligentiæ et consolationis copiosæ omnibus sincere obedientibus evangelio.}\footnote{Rom. xi. 5, 6, 20; 2 Pet. i. 10; Rom. viii. 33; Luke x. 20.}

\xsection{Chapter IV. Of Creation}

I. It pleased God the Father, Son, and Holy Ghost,\footnote{Heb. i. 2; John i. 2, 3; Gen. i. 2; Job xxvi. 13; xxxiii. 4.} for the manifestation of the glory of his eternal power, wisdom, and goodness,\footnote{Rom. i. 20; Jer. x. 12; Psa. civ. 24; xxxiii. 5, 6.} in the beginning, to create or make of nothing the world, and all things therein, whether visible or invisible, in the space of six days, and all very good.\footnote{Gen. ch. i.; Heb. xi. 3; Col. i. 16; Acts xvii. 24.}

I. \emph{Deo, Patri, Filio et Spiritui sancto, complacitum est,}\footnote{Heb. i. 2; John i. 2, 3; Gen. i. 2; Job xxvi. 13; xxxiii. 4.}\emph{ quo æternæ suæ cum potentiæ tum sapientiæ bonitatisque gloriam manifestaret,}\footnote{Rom. i. 20; Jer. x. 12; Psa. civ. 24; xxxiii. 5, 6.}\emph{ mundum hunc, et quæ in eo continentur universa tam visibilla quam invisibilia, in principio intra sex dierum spatium creare, seu ex nihilo condere, atque omnia quidem bona valde.}\footnote{Gen. ch. i.; Heb. xi. 3; Col. i. 16; Acts xvii. 24.}

II. After God had made all other creatures, he created man, male and female,\footnote{Gen. i. 27.} with reasonable and immortal souls,\footnote{Gen. ii. 7; Eccles. xii. 7; Luke xxiii. 43; Matt. x. 28.} endued with knowledge, righteousness, and true holiness, after his own image,\footnote{Gen. i. 26; Col. iii. 10; Eph. iv. 24.} having the law of God written in their hearts,\footnote{Rom. ii. 14, 15.} and power to fulfill it;\footnote{Eccles. vii. 29.} and yet under a possibility of transgressing, being left to the liberty of their own will, which was subject unto change.\footnote{Gen. iii. 6; Eccles. vii. 29.} Beside this law written in their hearts, they received a command not to eat of the tree of the knowledge of good

II. \emph{Postquam omnes alias creaturas condidisset Deus, creavit hominem marem et fœminam,}\footnote{Gen. i. 27.}\emph{ animabus inditis rationalibus ac immortalibus,}\footnote{Gen. ii. 7; Eccles. xii. 7; Luke xxiii. 43; Matt. x. 28.}\emph{ imbutos cognitione, justitia, veraque sanctitate, ad suam ipsius imaginem,}\footnote{Gen. i. 26; Col. iii. 10; Eph. iv. 24.}\emph{ habentes in cordibus suis inscriptam Divinam legem,}\footnote{Rom. ii. 14, 15.}\emph{ simul et eandem implendi vires;}\footnote{Eccles. vii. 29.}\emph{ non tamen sine quadam violandi possibilitate; libertati siquidem permissi erant voluntatis suæ haud immutabilis.}\footnote{Gen. iii. 6; Eccles. vii. 29.}\emph{ Præter autem hanc in cordibus eorum inscriptam legem de non comedendo ex arbore scientiæ boni malique }

and evil; which while they kept they were happy in their communion with God,\footnote{Gen. ii. 27; iii. 8–11, 23.} and had dominion over the creatures.\footnote{Gen. i. 26, 28; [Am. ed. Psa. viii. 6–8].}

\emph{mandatum insuper acceperunt; quod certe quam diu observabant, communione Dei beati erant,}\footnote{Gen. ii. 27; iii. 8–11, 23.}\emph{ dominiumque habebant in creaturas.}\footnote{Gen. i. 26, 28; [Am. ed. Psa. viii. 6–8].}

\xsection{Chapter V. Of Providence}

God, the great Creator of all things, doth uphold,\footnote{Heb. i. 3.} direct, dispose, and govern all creatures, actions, and things,\footnote{Dan. iv. 34, 35; Psa. cxxxv. 6; Acts xvii. 25, 26, 28; Job, chaps. xxxviii. xxxix. xl. xli.} from the greatest even to the least,\footnote{Matt. x. 29–31; [Am. ed. Matt. vi. 26, 30].} by his most wise and holy providence,\footnote{Prov. xv. 3; [Am. ed. 2 Chron. xvi. 9]; Psa. civ. 24; cxlv. 17.} according to his infallible foreknowledge\footnote{Acts xv. 18; Psa. xciv. 8–11.} and the free and immutable counsel of his own will,\footnote{Eph. i. 11; Psa. xxxiii. 10, 11.} to the praise of the glory of his wisdom, power, justice, goodness, and mercy.\footnote{Isa. lxiii. 14; Eph. iii. 10; Rom. ix. 17; Gen. xlv. 7; Psa. cxlv. 7.}

I. \emph{Magnus ille rerum omnium creator Deus sapientissima sua et sanctissima simul providentia}\footnote{Heb. i. 3.}\emph{ creaturas, actiones, resque}\footnote{Dan. iv. 34, 35; Psa. cxxxv. 6; Acts xvii. 25, 26, 28; Job, chaps. xxxviii. xxxix. xl. xli.}\emph{ a maximis usque ad minimas}\footnote{Matt. x. 29–31; [Am. ed. Matt. vi. 26, 30].} \emph{universas sustentat,}\footnote{Prov. xv. 3; [Am. ed. 2 Chron. xvi. 9]; Psa. civ. 24; cxlv. 17.}\emph{ dirigit, ordinat, gubernatque secundum infallibilem suam præscientiam,}\footnote{Acts xv. 18; Psa. xciv. 8–11.}\emph{ et voluntatis suæ consilium liberum ac immutabile,}\footnote{Eph. i. 11; Psa. xxxiii. 10, 11.}\emph{ ad laudem gloriæ sapientiæ suæ, potentiæ, justitiæ, bonitatis, ac misericordiæ.}\footnote{Isa. lxiii. 14; Eph. iii. 10; Rom. ix. 17; Gen. xlv. 7; Psa. cxlv. 7.}

II. Although in relation to the foreknowledge and decree of God, the first cause, all things come to pass immutably and infallibly,\footnote{Acts ii. 23.} yet by the same providence he ordereth them to fall out, according to the nature of second causes, either necessarily, freely, or contingently.\footnote{Gen. viii. 22; Jer. xxxi. 35; Exod. xxi. 13; Deut. xix. 5; 1 Kings xxii. 28, 34; Isa. x. 6, 7.}

II. \emph{Quamvis respectu præscientiæ ac decreti Dei} (\emph{causæ primæ}) \emph{omnia immutabiliter atque infallibiliter eveniant,}\footnote{Acts ii. 23.}\emph{ per eandem tamen ille providentiam eadem ordinat evenire necessario, libere, aut contingenter, pro natura causarum secundarum.}\footnote{Gen. viii. 22; Jer. xxxi. 35; Exod. xxi. 13; Deut. xix. 5; 1 Kings xxii. 28, 34; Isa. x. 6, 7.}

III. God, in his ordinary providence, maketh use of means,\footnote{Acts xxvii. 31, 44; Isa. lv. 10, 11; Hos. ii. 21, 22.} yet is free to work without,\footnote{Hos. i. 7; Matt. iv. 4; Job xxxiv. 10.}

III. \emph{Deus in providentia sua ordinaria mediis utitur,}\footnote{Acts xxvii. 31, 44; Isa. lv. 10, 11; Hos. ii. 21, 22.}\emph{ iis tamen non astringitur, quo minus absque eis,}\footnote{Hos. i. 7; Matt. iv. 4; Job xxxiv. 10.}

above,\footnote{Rom. iv. 19–21.} and against them, at his pleasure.\footnote{2 Kings vi. 6; Dan. iii. 27.}

\emph{supra}\footnote{Rom. iv. 19–21.}\emph{ aut etiam contra ea pro arbitrio suo operetur.}\footnote{2 Kings vi. 6; Dan. iii. 27.}

IV. The almighty power, unsearchable wisdom, and infinite goodness of God so far manifest themselves in his providence that it extendeth itself even to the first fall, and all other sins of angels and men,\footnote{Rom. xi. 32–34; 2 Sam. xxiv. 1; 1 Chron. xxi. 1; 1 Kings xxii. 22, 23; 1 Chron. x. 4, 13, 14; 2 Sam. xvi. 10; Acts ii. 23; iv. 27, 28.} and that not by a bare permission,\footnote{Acts xiv. 16.} but such as hath joined with it a most wise and powerful bounding,\footnote{Psa. lxxvi. 10; 2 Kings xix. 28.} and otherwise ordering and governing of them, in a manifold dispensation, to his own holy ends;\footnote{Gen. l. 20; Isa. x. 6, 7, 12.} yet so as the sinfulness thereof proceedeth only from the creature, and not from God; who, being most holy and righteous, neither is nor can be the author or approver of sin.\footnote{1 James i. 13, 14, 17; 1 John. ii. 16; Psa. l. 21.}

IV.\emph{ Omnipotentem Dei potentiam, sapientiam inscrutabilem, bonitatemque infinitam providentia ejus eo usque manifestat, ut vel ad primum lapsum, omniaque reliqua peccata, seu hominum sint sive angelorum, se extendat;}\footnote{Rom. xi. 32–34; 2 Sam. xxiv. 1; 1 Chron. xxi. 1; 1 Kings xxii. 22, 23; 1 Chron. x. 4, 13, 14; 2 Sam. xvi. 10; Acts ii. 23; iv. 27, 28.}\emph{ neque id quidem permissione nuda,}\footnote{Acts xiv. 16.}\emph{ verum cui conjuncta est sapientissima potentissimaque eorum limitatio,}\footnote{Psa. lxxvi. 10; 2 Kings xix. 28.}\emph{ ac aliusmodi ad sanctos sibi propositos fines dispensatione multiplici ordinatio et gubernatio;}\footnote{Gen. l. 20; Isa. x. 6, 7, 12.}\emph{ ita tamen ut omnis eorum vitiositas a sola proveniat creatura, a Deo neutiquam, qui sanctissimus quum sit justissimusque neque est, nec esse quidem potest peccati autor aut approbator.}\footnote{1 James i. 13, 14, 17; 1 John. ii. 16; Psa. l. 21.}

V. The most wise, righteous, and gracious God doth oftentimes leave for a season his own children to manifold temptations and the corruption of their own hearts, to chastise them for their former sins, or to discover unto them the hidden strength of corruption and deceitfulness of their hearts, that they may be humbled;\footnote{2 Chron. xxxii. 25, 26, 31; 2 Sam. xxiv. 1.} and to raise them to a more close and constant

V. \emph{Sapientissimus, justissimus, et gratiosissimus idem Deus, sæpenumero filios suos tentationibus multifariis, suorumque cordium corruptioni ad tempus permittit; quo ob admissa prius peccata castiget eos, vel corruptionis iis detegat vim occultam, cordiumque suorum fraudulentiam ut humilientur;}\footnote{2 Chron. xxxii. 25, 26, 31; 2 Sam. xxiv. 1.}\emph{ quoque eos excitet ad strictam magis et constantem a seipso proferendis suppetiis }

dependence for their support unto\footnote{[Am. ed. \emph{upon}.]} himself, and to make them more watchful against all future occasions of sin, and for sundry other just and holy ends.\footnote{2 Cor. xii. 7–9; Psa. lxxiii. throughout; lxxvii. 1–10, 12; Mark xiv. 66 to the end; John xxi. 15–17.}

\emph{dependentiam; Quo denique adversus onmes occasiones peccati de futuro reddat cautiores. Sed et ob alios etiam varios fines, justos sanctosque sibi propositos.}\footnote{2 Cor. xii. 7–9; Psa. lxxiii. throughout; lxxvii. 1–10, 12; Mark xiv. 66 to the end; John xxi. 15–17.}

VI. As for those wicked and ungodly men whom God, as a righteous judge, for former sins, doth blind and harden,\footnote{Rom. i. 24, 26, 28; xi. 7, 8.} from them he not only withholdeth his grace, whereby they might have been enlightened in their understandings and wrought upon in their hearts,\footnote{Deut. xxix. 4.} but sometimes also withdraweth the gifts which they had,\footnote{Matt. xiii. 12; xxv. 29.} and exposeth them to such objects as their corruption makes occasion of sin;\footnote{Deut. ii. 30; 2 Kings viii. 12, 13.} and withal, gives them over to their own lusts, the temptations of the world, and the power of Satan;\footnote{Psa. lxxxi. 11, 12; 2 Thess. ii. 10–12.} whereby it comes to pass that they harden themselves, even under those means which God useth for the softening of others.\footnote{Exod. vii. 3; viii. 15, 32; 2 Cor. ii. 15, 16; Isa. viii. 14; 1 Pet. ii. 7, 8; Isa. vi. 9, 10; Acts xxviii. 26, 27.}

VI. \emph{Quod scelestos illos spectat impiosque homines, quos Deus, ut justus judex, ob peccata præcedentia excæcat induratque;}\footnote{Rom. i. 24, 26, 28; xi. 7, 8.}\emph{ eis ille non solum gratiam suam non impertit, qua ipsis cum illuminari intellectus, tum affici corda potuissent;}\footnote{Deut. xxix. 4.}\emph{ sed interdum subtrahit eis quibus imbuti erant dona,}\footnote{Matt. xiii. 12; xxv. 29.}\emph{ et ipsos exponit illiusmodi objectis, unde corruptio eorum arripit sibi peccandi occasiones;}\footnote{Deut. ii. 30; 2 Kings viii. 12, 13.}\emph{ simulque tradit eos suis ipsorum concupiscentiis et tentationibus mundi, et potestati Satanæ;}\footnote{Psa. lxxxi. 11, 12; 2 Thess. ii. 10–12.}\emph{ ex quo fit ut seipsos ipsi indurent, et quidem sub iisdem mediis quibus utitur Deus ad alios emolliendos.}\footnote{Exod. vii. 3; viii. 15, 32; 2 Cor. ii. 15, 16; Isa. viii. 14; 1 Pet. ii. 7, 8; Isa. vi. 9, 10; Acts xxviii. 26, 27.}

VII. As the providence of God doth, in general, reach to all creatures, so, after a most special manner, it taketh care of his Church, and disposeth all things to the good thereof.\footnote{1 Tim. iv. 10; Amos ix. 8, 9; Rom. viii 28; Isa. xliii. 3–5,14.}

VII. \emph{Providentia Dei sicut ad omnes creaturas universali modo se extendit; ita modo plane peculiari Ecclesiæ suæ curam gerit, ac in ejus bonum disponit universa.}\footnote{1 Tim. iv. 10; Amos ix. 8, 9; Rom. viii 28; Isa. xliii. 3–5,14.}

\xsection{Chapter VI. Of the Fall of Man, of Sin, and of the Punishment thereof}

I. Our first parents, being seduced by the subtilty and temptation of Satan, sinned in eating the forbidden fruit.\footnote{Gen. iii. 13; 2 Cor. xi. 3.} This their sin God was pleased, according to his wise and holy counsel, to permit, having purposed to order it to his own glory.\footnote{Rom. xi. 32.}

I. \emph{Primi parentes, Satanæ subtilitate ac tentatione seducti, fructus vetiti esu peccaverunt.}\footnote{Gen. iii. 13; 2 Cor. xi. 3.}\emph{ Hoc eorum peccatum secundum sapiens suum sanctumque consilium Deo placuit permittere, non sine proposito illud ad suam ipsius gloriam ordinandi.}\footnote{Rom. xi. 32.}

II. By this sin they fell from their original righteousness and communion with God,\footnote{Gen. iii. 6–8; Eccles. vii. 29; Rom. iii. 23.} and so became dead in sin,\footnote{Gen. ii. 17; Eph. ii. 1; [Am. ed. Rom. v. 12].} and wholly defiled in all the faculties and parts of soul and body.\footnote{Tit. i. 15; Gen. vi. 5; Jer. xvii. 9; Rom. iii. 10–19.}

II. \emph{Hoc illi peccato, justitia sua originali et communione cum Deo exciderunt;}\footnote{Gen. iii. 6–8; Eccles. vii. 29; Rom. iii. 23.}\emph{ itaque facti sunt in peccato mortui,}\footnote{Gen. ii. 17; Eph. ii. 1; [Am. ed. Rom. v. 12].}\emph{ atque in omnibus facultatibus ac partibus animæ corporisque penitus contaminati.}\footnote{Tit. i. 15; Gen. vi. 5; Jer. xvii. 9; Rom. iii. 10–19.}

III. They being the root of all mankind,\footnote{Gen. i. 27, 28; ii. 16, 17; Acts xvii. 26; Rom. v. 12, 15–19; 1 Cor. xv. 21, 22, 45, 49.} the guilt of this sin was imputed, and the same death in sin and corrupted nature conveyed to all their posterity descending from them by ordinary generation.\footnote{Psa. li. 5; Gen. v. 3; Job xiv. 4; xv. 14.}

III. \emph{Quumque illi fuerint radix totius humani}\footnote{Gen. i. 27, 28; ii. 16, 17; Acts xvii. 26; Rom. v. 12, 15–19; 1 Cor. xv. 21, 22, 45, 49.}\emph{ generis, hujusce peccati reatus fuit imputatus, eademque in peccato mors ac natura corrupta propagata, omnibus illorum posteris, quotquot ab iis ordinaria quidem generatione procreantur.}\footnote{Psa. li. 5; Gen. v. 3; Job xiv. 4; xv. 14.}

IV. From this original corruption, whereby we are utterly indisposed, disabled, and made opposite to all good,\footnote{Rom. v. 6; vii. 18; viii. 7; Col. i. 21; [Am. ed. John iii. 6].} and wholly inclined to all evil,\footnote{Gen. vi. 5; viii. 21; Rom. iii. 10–12.} do proceed all actual transgressions.\footnote{James i. 14, 15; Eph. ii. 2, 3; Matt. xv. 19.}

IV. \emph{Ab hac originali labe} (\emph{qua ad omne bonum facti sumus inhabiles prorsus ac impotentes, eique plane oppositi,}\footnote{Rom. v. 6; vii. 18; viii. 7; Col. i. 21; [Am. ed. John iii. 6].}\emph{ ad malum autem omne proclives penitus})\footnote{Gen. vi. 5; viii. 21; Rom. iii. 10–12.}\emph{ proveniunt omnia peccata actualia.}\footnote{James i. 14, 15; Eph. ii. 2, 3; Matt. xv. 19.}

V. This corruption of nature,

V. \emph{Hæc naturæ corruptio durante }

during this life, doth remain in those that are regenerated;\footnote{1 John i. 8, 10; Rom. vii. 14, 17, 18, 23; James iii. 2; Prov. xx. 9; Eccles. vii. 20.} and although it be through Christ pardoned and mortified, yet both itself and all the motions thereof are truly and properly sin.\footnote{Rom. vii. 5, 7, 8, 25; Gal. v. 17.}

\emph{hac vita manet etiam in regenitis;}\footnote{1 John i. 8, 10; Rom. vii. 14, 17, 18, 23; James iii. 2; Prov. xx. 9; Eccles. vii. 20.}\emph{ et quamvis per Christum et condonata sit et mortificata; nihilo minus tam ipsa, quam ejus motus universi vere sunt ac proprie peccata.}\footnote{Rom. vii. 5, 7, 8, 25; Gal. v. 17.}

VI. Every sin, both original and actual, being a transgression of the righteous law of God, and contrary thereunto,\footnote{1 John iii. 4.} doth, in its own nature, bring guilt upon the sinner,\footnote{Rom. ii. 15; iii. 9, 19.} whereby he is bound over to the wrath of God\footnote{Eph. ii. 3.} and curse of the law,\footnote{Gal. iii. 10.} and so made subject to death,\footnote{Rom. vi. 23.} with all miseries spiritual,\footnote{Eph. iv. 18.} temporal,\footnote{Rom. viii. 20; Lam. iii. 39.} and eternal.\footnote{Matt. xxv. 41; 2 Thess. i. 9.}

VI. \emph{Peccatum omne cum originale tum actuale, quum justæ Dei legis transgressio sit eique contraria,}\footnote{1 John iii. 4.}\emph{ peccatori suapte natura reatum infert,}\footnote{Rom. ii. 15; iii. 9, 19.}\emph{ quo ad iram Dei,}\footnote{Eph. ii. 3.}\emph{ ac maledictionem legis}\footnote{Gal. iii. 10.}\emph{ subeundam obligatur, adeoque redditur obnoxius morti}\footnote{Rom. vi. 23.}\emph{ simul et miseriis omnibus spiritualibus,}\footnote{Eph. iv. 18.}\emph{ temporalibus,}\footnote{Rom. viii. 20; Lam. iii. 39.}\emph{ ac æternis.}\footnote{10 Matt. xxv. 41; 2 Thess. i. 9.}

\xsection{Chapter VII. Of God's Covenant with Man}

I. The distance between God and the creature is so great that although reasonable creatures do owe obedience unto him as their Creator, yet they could never have any fruition of him as their blessedness and reward but by some voluntary condescension on God's part, which he hath been pleased to express by way of covenant.\footnote{Isa. xl. 13–17; Job ix. 32, 33; 1 Sam. ii. 25; Psa. c. 2, 3; cxiii. 5, 6; Job xxii. 2, 3; xxxv. 7, 8; Luke xvii. 10; Acts xvii. 24, 25.}

I. \emph{Tanta est inter deum et creaturam distantia, ut licet creaturæ rationales obedientiam illi ut creatori suo debeant, nullam tamen fruitionem ejus tanquam suæ beatitudinis ac præmii habere unquam potuissent, ni voluntaria fuisset aliqua ex parte Dei condescentio; quam ipsi exprimere placuit icto fœdere.}\footnote{Isa. xl. 13–17; Job ix. 32, 33; 1 Sam. ii. 25; Psa. c. 2, 3; cxiii. 5, 6; Job xxii. 2, 3; xxxv. 7, 8; Luke xvii. 10; Acts xvii. 24, 25.}

II. The first covenant made with

II. \emph{Primum fœdus cum hominibus }

man was a covenant of works,\footnote{Gal. iii. 12; [Am. ed. Hos. vi. 7; Gen. ii. 16, 17].} wherein life was promised to Adam, and in him to his posterity,\footnote{Rom. v. 12–20; x. 5.} upon condition of perfect and personal obedience.\footnote{Gen. ii. 17; Gal. iii. 10.}

\emph{initum erat fœdus operum,}\footnote{Gal. iii. 12; [Am. ed. Hos. vi. 7; Gen. ii. 16, 17].}\emph{ quo vita Adamo promissa erat, ejusque in eo posteris,}\footnote{Rom. v. 12–20; x. 5.}\emph{ sub conditione obedientiæ perfectæ ac personalis.}\footnote{Gen. ii. 17; Gal. iii. 10.}

III. Man by his fall having made himself incapable of life by that covenant, the Lord was pleased to make a second,\footnote{Gal. iii. 21; Rom. iii. 20, 21; viii. 3; Gen. iii. 15; Isa. xlii. 6.} commonly called the covenant of grace: wherein he freely offered unto sinners life and salvation by Jesus Christ, requiring of them faith in him that they may be saved,\footnote{5 Mark xvi. 15, 16; John iii. 16; Rom. x. 6, 9; Gal. iii. 11.} and promising to give unto all those that are ordained unto life his Holy Spirit, to make them willing and able to believe.\footnote{Ezek. xxxvi. 26, 27; John vi. 44, 45; [Am. ed. v. 37].}

III. \emph{Quum autem homo lapsu suo omnem sibi præstruxisset ad vitam aditum per illud fœdus, complacuit Domino secundum inire,}\footnote{Gal. iii. 21; Rom. iii. 20, 21; viii. 3; Gen. iii. 15; Isa. xlii. 6.}\emph{ quod vulgo dicimus }Fœdus Gratiæ; \emph{in quo peccatoribus offert gratuito vitam ac salutem per Jesum Christum, fidem in illum ab iis requirens ut salventur;}\footnote{5 Mark xvi. 15, 16; John iii. 16; Rom. x. 6, 9; Gal. iii. 11.}\emph{ promittensque omnibus qui ad vitam ordinantur se spiritum suum sanctum daturum, qui in illis operetur credendi cum voluntatem tum potentiam.}\footnote{Ezek. xxxvi. 26, 27; John vi. 44, 45; [Am. ed. v. 37].}

IV. This covenant of grace is frequently set forth in the Scripture by the name of a testament, in reference to the death of Jesus Christ the testator, and to the everlasting inheritance, with all things belonging to it, therein bequeathed.\footnote{Heb. ix. 15–17; vii. 22; Luke xxii. 20; 1 Cor. xi. 25.}

IV. \emph{Hoc fœdus Gratiæ in Scriptura sæpe nomine} Testamenti \emph{indigitatur, respectu nimirum mortis Testatoris Jesu Christi, æternæque illius hæreditatis, quam is una cum omnibus eam spectantibus inibi legabat.}\footnote{Heb. ix. 15–17; vii. 22; Luke xxii. 20; 1 Cor. xi. 25.}

V. This covenant was differently administered in the time of the law and in the time of the gospel:\footnote{2 Cor. iii. 6–9.} under the law it was administered by promises, prophecies,

V. \emph{Hoc fœdus sub Lege atque sub Evangelio administratum est modo alio atque alio.}\footnote{2 Cor. iii. 6–9.}\emph{ Sub Lege quidem per promissiones, prophetias et sacrificia, per circumcisionem, agnum }

sacrifices, circumcision, the paschal lamb, and other types and ordinances delivered to the people of the Jews, all fore-signifying Christ to come,\footnote{Heb., chaps. viii. ix. x.; Rom. iv. 11; Col. ii. 11,12; 1 Cor. v. 7; [Am. ed. Col. ii. 17].} which were for that time sufficient and efficacious, through the operation of the Spirit, to instruct and build up the elect in faith in the promised Messiah,\footnote{1 Cor. x. 1–4; Heb. xi. 13; John viii. 56.} by whom they had full remission of sins and eternal salvation; and is called the Old Testament.\footnote{;Gal. iii. 7–9, 14.}

\emph{pascalem, aliosque typos ac instituta populo Judaico tradita, quæ omnia Venturum Christum præsignificabant;}\footnote{Heb., chaps. viii. ix. x.; Rom. iv. 11; Col. ii. 11,12; 1 Cor. v. 7; [Am. ed. Col. ii. 17].}\emph{ erantque pro ratione illorum temporum sufficientia, et per operationem spiritus efficacia ad electos instruendum ac ædificandum in fide in promissum Messiam,}\footnote{1 Cor. x. 1–4; Heb. xi. 13; John viii. 56.}\emph{ per quem plenum peccatorum remissionem et salutem æternam sunt consecuti; diciturque} Vetus Testamentum.\footnote{;Gal. iii. 7–9, 14.}

VI. Under the gospel, when Christ the substance\footnote{Gal. ii. 17; [Am. ed. Col. ii. 17].} was exhibited, the ordinances in which this covenant is dispensed are the preaching of the word and the administration of the sacraments of Baptism and the Lord's Supper;\footnote{Matt. xxviii. 19, 20; 1 Cor. xi. 23–25; [Am. ed. 2 Cor. iii. 7–11].} which, though fewer in number, and administered with more simplicity aud less outward glory, yet in them it is held forth in more fullness, evidence, and spiritual efficacy,\footnote{Heb. xii. 22–28; Jer. xxxi. 33, 34.} to all nations, both Jews and Gentiles;\footnote{Matt. xxviii. 19; Eph. ii. 15–19.} and is called the New Testament.\footnote{Luke xxii. 20; [Am. ed. Heb. viii. 7–9].} There are not, therefore, two covenants of grace differing in substance, but one and the same under various dispensations.\footnote{Gal. iii. 14, 16; Acts xv. 11; Rom. iii. 21–23, 30; Psa. xxxii. 1; Rom. iv. 3, 6, 16, 17, 23, 24; Heb. xiii. 8.}

VI. \emph{Sub evangelio autem, exhibito jam Christo, substantia}\footnote{Gal. ii. 17; [Am. ed. Col. ii. 17].}\emph{ scilicet ac antitypo, præscriptæ rationes in quibus hoc fœdus dispensatur, sunt prædicatio verbi, et administratio sacramentorum, baptismi nempe ac cœnæ Dominicæ;}\footnote{Matt. xxviii. 19, 20; 1 Cor. xi. 23–25; [Am. ed. 2 Cor. iii. 7–11].}\emph{ in quibus quidem utut numero paucioribus, iisque simplicius ac minore cum externa gloria administratis, cum majore tamen plenitudine, evidentia, et efficacia spirituali}\footnote{Heb. xii. 22–28; Jer. xxxi. 33, 34.}\emph{ populis cunctis tam Judæis quam Gentibus}\footnote{Matt. xxviii. 19; Eph. ii. 15–19.}\emph{ exhibetur; Diciturque} Novum Testamentum.\footnote{Luke xxii. 20; [Am. ed. Heb. viii. 7–9].}\emph{ Non sunt ergo duo fœdera gratiæ, re atque natura discrepantia; sed unum idemque, licet non uno modo dispensatum.}\footnote{Gal. iii. 14, 16; Acts xv. 11; Rom. iii. 21–23, 30; Psa. xxxii. 1; Rom. iv. 3, 6, 16, 17, 23, 24; Heb. xiii. 8.}.

\xsection{Chapter VIII. Of Christ the Mediator}

I. It pleased God, in his eternal purpose, to choose and ordain the Lord Jesus, his only-begotten Son, to be the Mediator between God and man,\footnote{Isa. xlii. 1; 1 Pet. i. 19, 20; John iii. 16; 2 Tim. ii. 5.} the Prophet,\footnote{Acts iii. 22; [Am. ed. Deut. xviii. 15].} Priest,\footnote{Heb. v. 5, 6.} and King;\footnote{Psa. ii. 6; Luke i. 33.} the Head and Saviour of his Church,\footnote{Eph. v. 23.} the Heir of all things,\footnote{Heb. i. 2.} and Judge of the world;\footnote{Acts xvii. 31.} whom he did, from all eternity, give a people to be his seed,\footnote{John xvii. 6; Psa. xxii. 30; Isa. liii. 10.} and to be by him in time redeemed, called, justified, sanctified, and glorified.\footnote{1 Tim. ii. 6; Isa. lv. 4, 5; 1 Cor. i. 30.}

I. \emph{Complacitum est Deo Filium ejus unigenitum Dominum Jesum in æterno suo proposito eligere atque ordinare ut Mediator esset inter Deum et hominem,}\footnote{Isa. xlii. 1; 1 Pet. i. 19, 20; John iii. 16; 2 Tim. ii. 5.} \emph{Propheta},\footnote{Acts iii. 22; [Am. ed. Deut. xviii. 15].} \emph{Sacerdos},\footnote{Heb. v. 5, 6.} \emph{et Rex},\footnote{Psa. ii. 6; Luke i. 33.} \emph{caput idem et salvator Ecclesiæ suæ};\footnote{Eph. v. 23.} \emph{rerum omnium hæres},\footnote{Heb. i. 2.} \emph{Mundique Judex};\footnote{Acts xvii. 31.} \emph{cui ab æterno populum dedit futurum illi in semen},\footnote{John xvii. 6; Psa. xxii. 30; Isa. liii. 10.} \emph{ac per illum stato tempore redimendum, vocandum, justificandum, sanctificandum ac glorificandum}.\footnote{1 Tim. ii. 6; Isa. lv. 4, 5; 1 Cor. i. 30.}

II. The Son of God, the second person in the Trinity, being very and eternal God, of one substance, and equal with the Father did, when the fullness of time was come, fake upon him man's nature,\footnote{John i. 1, 14; 1 John v. 20; Phil. ii. 6; Gal. iv. 4.} with all the essential properties and common infirmities thereof, yet without sin:\footnote{Heb. ii. 14, 16, 17; iv. 15.} being conceived by the power of the Holy Ghost in the womb of the Virgin Mary, of her substance.\footnote{Luke i. 27, 31, 35; Gal. iv. 4.} So that two whole, perfect, and distinct natures, the Godhead and the manhood, were inseparably joined together

II. \emph{Filius Dei persona secunda in Trinitate, verus nempe idem æternusque Deus, substantiæ cum Patre unius ejusdemque, eique coæqualis, cum advenerat temporis plenitudo, assumpsit naturam humanam,}\footnote{John i. 1, 14; 1 John v. 20; Phil. ii. 6; Gal. iv. 4.} \emph{una cum omnibus ejus proprietatibus essentialibus, communibusque infirmitatibus, immunem tamen a peccato,}\footnote{Heb. ii. 14, 16, 17; iv. 15.}\emph{ conceptus scilicet in utero eque substantia Mariæ Virginis,}\footnote{Luke i. 27, 31, 35; Gal. iv. 4.}\emph{ virtute Spiritus Sancti. Adeo sane ut naturæ duæ, integræ, perfectæ, distinctæque Deitas ac humanitas in una eademque }

in one person, without conversion, composition, or confusion.\footnote{Luke i. 35; Col. ii. 9; Rom. ix. 5; 1 Pet. iii. 18; 1 Tim. iii. 16.} Which person is very God and very man, yet one Christ, the only mediator between God and man.\footnote{Rom. i. 3, 4; 1 Tim. ii. 5.}

\emph{persona indissolubili nexu conjunctæ fuerint, sine conversione, compositione, aut confusione.}\footnote{Luke i. 35; Col. ii. 9; Rom. ix. 5; 1 Pet. iii. 18; 1 Tim. iii. 16.}\emph{ Quæ quidem persona vere Deus est ac vere homo, unus tamen Christus, unicus inter Deum et hominem Mediator.}\footnote{Rom. i. 3, 4; 1 Tim. ii. 5.}

III. The Lord Jesus, in his human nature thus united to the divine, was sanctified and anointed with the Holy Spirit above measure;\footnote{Psa. xlv. 7; John iii. 34.} having in him all the treasures of wisdom and knowledge,\footnote{Col. ii. 3.} in whom it pleased the Father that all fullness should dwell;\footnote{Col. i. 19.} to the end that, being holy, harmless, undefiled, and full of grace and truth,\footnote{Heb. vii. 26; John i. 14.} he might be thoroughly furnished to execute the office of a mediator and surety.\footnote{Acts x. 38; Heb. xii. 24; vii. 22.} Which office he took not unto himself, but was thereunto called by his Father,\footnote{Heb. v. 4, 5.} who put all power and judgment into his hand, and gave him commandment to execute the same.\footnote{John v. 22, 27; Matt. xxviii. 18; Acts ii. 36.}

III. \emph{Dominus Jesus in humana sua natura divinæ hunc modum conjuncta sanctificatus est, ac Spiritu sancto supra mensuram unctus,}\footnote{Psa. xlv. 7; John iii. 34.}\emph{ in se habens omnes sapientiæ notitiæqum thesauros;}\footnote{Col. ii. 3.}\emph{ in quo Patri visum est ut omnis plenitudo inhabitaret,}\footnote{Col. i. 19.}\emph{ atque eo quidem fine ut sanctus, innocuus, intaminatus, plenusque gratiæ ac veritatis existens,}\footnote{Heb. vii. 26; John i. 14.}\emph{ ad Mediatoris Vadisque munus exequendum perfecte esset instructus.}\footnote{Acts x. 38; Heb. xii. 24; vii. 22.}\emph{ Quod ille officium non arripuit sibi, verum a Patre erat ad id vocatus,}\footnote{Heb. v. 4, 5.}\emph{ qui omnem ei potestatem ac judicium in manus dedit, und cum mandato exercendi.}\footnote{John v. 22, 27; Matt. xxviii. 18; Acts ii. 36.}

IV. This office the Lord Jesus did most willingly undertake,\footnote{Psa. xl. 7, 8; Heb. x. 5–10; John x. 18; Phil. ii. 8.} which, that he might discharge, he was made under the law,\footnote{Gal. iv. 4.} and did perfectly fulfill it;\footnote{Matt. iii. 15; v. 17.} endured most grievous torments immediately in his soul,\footnote{Matt. xxvi. 37, 38; Luke xxii. 44; Matt. xxvii. 46.}

IV. \emph{Hoc munus promtissima voluntate in se suscepit Dominus Jesus,}\footnote{Psa. xl. 7, 8; Heb. x. 5–10; John x. 18; Phil. ii. 8.}\emph{ quod ut expleret factus est sub Lege,}\footnote{Gal. iv. 4.}\emph{ eam perfecte implevit,}\footnote{Matt. iii. 15; v. 17.}\emph{ immediate in anima,}\footnote{Matt. xxvi. 37, 38; Luke xxii. 44; Matt. xxvii. 46.}\emph{ sua gravissimos subiit cruciatus, in corpore}\footnote{Matt., chaps. xxvi. xxvii.}

and most painful sufferings in his body;\footnote{Matt., chaps. xxvi. xxvii.} was crucified, and died;\footnote{Phil. ii. 8.} was buried, and remained under the power of death, yet saw no corruption.\footnote{Acts ii. 23, 24, 27; xiii. 37; Rom. vi. 9.} On the third day he arose from the dead,\footnote{1 Cor. xv. 3, 4.} with the same body in which he suffered;\footnote{John xx. 25, 27.} with which also he ascended into heaven, and there sitteth at the right hand of his Father,\footnote{Mark xvi. 19.} making intercession;\footnote{Rom. viii. 34; Heb. ix. 24; vii. 25.} and shall return to judge men and angels at the end of the world.\footnote{Rom. xiv. 9, 10; Acts i. 11; x. 42; Matt. xiii. 40–42; Jude 6; 2 Pet. ii. 4.}

\emph{vero perpessiones quam maxime dolorificas; crucifixus est, ac mortuus;}\footnote{Phil. ii. 8.}\emph{ sepultus est, mansitque sub mortis potestate; nec tamen ullam vidit corruptionem.}\footnote{Acts ii. 23, 24, 27; xiii. 37; Rom. vi. 9.}\emph{ Tertio die surrexit a mortuis,}\footnote{1 Cor. xv. 3, 4.}\emph{ cum eodem in quo passus fuerat corpore,}\footnote{John xx. 25, 27.}\emph{ cum quo etiam ascendit in cœlum, ibique sedens ad dextram Patris}\footnote{Mark xvi. 19.}\emph{ intercedit,}\footnote{Rom. viii. 34; Heb. ix. 24; vii. 25.}\emph{ rediturus inde in consummatione mundi, ad homines angelosque judicandum.}\footnote{Rom. xiv. 9, 10; Acts i. 11; x. 42; Matt. xiii. 40–42; Jude 6; 2 Pet. ii. 4.}

V. The Lord Jesus, by his perfect obedience and sacrifice of himself, which he through the eternal Spirit once offered up unto God, hath fully satisfied the justice of his Father,\footnote{Rom. v. 19; Heb. ix. 14, 16; x. 14; Eph. v. 2; Rom. iii. 25, 26.} and purchased not only reconciliation, but an everlasting inheritance in the kingdom of heaven, for all those whom the Father hath given unto him.\footnote{Dan. ix. 24, 26; Col. i. 19, 20; Eph. i. 11, 14; John xvii. 2; Heb. ix. 12, 15.}

V. \emph{Dominus Jesus obedientia sua perfecta, suique ipsius sacrificio; quod per æternum Spiritum Deo semel obtulit, justitiæ Patris plene satisfecit,}\footnote{Rom. v. 19; Heb. ix. 14, 16; x. 14; Eph. v. 2; Rom. iii. 25, 26.}\emph{ ac omnibus ei a Patre datis non modo reconciliationem; verum etiam æternam hæreditatem in regno cœlorum acquisivit.}\footnote{Dan. ix. 24, 26; Col. i. 19, 20; Eph. i. 11, 14; John xvii. 2; Heb. ix. 12, 15.}

VI. Although the work of redemption was not actually wrought by Christ till after his incarnation, yet the virtue, efficacy, and benefits thereof were communicated unto the elect, in all ages successively from the beginning of the world, in and by those promises, types,

VI. \emph{Quamvis redemptionis opus non nisi post incarnationem ejus, a Christo quidem actu effectum fuerit, vis tamen ejus, efficacia, et beneficia per omnia iam inde a mundi primordiis elapsa secula electis sunt communicata, in et per promissiones illas, typos, et sacrifica, }

and sacrifices, wherein he was revealed, and signified to be the seed of the woman which should bruise the serpent's head, and the lamb slain from the beginning of the world, being yesterday and today the same and forever.\footnote{Gal. iv. 4, 5; Gen. iii. 15; Rev. xiii. 8; Heb. xiii. 8.}

\emph{quibus revelatum erat et significatum hunc esse semen illud mulieris, quod contriturum erat serpentis caput, agnumque illum mactatum ab initio mundi; ut qui heri ac hodie idem est et in sempiternum.}\footnote{Gal. iv. 4, 5; Gen. iii. 15; Rev. xiii. 8; Heb. xiii. 8.}

VII. Christ, in the work of mediation, acteth according to both natures; by each nature doing that which is proper to itself;\footnote{Heb. ix. 14; 1 Pet. iii. 18.} yet, by reason of the unity of the person, that which is proper to one nature is sometimes, in Scripture, attributed to the person denominated by the other nature.\footnote{Acts xx. 28; John iii. 13; 1 John iii. 16.}

VII. \emph{Christus in opere Mediatorio agit secundum utramque naturam, id agens per utramvis, quod eidem proprium est,}\footnote{Heb. ix. 14; 1 Pet. iii. 18.}\emph{ nonnunquam tamen fit propter personæ unitatem ut quod uni naturæ proprium est, personæ ab altera natura denominatæ in Scriptura tribuatur.}\footnote{Acts xx. 28; John iii. 13; 1 John iii. 16.}

VIII. To all those for whom Christ hath purchased redemption he doth certainly and effectually apply and communicate the same;\footnote{John vi. 37, 39; x. 15, 16.} making intercession for them,\footnote{1 John ii. 1, 2; Rom. viii. 34.} and revealing unto them, in and by the Word, the mysteries of salvation;\footnote{John xv. 13, 15; Eph. i. 7–9; John xvii. 6.} effectually persuading them by his Spirit to believe and obey; and governing their hearts by his Word and Spirit;\footnote{John xiv. 16; Heb. xii. 2; 2 Cor. iv. 13; Rom. viii. 9, 14; xv. 18, 19; John xvii. 17.} overcoming all their enemies by his almighty power and wisdom, in such manner and ways as are most consonant to his wonderful and unsearchable dispensation.\footnote{Psa. cx. 1; 1 Cor. xv. 25, 26; Mal. iv. 2, 3; Col. ii. 15.}

VIII. \emph{Pro quibus Christus redemptionem acquisivit, iis omnibus certo quidem ac efficaciter eam applicat impertitque,}\footnote{John vi. 37, 39; x. 15, 16.}\emph{ pro eis intercedens,}\footnote{1 John ii. 1, 2; Rom. viii. 34.}\emph{ eisque in verbo et per verbum revelans mysterium salutis,}\footnote{John xv. 13, 15; Eph. i. 7–9; John xvii. 6.}\emph{ per Spiritum suum eis ut credere velint ac obedire persuadens efficaciter,}\footnote{John xiv. 16; Heb. xii. 2; 2 Cor. iv. 13; Rom. viii. 9, 14; xv. 18, 19; John xvii. 17.}\emph{ eorumque gubernans corda verbo suo spirituque; sed et vi sua omnipotenti, ac sapientia debellans omnes eorum hostes, iis autem modis mediisque quæ admirabili et inscrutabili ejus dispensationi sunt maxime consentanea.}\footnote{Psa. cx. 1; 1 Cor. xv. 25, 26; Mal. iv. 2, 3; Col. ii. 15.}

\xsection{Chapter IX. Of Free-will}

I. God hath endued the will of man with that natural liberty, that \footnote{[Am. ed. inserts \emph{it}.]}is neither forced nor by any absolute necessity of nature determined to good or evil.\footnote{Matt. xvii. 12; James i. 14; Deut. xxx. 19; [Am. ed. John v. 40].}

I. \emph{Eam humanæ voluntati naturalem Deus indidit libertatem, ut nec cogatur unquam, neque absoluta ulla naturæ necessitate ad bonum aut malum determinetur.}\footnote{Matt. xvii. 12; James i. 14; Deut. xxx. 19; [Am. ed. John v. 40].}

II. Man, in his state of innocency, had freedom and power to will and to do that which is good and well-pleasing to God, but yet mutably, so that he might fall from it.\footnote{Gen. ii. 16, 17; iii. 6.}

II. \emph{Homo in statu innocentiæ libertatem habuit ac potentiam, quod bonum erat Deoque gratum volendi agendique;}\footnote{Eccles. vii. 29; Gen. i. 26.}\emph{ mutabiliter tamen, ita ut illa potuerit excidere.}\footnote{Gen. ii. 16, 17; iii. 6.}

III. Man, by his fall into a state of sin, hath wholly lost all ability of will to any spiritual good accompanying salvation;\footnote{Rom. v. 6; viii. 7; John xv. 5.} so as a natural man, being altogether averse from that good,\footnote{Rom. iii. 10, 12.} and dead in sin,\footnote{Eph. ii. 1, 5; Col. ii. 13.} is not able, by his own strength, to convert himself, or to prepare himself thereunto.\footnote{John vi. 44, 65; 1 Cor. ii. 14; Eph. ii. 2–5; Titus iii. 3–5.}

III. \emph{Homo per lapsum suum in statum peccati, potentiam omnem quam habuerat voluntas ejus ad bonum aliquod spirituals et saluti contiguum amisit penitus;}\footnote{Rom. v. 6; viii. 7; John xv. 5.}\emph{ adeo sane ut naturalis homo, utpote ab ejusmodi bono abhorrens prorsus,}\footnote{Rom. iii. 10, 12.}\emph{ ac in peccato mortuus,}\footnote{Eph. ii. 1, 5; Col. ii. 13.}\emph{ non possit unquam suis ipsius viribus convertere semet, sed ne quidem ad conversionem se vel præparare.}\footnote{John vi. 44, 65; 1 Cor. ii. 14; Eph. ii. 2–5; Titus iii. 3–5.}

IV. When God converts a sinner, and translates him into the state of grace, he freeth him from his natural bondage under sin,\footnote{Col. i. 13; John viii. 34, 36.} and by his grace alone enables him freely to will and to do that which is spiritually good;\footnote{Phil. ii. 13; Rom. vi. 18, 22.} yet so

IV. \emph{Quandocunque Deus convertit ac in statum gratiæ transfert peccatorem, eundem eximit naturali sua sub peccato servitute,}\footnote{Col. i. 13; John viii. 34, 36.}\emph{ solaque gratia sua potentem reddit ad spirituale bonum volendum præstandumque;}\footnote{Phil. ii. 13; Rom. vi. 18, 22.}\emph{ ita tamen ut propter }

as that, by reason of his remaining corruption, he doth not perfectly, nor only, will that which is good, but doth also will that which is evil.\footnote{Gal. v. 17; Rom. vii. 15, 18, 19, 21, 23.}

\emph{manentem adhuc in eo corruptionem, bonum nec perfecte velit; neque id tantummodo, verum etiam quandoque malum.}\footnote{Gal. v. 17; Rom. vii. 15, 18, 19, 21, 23.}

V. The will of man is made perfectly and immutably free to good alone, in the state of glory only.\footnote{Eph. iv. 13; Heb. xii. 23; 1 John iii. 2; Jude 24.}

V. \emph{Voluntas humana perfecte ac immutabiliter libera ad bonum solum redditur non nisi in statu gloriæ.}\footnote{Eph. iv. 13; Heb. xii. 23; 1 John iii. 2; Jude 24.}

\xsection{Chapter X. Of Effectual Calling}

I. All those whom God hath predestinated unto life, and those only, he is pleased, in his appointed and accepted time, effectually to call,\footnote{Rom. viii. 30; xi. 7; Eph. i. 10, 11.} by his Word and Spirit,\footnote{2 Thess. ii. 13, 14; 2 Cor. iii. 3, 6.} out of that state of sin and death, in which they are by nature, to grace and salvation by Jesus Christ;\footnote{Rom. viii. 2; Eph. ii. 1–5; 2 Tim. i. 9, 10.} enlightening their minds, spiritually and savingly, to understand the things of God;\footnote{Acts xxvi. 18; 1 Cor. ii. 10, 12; Eph. i. 17, 18.} taking away their heart of stone, and giving unto them an heart of flesh;\footnote{Ezek. xxxvi. 26.} renewing their wills, and by his almighty power determining them to that which is good,\footnote{Ezek. xi. 19; Phil. ii. 13; Deut. xxx. 6; Ezek. xxxvi. 27.} and effectually drawing them to Jesus Christ;\footnote{Eph. i. 19; John vi. 44, 45.} yet so as they come most freely, being made willing by his grace.\footnote{Cant. i. 4; Psa. cx. 3; John vi. 37; Rom. vi. 16–18.}

I. \emph{Deus quos ad vitam prædestinavit omnes, eosque solos dignatur per verbum suum et spiritum}\footnote{Rom. viii. 30; xi. 7; Eph. i. 10, 11.}\emph{ constituto suo acceptoque tempore vocare efficaciter}\footnote{2 Thess. ii. 13, 14; 2 Cor. iii. 3, 6.}\emph{ e statu illo peccati et mortis in quo sunt natura constituti, ad gratiam ac salutem per Jesum Christum;}\footnote{Rom. viii. 2; Eph. ii. 1–5; 2 Tim. i. 9, 10.}\emph{ idque mentes eorum illuminando, ut modo spirituali et salutari quæ Dei sunt intelligant;}\footnote{Acts xxvi. 18; 1 Cor. ii. 10, 12; Eph. i. 17, 18.} tollendo eorum cor lapideum, donandoque eis cor carneum;\footnote{Ezek. xxxvi. 26.}\emph{ voluntates eorum renovando ac pro potentia sua omnipotente ad bonum determinando,}\footnote{Ezek. xi. 19; Phil. ii. 13; Deut. xxx. 6; Ezek. xxxvi. 27.}\emph{ et ad Jesum Christum trahendo efficaciter;}\footnote{Eph. i. 19; John vi. 44, 45.}\emph{ ita tamen ut illi nihilominus liberrime veniant, volentes nempe facti per illius gratiam.}\footnote{Cant. i. 4; Psa. cx. 3; John vi. 37; Rom. vi. 16–18.}

II. This effectual call is of God's

II. \emph{Efficax hæc vocatio est a sola }

free and special grace alone, not from any thing at all foreseen in man;\footnote{2 Tim. i. 9; Titus iii. 4, 5; Eph. ii. 4, 5, 8, 9; Rom. ix. 11.} who is altogether passive therein, until, being quickened and renewed by the Holy Spirit,\footnote{1 Cor. ii. 14; Rom. viii. 7; Eph. ii. 5.} he is thereby enabled to answer this call, and to embrace the grace offered and conveyed in it.\footnote{John vi. 37; Ezek. xxxvi. 27; John v. 25.}

\emph{Dei gratia:, gratuita illa et speciali; a nulla autem re in homine prævisa;}\footnote{2 Tim. i. 9; Titus iii. 4, 5; Eph. ii. 4, 5, 8, 9; Rom. ix. 11.}\emph{ qui in hoc negotio se habet omnino passive, donec per spiritum sanctum vivificatus ac renovatus,}\footnote{1 Cor. ii. 14; Rom. viii. 7; Eph. ii. 5.}\emph{ potis inde factus sit vocationi huic respondere, gratiamque inibi oblatam et exhibitam amplexari.}\footnote{John vi. 37; Ezek. xxxvi. 27; John v. 25.}

III. Elect infants, dying in infancy, are regenerated and saved by Christ through the Spirit,\footnote{Luke xviii. 15, 16 and Acts ii. 38, 39, and John iii. 3, 5, and 1 John v. 12, and Rom. viii. 9, compared.} who worketh when, and where, and how he pleaseth.\footnote{John iii. 8.} So also are all other elect persons, who are incapable of being outwardly called by the ministry of the Word.\footnote{1 John v. 12; Acts iv. 12.}

III. \emph{Electi infantes in infantia sua morientes regenerantur salvanturque a Christo per spiritum}\footnote{Luke xviii. 15, 16 and Acts ii. 38, 39, and John iii. 3, 5, and 1 John v. 12, and Rom. viii. 9, compared.} \{\emph{qui quando et ubi, et quo sibi placuerit modo operator});\footnote{John iii. 8.}\emph{ sicut et reliqui electi omnes, quotquot externæ vocationis per ministerium verbi sunt incapaces.}\footnote{1 John v. 12; Acts iv. 12.}

IV. Others, not elected, although they may be called by the ministry of the Word,\footnote{Matt. xxii. 14.} and may have some common operations of the Spirit,\footnote{Matt. vii. 22; xiii. 20, 21; Heb. vi. 4, 5.} yet they never truly come unto\footnote{[Am. ed. \emph{to.}]} Christ, and therefore can not be saved:\footnote{John vi. 64–66; viii. 24.} much less can men, not professing the Christian religion, be saved in any other way whatsoever, be they never so diligent to frame their lives according to the light of nature and the law of that religion they do profess;\footnote{Acts iv. 12; John xiv. 6; Eph. ii. 12; John iv. 22; xvii. 3.} and to assert and maintain that they may is

IV. \emph{Alii autem, qui non electi sunt, ut ut verbi ministerio vocari possint,}\footnote{Matt. xxii. 14.}\emph{ communesque nonnullas operationes Spiritus experiri,}\footnote{Matt. vii. 22; xiii. 20, 21; Heb. vi. 4, 5.}\emph{ nunquam tamen vere ad Christum accedunt, proindeque nec salvari possunt.}\footnote{John vi. 64–66; viii. 24.}\emph{ Multo quidem minus poterunt illi, quotquot religionem Christianam non profitentur (summam licet operam navaverint moribus suis ad naturæ lumen, istiusque quam profitentur religionis legem componendis), extra hanc unicam viam salutem unquam obtinere.}\footnote{Acts iv. 12; John xiv. 6; Eph. ii. 12; John iv. 22; xvii. 3.}\emph{ Atque huic quidem contrarium }

very pernicious, and to be detested.\footnote{2 John 9–11; 1 Cor. xvi. 22; Gal. i. 6–8.}

\emph{statuere ac defendere, perniciosum admodum est ac detestandum.}\footnote{2 John 9–11; 1 Cor. xvi. 22; Gal. i. 6–8.}

\xsection{Chapter XI. Of Justification}

I. Those whom God effectually calleth he also freely justifieth;\footnote{Rom. viii. 30; iii. 24.} not by infusing righteousness into them, but by pardoning their sins, and by accounting and accepting their persons as righteous: not for any thing wrought in them, or done by them, but for Christ's sake alone; nor\footnote{[Am. ed. \emph{not.}]} by imputing faith itself, the act of believing, or any other evangelical obedience to them, as their righteousness; but by imputing the obedience and satisfaction of Christ unto them,\footnote{Rom. iv. 5–8; 2 Cor. v. 19, 21; Rom. iii. 22, 24, 25, 27, 28; Titus iii. 5, 7; Eph. i. 7; Jer. xxiii. 6; 1 Cor. i. 30, 31; Rom. v. 17–19.} they receiving and resting on him and his righteousness by faith; which faith they have not of themselves, it is the gift of God.\footnote{Acts x. 44; Gal. ii. 16; Phil. iii. 9; Acts xiii. 38, 39; Eph. ii. 7, 8.}

I. \emph{Quos Deus vocat efficaciter, eosdem etiam gratis justificat,}\footnote{Rom. viii. 30; iii. 24.}\emph{ non quidem justitiam iis infundendo, sed eorum peccata condonando, personasque pro justis reputando atque acceptando; neque id certe propter quicquam aut in iis productum, aut ab iis præstitum, verum Christi solius ergo; eisque ad justitiam non fidem ipsam, non credendi actum, aut aliam quamcunque obedientiam evangelicam, verum obedientiam ac satisfactionem Christi imputando,}\footnote{Rom. iv. 5–8; 2 Cor. v. 19, 21; Rom. iii. 22, 24, 25, 27, 28; Titus iii. 5, 7; Eph. i. 7; Jer. xxiii. 6; 1 Cor. i. 30, 31; Rom. v. 17–19.}\emph{ eum nempe recipientibus, eique ac justitiæ ejus per fidem innitentibus; quam illi fidem ex dono Dei, non a seipsis, habent.}\footnote{Acts x. 44; Gal. ii. 16; Phil. iii. 9; Acts xiii. 38, 39; Eph. ii. 7, 8.}

II. Faith, thus receiving and resting on Christ and his righteousness, is the alone instrument of justification;\footnote{John i. 12; Rom. iii. 28; v. 1.} yet is it not alone in the person justified, but is ever accompanied with all other saving graces, and is no dead faith, but worketh by love.\footnote{James ii. 17, 22, 26; Gal. v. 6.}

II. \emph{Fides hoc modo Christum recipiens, eique innitens ac justitiæ ejus, est justificationis unicum instrumentum;}\footnote{John i. 12; Rom. iii. 28; v. 1.}\emph{ in homine tamen justificato hæc non est solitaria, verum gratiis aliis omnibus salutaribus semper comitata; neque est hæc fides mortua, sed quæ per charitatem operatur.}\footnote{James ii. 17, 22, 26; Gal. v. 6.}

III. Christ, by his obedience and

III. \emph{Qui hunc in modum justificantur, }

death, did fully discharge the debt of all those that are thus justified, and did make a proper, real, and full satisfaction to his Father's justice in their behalf.\footnote{Rom. v. 8–10, 19; 1 Tim. ii. 5, 6; Heb. x. 10, 14; Dan. ix. 24, 26; Isa. liii. 4–6, 10–12.} Yet inasmuch as he was given by the Father for them,\footnote{Rom. viii. 32.} and his obedience and satisfaction accepted in their stead,\footnote{2 Cor. v. 21; Matt. iii. 17; Eph. v. 2.} and both freely, not for any thing in them, their justification is only of free grace;\footnote{Rom. iii. 24; Eph. i. 7.} that both the exact justice and rich grace of God might be glorified in the justification of sinners.\footnote{Rom. iii. 26; Eph. ii. 7.}

\emph{eorum omnium debita Christus per obedientiam suam mortemque prorsus dissolvit; eorumque vice justitiæ Patris sui realem, plenum, et proprie dictam satisfactionem præstitit.}\footnote{Rom. v. 8–10, 19; 1 Tim. ii. 5, 6; Heb. x. 10, 14; Dan. ix. 24, 26; Isa. liii. 4–6, 10–12.}\emph{ Quum tamen non propter in iis quicquam, verum gratuito Pater cum Christum ipsum pro eis dederit,}\footnote{Rom. viii. 32.}\emph{ tum obedientiam ejus ac satisfactionem tanquam eorum loco constituti}\footnote{2 Cor. v. 21; Matt. iii. 17; Eph. v. 2.}\emph{ acceptaverit; omnino a gratia gratuita est eorum justificatio;}\footnote{Rom. iii. 24; Eph. i. 7.}\emph{ Quo nimirum Dei tum accurata justitia tum locuples gratia glorificata foret in justificatione peccatorum.}\footnote{Rom. iii. 26; Eph. ii. 7.}

IV. God did, from all eternity, decree to justify all the elect,\footnote{Gal. iii. 8; 1 Pet. i. 2, 19, 20; Rom. viii. 30.} and Christ did, in the fullness of time, die for their sins, and rise again for their justification:\footnote{Gal. iv. 4; 1 Tim. ii. 6; Rom. iv. 25.} nevertheless, they are not justified until the Holy Spirit doth, in due time, actually apply Christ unto them.\footnote{Col. i. 21, 22; Gal. ii. 16; Titus iii. 4–7.}

IV. \emph{Ab æterno decrevit Deus electos omnes justificare,}\footnote{Gal. iii. 8; 1 Pet. i. 2, 19, 20; Rom. viii. 30.}\emph{ Christusque in temporis plenitudine mortuus est pro eorum peccatis, et in justificationem eorum resurrexit:}\footnote{Gal. iv. 4; 1 Tim. ii. 6; Rom. iv. 25.}\emph{ nihilo minus tamen justificati prius non sunt, quam Christum eis in tempore suo opportuno Spiritus Sanctus actu applicuerit.}\footnote{Col. i. 21, 22; Gal. ii. 16; Titus iii. 4–7.}

V. God doth continue to forgive the sins of those that are justified;\footnote{Matt. vi. 12; 1 John i. 7, 9; ii. 1, 2.} and although they can never fall from the state of justification,\footnote{Luke xxii. 32; John x. 28; Heb. x. 14.} yet they may by their sins fall under God's fatherly displeasure, and not have the light of his countenance restored unto them, until they humble themselves,

V. \emph{Perseverat Deus eorum peccata condonare quos semel justificavit,}\footnote{Matt. vi. 12; 1 John i. 7, 9; ii. 1, 2.}\emph{ quin et etiamsi excidere statu justificationis nunquam possint;}\footnote{Luke xxii. 32; John x. 28; Heb. x. 14.}\emph{ fieri tamen potest ut iræ Dei, paternæ quidem illi, per peccata sua se exponant, nec lumen paterni vultus prius sibi habeant restitutum, quam semet ipsos humiliaverint, }

confess their sins, beg pardon, and renew their faith and repentance.\footnote{Psa. lxxxix. 31–33; li. 7–12; xxxii. 5; Matt. xxvi. 75; 1 Cor. xi. 30, 32; Luke i. 20.}

\emph{peccata agnoverint, imploraverint veniam, fidem denique et pœnitentiam suam renovaverint.}\footnote{Psa. lxxxix. 31–33; li. 7–12; xxxii. 5; Matt. xxvi. 75; 1 Cor. xi. 30, 32; Luke i. 20.}

VI. The justification of believers under the Old Testament was, in all these respects, one and the same with the justification of believers under the New Testament.\footnote{Gal. iii. 9, 13, 14; Rom. iv. 22–24; Heb.xiii. 8.}

VI. \emph{Justificatio fidelium sub Vetere ac Novo}\footnote{Gal. iii. 9, 13, 14; Rom. iv. 22–24; Heb.xiii. 8.}\emph{ Testamento quoad isthæc omnia est una eademque.}\footnote{Gal. iii. 9, 13, 14; Rom. iv. 22–24; Heb.xiii. 8.}

\xsection{Chapter XII. Of Adoption}

All those that are justified God vouchsafeth, in and for his only Son Jesus Christ, to make partakers of the grace of adoption;\footnote{Eph. i. 5; Gal. iv. 4, 5.} by which they are taken into the number, and enjoy the liberties and privileges of the children of God;\footnote{Rom. viii. 17; John i. 12.} have his name put upon them;\footnote{Jer. xiv. 9; 2 Cor. vi. 18; Rev. iii. 12.} receive the Spirit of adoption;\footnote{Rom. viii. 15.} have access to the throne of grace with boldness;\footnote{Eph. iii. 12; Rom. v. 2.} are enabled to cry, Abba, Father;\footnote{Gal. iv. 6.} are pitied,\footnote{Psa. ciii. 13.} protected,\footnote{Prov. xiv. 26.} provided for,\footnote{Matt. vi. 30, 32; 1 Pet. v. 7.} and chastened by him as by a father;\footnote{Heb. xii. 6.} yet never cast off,\footnote{Lam. iii. 31.} but sealed to the day of redemption,\footnote{Eph. iv. 30.} and inherit the promises,\footnote{Heb. vi. 12.} as heirs of everlasting salvation.\footnote{l Pet. i. 3, 4; Heb. i. 14.}

\emph{Deus justificatos omnes dignatur in filio suo unigenito Jesu Christo, et propter eundem participes facere gratiæ Adoptionis;}\footnote{Eph. i. 5; Gal. iv. 4, 5.}\emph{ per quam in numerum filiorum Dei assumuntur, taliumque immunitatibus ac privilegiis potiuntur,}\footnote{Rom. viii. 17; John i. 12.}\emph{ impositum sibi habent nomen Dei,}\footnote{Jer. xiv. 9; 2 Cor. vi. 18; Rev. iii. 12.}\emph{ Spiritum adoptionis accipiunt,}\footnote{Rom. viii. 15.}\emph{ aditum habent ad thronum gratiæ cum confidentia,}\footnote{Eph. iii. 12; Rom. v. 2.}\emph{ potestatem consequuntur clamandi Abba Pater,}\footnote{Gal. iv. 6.}\emph{ commiserationem,}\footnote{Psa. ciii. 13.}\emph{ tutelam,}\footnote{Prov. xiv. 26.}\emph{ et providentiam}\footnote{Matt. vi. 30, 32; 1 Pet. v. 7.}\emph{ sortiuntur; quin et castigationem Dei paternam experiuntur;}\footnote{Heb. xii. 6.}\emph{ nunquam tamen abdicantur,}\footnote{Lam. iii. 31.}\emph{ verum in diem redemptionis consignati}\footnote{Eph. iv. 30.}\emph{ promissiones obtinent hæreditario jure,}\footnote{Heb. vi. 12.}\emph{ ut qui hæredes sunt æternæ salutis.}\footnote{l Pet. i. 3, 4; Heb. i. 14.}

\xsection{Chapter XIII. Of Sanctification}

I. They who are effectually called and regenerated, having a new heart and a new spirit created in them, are further sanctified, really and personally, through the virtue of Christ's death and resurrection,\footnote{1 Cor. vi. 11; Acts xx. 32; Phil. iii. 10; Rom. vi. 5, 6.} by his Word and Spirit dwelling in them;\footnote{John xvii. 17; Eph. v. 26; 2 Thess. ii. 13.} the dominion of the whole body of sin is destroyed,\footnote{Rom. vi. 6, 14.} and the several lusts thereof are more and more weakened and mortified,\footnote{Gal. v. 24; Rom. viii. 13.}and they more and more quickened and strengthened, in all saving graces,\footnote{Col. i. 11; Eph. iii. 16–19.} to the practice of true holiness, without which no man shall see the Lord.\footnote{2 Cor. vii. 1; Heb. xii. 14.}

I. \emph{Quotquot efficaciter vocantur, ac regenerantur, cor novum habentes novumque spiritum in se creatum, sunt virtute mortis et resurrectionis Christi}\footnote{1 Cor. vi. 11; Acts xx. 32; Phil. iii. 10; Rom. vi. 5, 6.}\emph{ per verbum ejus spiritumque in eis inhabitantem}\footnote{John xvii. 17; Eph. v. 26; 2 Thess. ii. 13.}\emph{ ulterius sandificati, realiter quidem ac personaliter: totius corporis peccati dominium in eos destruitur,}\footnote{Rom. vi. 6, 14.}\emph{ ejusque variæ libidines debilitantur indies magis magisque ac mortificantur;}\footnote{Gal. v. 24; Rom. viii. 13.}\emph{ illi interim magis magisque in omni gratia salutari vivificantur et corroborantur indies,}\footnote{Col. i. 11; Eph. iii. 16–19.}\emph{ ad praxim veræ sanctimoniæ, qua quidem destitutus nemo unquam videbit Dominum.}\footnote{2 Cor. vii. 1; Heb. xii. 14.}

II. This sanctification is throughout in the whole man,\footnote{1 Thess. v. 23.} yet imperfect in this life; there abideth still some remnants of corruption in every part,\footnote{1 John i. 10; Rom. vii. 18, 23; Phil. iii. 12.} whence ariseth a continual and irreconcilable war, the flesh lusting against the spirit, and the spirit against the flesh.\footnote{Gal. v. 17; 1 Pet. ii. 11.}

II. \emph{Universalis est hæc et per totum hominem diffusa sanctificatio,}\footnote{1 Thess. v. 23.}\emph{ verum in hac vita est imperfecta nonnullis corruptionis reliquiis adhuc in omni parte remanentibus,}\footnote{1 John i. 10; Rom. vii. 18, 23; Phil. iii. 12.}\emph{ unde bellum exoritur perpetuum et implacabile; hinc carne adversus spiritum, illinc spiritu adversus carnem concupiscente.}\footnote{Gal. v. 17; 1 Pet. ii. 11.}

III. In which war, although the remaining corruption for a time may much prevail,\footnote{Rom. vii. 23.} yet, through the continual supply of strength from the sanctifying Spirit of

III. \emph{In quo quidem bello licet corruptio residua possit aliquandiu prævalere plurimum,}\footnote{Rom. vii. 23.}\emph{ pars tamen regenita, sanctificante Christi spiritu perpetuas ferente suppetias, }

Christ, the regenerate part doth overcome;\footnote{Rom. vi. 14; 1 John v. 4; Eph. iv. 15, 16.} and so the saints grow in grace,\footnote{2 Pet. iii. 18; 2 Cor. iii. 18.} perfecting holiness in the fear of God.\footnote{2 Cor. vii. 1.}

\emph{evadit victrix});\footnote{Rom. vi. 14; 1 John v. 4; Eph. iv. 15, 16.}\emph{ adeoque sancti in gratia crescunt,}\footnote{2 Pet. iii. 18; 2 Cor. iii. 18.}\emph{ sanctitatem in timore Domini perficientes.}\footnote{2 Cor. vii. 1.}

\xsection{Chapter XIV. Of Saving Faith}

I. The grace of faith, whereby the elect are enabled to believe to the saving of their souls,\footnote{Heb. x. 39.} is the work of the Spirit of Christ in their hearts,\footnote{2 Cor. iv. 13; Eph. i. 17–19; ii. 8.} and is ordinarily wrought by the ministry of the Word;\footnote{Rom. x. 14, 17.} by which also, and by the administration of the sacraments and prayer, it is increased and strengthened.\footnote{1 Pet. ii. 2; Acts xx. 32; Rom. iv. 11; Luke xvii. 5; Rom. i. 16, 17.}

I. \emph{Gratia Fidei, qua electi credere valent ad animarum suarum salutem,}\footnote{Heb. x. 39.}\emph{ Spiritus Christi opus est in eorum cordibus operantis,}\footnote{2 Cor. iv. 13; Eph. i. 17–19; ii. 8.}\emph{ effectum plerumque verbi Dei ministerio,}\footnote{Rom. x. 14, 17.}\emph{ quo eodem etiam, ut et administratione Sacramentorum atque oratione robur ei accedit ac incrementum.}\footnote{1 Pet. ii. 2; Acts xx. 32; Rom. iv. 11; Luke xvii. 5; Rom. i. 16, 17.}

II. By this faith a Christian believeth to be true whatsoever is revealed in the Word, for the authority of God himself speaking therein;\footnote{John iv. 42; 1 Thess. ii. 13; 1 John v. 10; Acts xxiv. 14.} and acteth differently upon that which each particular passage thereof containeth; yielding obedience to the commands,\footnote{Rom. xvi. 26.} trembling at the threatenings,\footnote{Isa. lxvi. 2.} and embracing the promises of God for this life and that which is to come.\footnote{Heb. xi. 13; 1 Tim. iv. 8.} But the principal acts of saving faith are accepting, receiving, and resting upon Christ alone

II. \emph{Hac Fide credit Christianus verum esse quicquid in verbo revelatur, propter authoritatem ipsius inibi loquentis Dei;}\footnote{John iv. 42; 1 Thess. ii. 13; 1 John v. 10; Acts xxiv. 14.}\emph{ et varie quidem in illud agit tum obsequendo mandatis,}\footnote{Rom. xvi. 26.}\emph{ tum ad minas contremiscens,}\footnote{Isa. lxvi. 2.}\emph{ tum etiam promissa Dei, seu præsentem hanc vitam seu futuram spectent, amplexando,}\footnote{Heb. xi. 13; 1 Tim. iv. 8.}\emph{ pro varia nempe ratione illarum rerum, quæ in singulis verbi partibus continentur. Verum fidei salvificæ actus illi sunt præcipui, Christi acceptatio et receptio, in eumque solum }

for justification, sanctification, and eternal life, by virtue of the covenant of grace.\footnote{John i. 12; Acts xvi. 31; Gal. ii. 20; Acts xv. 11.}

\emph{recumbentia pro justificatione, sanctificatione, ipsaque adeo vita æterna, virtute fœderis gratiæ consequendis.}\footnote{John i. 12; Acts xvi. 31; Gal. ii. 20; Acts xv. 11.}

III. This faith is different in degrees, weak or strong;\footnote{Heb. v. 13, 14; Rom. iv. 19, 20; Matt. vi. 30; viii. 10.} may be often and many ways assailed and weakened, but gets the victory;\footnote{Luke xxii. 31, 32; Eph. vi. 16; 1 John v. 4, 5.} growing up in many to the attainment of a full assurance through Christ,\footnote{Heb. vi. 11, 12; x. 22; Col. ii. 2.} who is both the author and finisher of our faith.\footnote{Heb. xii. 2.}

III. \emph{Fides hæc pro diversis ejus gradibus debilior est aut fortior;}\footnote{Heb. v. 13, 14; Rom. iv. 19, 20; Matt. vi. 30; viii. 10.}\emph{ impugnari quidem sæpenumero multisque modis ac debilitari potest, non ita tamen quin victrix evadat;}\footnote{Luke xxii. 31, 32; Eph. vi. 16; 1 John v. 4, 5.}\emph{ et quidem in multis ad plenum usque certitudinem per Christum adolescit,}\footnote{Heb. vi. 11, 12; x. 22; Col. ii. 2.}\emph{ qui fidei nostræ idem author est et consummator.}\footnote{Heb. xii. 2.}

\xsection{Chapter XV. Of Repentance unto Life}

I. Repentance unto life is an evangelical grace,\footnote{Zech. xii. 10; Acts xi. 18.} the doctrine whereof is to be preached by every minister of the gospel, as well as that of faith in Christ.\footnote{Luke xxiv. 47; Mark i. 15; Acts xx. 21.}

I. \emph{Resipiscentia ad vitam est gratia Evangelica,}\footnote{Zech. xii. 10; Acts xi. 18.}\emph{ cuius quidem doctrina pariter ac illa de fide in Christum est a singulis ministris Evangelii prædicanda.}

II. By it a sinner, out of the sight and sense, not only of the danger, but also of the filthiness and odiousness of his sins, as contrary to the holy nature and righteous law of God, and upon the apprehension of his mercy in Christ to such as are penitent, so grieves for and hates his sins as to turn from them all unto God,\footnote{Ezek. xviii. 30, 31; xxxvi. 31; Isa. xxx. 22; Psa. li. 4; Jer. xxxi. 18, 19; Joel ii. 12, 13; Amos v. 15; Psa. cxix. 128; 2 Cor. vii. 11.} purposing and endeavoring

II. \emph{Per eam peccator ex inspectu sensuque non solum periculi verum etiam turpitudinis, ac naturæ peccatorum suorum prorsus abominandæ.}\footnote{Luke xxiv. 47; Mark i. 15; Acts xx. 21.}\emph{ utpote sanctæ Dei naturæ, justæque legi adversantium, atque e perspecta ejus erga pœnitentes in Christo misericordia, ita peccata sua deflet ac detestatur, ut ab eis omnibus ad Deum convertatur}\footnote{Ezek. xviii. 30, 31; xxxvi. 31; Isa. xxx. 22; Psa. li. 4; Jer. xxxi. 18, 19; Joel ii. 12, 13; Amos v. 15; Psa. cxix. 128; 2 Cor. vii. 11.}\emph{ cum proposito conatuque in cunctis mandatorum }

to walk with him in all the ways of his commandments.\footnote{Psa. cxix. 6, 59, 106; Luke i. 6; 2 Kings xxiii. 25.}

\emph{ejus viis cum eodem ambulandi.}\footnote{Psa. cxix. 6, 59, 106; Luke i. 6; 2 Kings xxiii. 25.}

III. Although repentance be not to be rested in as any satisfaction for sin, or any cause of the pardon thereof,\footnote{Ezek. xxxvi. 31, 32; xvi. 61–63.} which is the act of God's free grace in Christ;\footnote{Hos. xiv. 2, 4; Rom. iii. 24; Eph. i. 7.} yet is it of such necessity to all sinners that none may expect pardon without it.\footnote{Luke xiii. 3, 5; Acts xvii. 30, 31.}

III. \emph{Etsi resipiscentiæ nobis fidendum non sit, ac si ea esset ulla aut pro peccatis satisfactio, aut causa remissionis peccatorum}\footnote{Ezek. xxxvi. 31, 32; xvi. 61–63.} (\emph{qui gratiæ Dei in Christo gratuitæ actus est}),\footnote{Hos. xiv. 2, 4; Rom. iii. 24; Eph. i. 7.}\emph{ est nihilominus cunctis peccatoribus usque adeo necessaria, ut sine ea nulla cuivis unquam remissio sit expectanda.}\footnote{Luke xiii. 3, 5; Acts xvii. 30, 31.}

IV. As there is no sin so small but it deserves damnation,\footnote{Rom. vi. 23; v. 12; Matt. xii. 36.} so there is no sin so great that it can bring damnation upon those who truly repent.\footnote{Isa. lv. 7; Rom. viii. 1; Isa. i. 16, 18.}

IV. \emph{Quemadmodum nullum est peccatum adeo exiguum ut damnationem non mereatur,}\footnote{Rom. vi. 23; v. 12; Matt. xii. 36.}\emph{ ita neque magnum adeo peccatum ullum est, ut damnationem inferre possit vere pœnitentibus.}\footnote{Isa. lv. 7; Rom. viii. 1; Isa. i. 16, 18.}

V. Men ought not to content themselves with a general repentance, but it is every man's duty to endeavor to repent of his particular sins particularly.\footnote{Psa. xix. 13; Luke xix. 8; 1 Tim. i. 13, 15.}

V. \emph{In resipiscentia generali acquiescendum non est, verum ad id contendere tenetur quisque, ut singulorum suorum peccatorum quam particularem agat pœnitentiam.}\footnote{Psa. xix. 13; Luke xix. 8; 1 Tim. i. 13, 15.}

VI. As every man is bound to make private confession of his sins to God, praying for the pardon thereof,\footnote{Psa. li. 4, 5, 7, 9, 14; xxxii. 5, 6.} upon which, and the forsaking of them, he shall find mercy;\footnote{Prov. xxviii. 13; 1 John i. 9.} so he that scandalizeth his brother, or the Church of Christ, ought to be willing, by a private or public confession and sorrow for his sin, to declare his repentance

VI. \emph{Quemadmodum autem tenetur quivis peccata sua Deo privatim confiteri, et pro remissione illorum precibus contendere:}\footnote{Psa. li. 4, 5, 7, 9, 14; xxxii. 5, 6.} (\emph{quod si præstiterit et peccata simul dereliquerit, misericordiam consequetur})\footnote{Prov. xxviii. 13; 1 John i. 9.}\emph{ ita qui fratri suo, aut Ecclesiæ Christi, scandalo fuerit, promptus et paratus esse debet qua confessione sive privata, sive etiam publica, qua de peccatis }

to those that are offended,\footnote{James v. 16; Luke xvii. 3, 4; Josh. vii. 19; Psa. li. throughout.} who are thereupon to be reconciled to him, and in love to receive him.\footnote{2 Cor. ii. 8; [Amer. ed. Gal. vi. 1, 2].}

\emph{suis dolore, resipiscentiam suam eis quibus offendiculo fuerit declarare,}\footnote{James v. 16; Luke xvii. 3, 4; Josh. vii. 19; Psa. li. throughout.}\emph{ quo præstito illi redire cum eo in gratiam debent, eumque denuo cum charitate recipere.}\footnote{2 Cor. ii. 8; [Amer. ed. Gal. vi. 1, 2].}

\xsection{Chapter XVI. Of Good Works}

I. Good works are only such as God hath commanded in his holy Word,\footnote{Micah vi. 8; Rom. xii. 2; Heb. xiii. 21.} and not such as, without the warrant thereof, are devised by men out of blind zeal, or upon any pretense of good intention.\footnote{Matt. xv. 9; Isa. xxix. 13; 1 Pet. i. 18; Rom. x. 2; John xvi. 2; 1 Sam. xv. 21–23.}

I. \emph{Bona opera ea tantum sunt quæ in verbo suo sancto præcepit Deus;}\footnote{Micah vi. 8; Rom. xii. 2; Heb. xiii. 21.}\emph{ minime autem ea quæ absque ulla illius authoritate, sunt ab hominibus excogitata, sive e cæco zelo id factum fuerit, seu bonæ intentionis prætextu quoviscunque.}\footnote{Matt. xv. 9; Isa. xxix. 13; 1 Pet. i. 18; Rom. x. 2; John xvi. 2; 1 Sam. xv. 21–23.}

II. These good works, done in obedience to God's commandments, are the fruits and evidences of a true and lively faith;\footnote{James ii. 18, 22.} and by them believers manifest their thankfulness,\footnote{Psa. cxvi. 12, 13; 1 Pet. ii. 9.} strengthen their assurance,\footnote{1 John ii. 3, 5; 2 Pet i. 5–10.} edify their brethren,\footnote{2 Cor. ix. 2; Matt. v. 16.} adorn the profession of the gospel,\footnote{Tit. ii. 5, 9–12; 1 Tim. vi. 1.} stop the mouths of the adversaries,\footnote{1 Pet. ii. 15.} and glorify God,\footnote{1 Pet. ii. 12; Phil. i. 11; John xv. 8.} whose workmanship they are, created in Christ Jesus thereunto,\footnote{Eph. ii. 10.} that, having their fruit unto holiness, they may have the end, eternal life.\footnote{Rom. vi. 22.}

II. \emph{Bona hæc opera e conscientia mandatorum Dei præstita vivæ veræque fidei fructus sunt ac evidentiæ;}\footnote{James ii. 18, 22.}\emph{ per hæc fideles gratitudinem suam manifestant,}\footnote{Psa. cxvi. 12, 13; 1 Pet. ii. 9.}\emph{ de salute certitudinem suam augent,}\footnote{1 John ii. 3, 5; 2 Pet i. 5–10.}\emph{ fratres suos ædificant,}\footnote{2 Cor. ix. 2; Matt. v. 16.}\emph{ Evangelii professionem ornant,}\footnote{Tit. ii. 5, 9–12; 1 Tim. vi. 1.}\emph{ obturant ora adversantibus,}\footnote{1 Pet. ii. 15.}\emph{ ac Deum denique glorificant,}\footnote{1 Pet. ii. 12; Phil. i. 11; John xv. 8.}\emph{ cuius opificium sunt in Jesu Christo ad hæc creati,}\footnote{Eph. ii. 10.}\emph{ quo fructum habentes ad sanctimoniam, finem consequantur æternam vitam.}\footnote{Rom. vi. 22.}

III. Their ability to do good

III. \emph{Quod bonis operibus idonei }

works is not at all of themselves, but wholly from the Spirit of Christ.\footnote{John xv. 4–6; Ezek. xxxvi. 26, 27.} And that they may be enabled thereunto, besides the graces they have already received, there is required an actual influence of the same Holy Spirit to work in them to will and to do of his good pleasure;\footnote{Phil. ii. 13; iv. 13; 2 Cor. iii. 5.} yet are they not hereupon to grow negligent, as if they were not bound to perform any duty unless upon a special motion of the Spirit; but they ought to be diligent in stirring up the grace of God that is in them.\footnote{Phil. ii. 12; Heb. vi. 11, 12; 2 Pet. i. 3, 5, 10, 11; Isa. lxiv. 7; 2 Tim. i. 6; Acts xxvi. 6, 7; Jude 20, 21.}

\emph{sint præstandis omnino id a spiritu Christi est, nullatenus autem e seipsis.}\footnote{John xv. 4–6; Ezek. xxxvi. 26, 27.}\emph{ Et quo eis præstandis pares fiant, prater habitus gratiæ iam infusos, ejusdem Spiritus sancti actualis porro requiritur influentia, qua nempe in iis operetur tum velle tum etiam efficere pro suo ipsius beneplacito:}\footnote{Phil. ii. 13; iv. 13; 2 Cor. iii. 5.}\emph{ sed neque tamen iis proinde socordiæ sese licet permittere; ac si nisi specialiter eos excitante Spiritu ad nulla pietatis officia præstanda tenerentur; verum sedulam debent navare operam sustitandæ illi quæ in iis est divinæ gratiæ.}\footnote{Phil. ii. 12; Heb. vi. 11, 12; 2 Pet. i. 3, 5, 10, 11; Isa. lxiv. 7; 2 Tim. i. 6; Acts xxvi. 6, 7; Jude 20, 21.}

IV. They who in their obedience attain to the greatest height which is possible in this life, are so far from being able to supererogate and to do more than God requires, as\footnote{[Amer. ed. omits \emph{as.}]} that they fall short of much which in duty they are bound to do.\footnote{Luke xvii. 10; Neh. xiii. 22; Job ix. 2, 3; Gal. v. 17.}

IV. \emph{Qui gradum obedientiæ summum quidem in hac vita possibilem assequuntur, tantum abest ut supererogare quicquam possint ac plus præstare quam quod Deus requisiverit, ut multum sane subsidant infra illud, quod ex officio præstare obligantur.}\footnote{Luke xvii. 10; Neh. xiii. 22; Job ix. 2, 3; Gal. v. 17.}

V. We can not, by our best works, merit pardon of sin, or eternal life at the hand of God, by reason of the great disproportion that is between them and the glory to come, and the infinite distance that is between us and God, whom by them we can neither profit nor satisfy

V. \emph{Peccatorum veniam, aut vitam æternam de Deo mereri non valemus, ne optimis quidem operibus nostris; cum propter summam illam inter ea et futuram gloriam disparitatem; tum etiam propter infinitam distantiam quæ inter nos ac Deum intercedit; cui nos per illa nec prodesse }

for the debt of our former sins;\footnote{Rom. iii. 20; iv. 2, 4, 6; Eph. ii. 8, 9; Titus iii. 5–7; Rom. viii. 18; Psa. xvi. 2; Job xxii. 2, 3; xxxv. 7, 8.} but when we have done all we can, we have done but our duty, and are unprofitable servants;\footnote{Luke xvii. 10.} and because, as they are good, they proceed from his Spirit;\footnote{Gal. v. 22, 23.} and as they are wrought by us, they are defiled and mixed with so much weakness and imperfection that they can not endure the severity of God's judgment.\footnote{Isa. lxiv. 6; Gal. v. 17; Rom. vii. 15, 18; Psa. cxliii. 2; cxxx. 3.}

\emph{quicquam possumus, neque pro antecedentium peccatorum nostrorum debito satisfacere;}\footnote{Rom. iii. 20; iv. 2, 4, 6; Eph. ii. 8, 9; Titus iii. 5–7; Rom. viii. 18; Psa. xvi. 2; Job xxii. 2, 3; xxxv. 7, 8.}\emph{ verum cum quantum possumus fecerimus, non nisi quod debemus præstiterimus, ac servi inutiles futuri sumus;}\footnote{Luke xvii. 10.}\emph{ tum denique quoniam a spiritu Dei in quantum bona sunt proficiscuntur,}\footnote{Gal. v. 22, 23.}\emph{ ita vero sunt coinquinata, tantumque imperfectionis ac infirmitas admistum habent, prout a nobis efficiuntur, ut strictum Dei judicium non sint ferendo.}\footnote{Isa. lxiv. 6; Gal. v. 17; Rom. vii. 15, 18; Psa. cxliii. 2; cxxx. 3.}

VI. Yet notwithstanding, the persons of believers being accepted through Christ, their good works also are accepted in him,\footnote{Eph. i. 6; 1 Pet. ii. 5; Exod. xxviii. 38; Gen. iv. 4 with Heb. xi. 4.} not as though they were in this life wholly unblamable and unreprovable in God's sight;\footnote{Job ix. 20; Psa. cxliii. 2.} but that he, looking upon them in his Son, is pleased to accept and reward that which is sincere, although accompanied with many weaknesses and imperfections.\footnote{Heb. xiii. 20, 21; 2 Cor. viii. 12; Heb. vi. 10; Matt. xxv. 21, 23.}

VI. \emph{Nihilominus tamen acceptis in gratiam per Christum fidelium personis, eorum etiam opera bona per eundem accepta sunt;}\footnote{Eph. i. 6; 1 Pet. ii. 5; Exod. xxviii. 38; Gen. iv. 4 with Heb. xi. 4.}\emph{ non quod in hac vita sint omnis culpæ prorsus immunia, quæque in conspectu Dei nullam reprehensionem mereantur;}\footnote{Job ix. 20; Psa. cxliii. 2.}\emph{ verum quod illa respiciens in filio suo Deus, quod sincerum est, utcunque multis infirmitatibus ac imperfectionibus involutum, acceptare dignetur ac remunerari.}\footnote{Heb. xiii. 20, 21; 2 Cor. viii. 12; Heb. vi. 10; Matt. xxv. 21, 23.}

VII. Works done by unregenerate men, although for the matter of them they may be things which God commands, and of good use both to themselves and others;\footnote{2 Kings x. 30, 31; 1 Kings xxi. 27, 29; Phil. i. 15, 16, 18.}

VII. \emph{Opera nondum regenitorum, licet, quoad materiam præcepto divino conformia esse possint, sibique ipsis et aliis item utilia;}\footnote{2 Kings x. 30, 31; 1 Kings xxi. 27, 29; Phil. i. 15, 16, 18.}\emph{ cum tamen neque a corde profluant per fidem }

yet because they proceed not from a heart purified by faith,\footnote{Gen. iv. 3–5 with Heb. xi. 4, 6.} nor are done in a right manner, according to the Word,\footnote{1 Cor. xiii. 3; Isa. i. 12.} nor to a right end, the glory of God;\footnote{Matt. vi. 2, 5, 16.} they are therefore sinful, and can not please God, or make a man meet to receive grace from God.\footnote{Hag. ii. 14; Titus i. 15; Amos v. 21, 22; Hos. i. 4; Rom. ix. 16; Titus iii. 5.} And yet their neglect of them is more sinful and displeasing unto God.\footnote{Psa. xiv. 4; xxxvi. 3; Job xxi. 14, 15; Matt. xxv. 41–45; xxiii. 23.}

\emph{depurato,}\footnote{Gen. iv. 3–5 with Heb. xi. 4, 6.}\emph{ nec secundum verbum eo quo par est præstentur modo,}\footnote{1 Cor. xiii. 3; Isa. i. 12.}\emph{ sed neque ad finem debitum, Dei nempe gloriam, destinentur;}\footnote{Matt. vi. 2, 5, 16.}\emph{ sunt proinde peccata, nec Deo grata esse possunt, nec reddere quenquam valent idoneum ad gratiam a Deo recipiendum.}\footnote{Hag. ii. 14; Titus i. 15; Amos v. 21, 22; Hos. i. 4; Rom. ix. 16; Titus iii. 5.}\emph{ Ejusmodi tamen operum neglectu, gravius quidem illi peccant Deumque offendunt vehementius.}\footnote{Psa. xiv. 4; xxxvi. 3; Job xxi. 14, 15; Matt. xxv. 41–45; xxiii. 23.}

\xsection{Chapter XVII. Of the Perseverance of the Saints}

I. They whom God hath accepted in his Beloved, effectually called and sanctified by his Spirit, can neither totally nor finally fall away from the state of grace; but shall certainly persevere therein to the end, and be eternally saved.\footnote{Phil. i. 6; 2 Pet. i. 10; John x. 28, 29; 1 John iii. 9; 1 Pet. i. 5, 9; [Am. ed. Job xvii. 9].}

I. \emph{Quotquot Deus in dilecto suo acceptavit, vocavit efficaciter ac per Spiritum suum sanctificavit, non possunt illi statu gratiæ aut finaliter excidere aut totaliter; verum in eo ad finem usque certo perseverabunt, ac salutem æternam consequentur.}\footnote{Phil. i. 6; 2 Pet. i. 10; John x. 28, 29; 1 John iii. 9; 1 Pet. i. 5, 9; [Am. ed. Job xvii. 9].}

II. This perseverance of the saints depends, not upon their own freewill, but upon the immutability of the decree of election, flowing from the free and unchangeable love of God the Father;\footnote{2 Tim. ii. 18, 19; Jer. xxxi. 3.} upon the efficacy of the merit and intercession of Jesus Christ;\footnote{Heb. x. 10, 14; xiii. 20, 21; ix. 12–15; Rom. viii. 33, to the end; John xvii. 11, 24; Luke xxii. 32; Heb. vii. 25.} the abiding

II. \emph{hæc autem sanctorum perseverantia, non pendet a libero ipsorum arbitrio, verum a decreti electionis immutabilitate} (\emph{quod ex amore Dei Patris fluxit, gratuito illo ac immutabili}),\footnote{2 Tim. ii. 18, 19; Jer. xxxi. 3.}\emph{ a meriti Jesu Christi ac intercessionis efficacia,}\footnote{Heb. x. 10, 14; xiii. 20, 21; ix. 12–15; Rom. viii. 33, to the end; John xvii. 11, 24; Luke xxii. 32; Heb. vii. 25.}\emph{ a Spiritus et seminis Dei in iis permansione;}\footnote{John xiv. 16, 17; 1 John ii. 27; iii. 9}

of the Spirit and of the seed of God within them;\footnote{John xiv. 16, 17; 1 John ii. 27; iii. 9.} and the nature of the covenant of grace:\footnote{Jer. xxxii. 40; [Am. ed. Heb. viii. 10–12].} from all which ariseth also the certainty and infallibility thereof.\footnote{John x. 28; 2 Thess. iii. 3; 1 John ii. 19; [Am. ed. 1 Thess. v. 23, 24].}

a natura denique fœderis gratiæ;\footnote{Jer. xxxii. 40; [Am. ed. Heb. viii. 10–12].} e quibus omnibus etiam emergit certitudo ejusdem et infallibilitas.\footnote{John x. 28; 2 Thess. iii. 3; 1 John ii. 19; [Am. ed. 1 Thess. v. 23, 24].}

III. Nevertheless they may, through the temptations of Satan and of the world, the prevalency of corruption remaining in them, and the neglect of the means of their preservation, fall into grievous sins;\footnote{Matt. xxvi. 70, 72, 74.} and for a time continue therein:\footnote{Psa. li. title and verse 14; [Am. ed. 2 Sam. xii. 9, 13].} whereby they incur God's displeasure,\footnote{Isa. lxiv. 5, 7, 9; 2 Sam. xi. 27.} and grieve his Holy Spirit;\footnote{Eph. iv. 30.} come to be deprived of some measure of their graces and comforts;\footnote{Psa. li. 8, 10, 12; Rev. ii. 4; Cant. v. 2, 3, 4, 6.} have their hearts hardened,\footnote{Isa. xxxvi. 17; Mark vi. 52; xvi. 14; [Am. ed. Psa. xcv. 8].} and their consciences wounded;\footnote{Psa. xxxii. 3, 4; li. 8.} hurt and scandalize others,\footnote{2 Sam. xii. 14.} and bring temporal judgments upon themselves.\footnote{Psa. lxxxix. 31, 32; 1 Cor. xi. 32.}

III. \emph{Nihilo tamen minus fieri potest ut iidem illi, qua Satanæ mundique tentatione, qua manentis adhuc in iis corruptionis prævalentia, et neglectu mediorum conservationis suæ, in peccata gravia incidant,}\footnote{Matt. xxvi. 70, 72, 74.}\emph{ in eisque ad tempus commorentur;}\footnote{Psa. li. title and verse 14; [Am. ed. 2 Sam. xii. 9, 13].}\emph{ unde iram Dei sibi ipsis contrahunt,}\footnote{Isa. lxiv. 5, 7, 9; 2 Sam. xi. 27.}\emph{; ejusque Spiritum Sanctum contristant,}\footnote{Eph. iv. 30.}\emph{ gratias suas et consolationes quadantenus et quoad gradus nonnullos amittunt,}\footnote{Psa. li. 8, 10, 12; Rev. ii. 4; Cant. v. 2, 3, 4, 6.}\emph{ corda sibi habent indurata,}\footnote{Isa. xxxvi. 17; Mark vi. 52; xvi. 14; [Am. ed. Psa. xcv. 8].}\emph{ et vulneratas conscientias;}\footnote{Psa. xxxii. 3, 4; li. 8.}\emph{ aliis nocumento sunt et offendiculo,}\footnote{2 Sam. xii. 14.}\emph{ sibimet ipsis denique accersunt judicia Dei temporalia.}\footnote{Psa. lxxxix. 31, 32; 1 Cor. xi. 32.}

\xsection{Chapter XVIII. Of the Assurance of Grace and Salvation}

I. Although hypocrites and other unregenerate men may vainly deceive themselves with false hopes and carnal presumptions of being in the favor of God and estate of salvation,\footnote{Job viii. 13, 14; Micah iii. 11; Deut. xxix. 19; John viii. 41.} which hope of theirs

I. \emph{Quamvis fieri potest ut hypocritæ aliique homines non regeniti spe vana falsisque (pro corruptæ naturæ more) opinionibus præsumptis, se decipiant, favorem Dei, statumque salutis sibi falso arrogantes;}\footnote{Job viii. 13, 14; Micah iii. 11; Deut. xxix. 19; John viii. 41.}\emph{ quæ illorum }

shall perish:\footnote{Matt. vii. 22, 23; [Am. ed. Job viii. 13].} yet such as truly believe in the Lord Jesus, and love him in sincerity, endeavoring to walk in all good conscience before him, may in this life be certainly assured that they are in a state of grace,\footnote{1 John ii. 3; iii. 14, 18, 19, 21, 24; v. 13.} and may rejoice in the hope of the glory of God, which hope shall never make them ashamed.\footnote{Rom. v. 2, 5.}

\emph{spes peribit:}\footnote{Matt. vii. 22, 23; [Am. ed. Job viii. 13].}\emph{ qui tamen in Dominum Jesum vere credunt, eumque sincere diligunt, studentes coram ipso in omni bona conscientia ambulare; evadere possunt in hac vita certi se in statu gratiæ esse constitutos;}\footnote{1 John ii. 3; iii. 14, 18, 19, 21, 24; v. 13.}\emph{ quin etiam lætari possunt spe gloriæ Dei, quæ quidem spes nunquam eos pudefaciet.}\footnote{Rom. v. 2, 5.}

II. This certainty is not a bare conjectural and probable persuasion, grounded upon a fallible hope;\footnote{Heb. vi. 11, 19.} but an infallible assurance of faith, founded upon the divine truth of the promises of salvation,\footnote{Heb. vi. 17, 18.} the inward evidence of those graces unto which these promises are made,\footnote{2 Pet. i. 4, 5, 10, 11; 1 John ii. 3; iii. 14; 2 Cor. i. 12.} the testimony of the Spirit of adoption witnessing with our spirits that we are the children of God:\footnote{Rom. viii. 15, 16.} which Spirit is the earnest of our inheritance, whereby we are sealed to the day of redemption.\footnote{Eph. i. 13, 14; iv. 30; 2 Cor. i. 21, 22.}

II. \emph{Hæc certitudo non est persuasio mere conjecturalis et probabilis, innixa spe fallaci;}\footnote{Heb. vi. 11, 19.}\emph{ verum infallibilis quædam fidei certitudo, fundamentum habens divinam promissionum salutis veritatem;}\footnote{Heb. vi. 17, 18.}\emph{ gratiarum, quibus promissiones illæ fiunt internam evidentiam;}\footnote{2 Pet. i. 4, 5, 10, 11; 1 John ii. 3; iii. 14; 2 Cor. i. 12.}\emph{ testimonium denique spiritus adoptionis una cum spiritibus nostris testificantis nos esse filios Dei;}\footnote{Rom. viii. 15, 16.}\emph{ qui quidem spiritus arrhabo est hæreditatis nostræ, quo in diem redemtionis sigillamur.}\footnote{Eph. i. 13, 14; iv. 30; 2 Cor. i. 21, 22.}

III. This infallible assurance doth not so belong to the essence of faith, but that a true believer may wait long, and conflict with many difficulties before he be partaker of it:\footnote{1 John v. 13; Isa. l. 10; Mark ix. 24; Psa. lxxxviii. throughout; lxxvii. to ver. 12.} yet, being enabled by the Spirit to know the things which are freely given him of God, he

III. \emph{Hæc certitudo infallibilis, non ita spectat essentiam fidei, quin vere fidelis expectare quandoque diutius, et cum variis difficultatibus confligere prius possit, quam illius compos fiat,}\footnote{1 John v. 13; Isa. l. 10; Mark ix. 24; Psa. lxxxviii. throughout; lxxvii. to ver. 12.}\emph{ verum poterit idem ordinariorum usu debito mediorum, absque revelatione ulla extraordinaria }

may, without extraordinary revelation, in the right use of ordinary means, attain thereunto.\footnote{1 Cor. ii. 12; 1 John iv. 13; Heb. vi. 11, 12; Eph. iii. 17–19.} And therefore it is the duty of every one to give all diligence to make his calling and election sure;\footnote{2 Pet. i. 10.} that thereby his heart may be enlarged in peace and joy in the Holy Ghost, in love and thankfulness to God, and in strength and cheerfulness in the duties of obedience, the proper fruits of this assurance:\footnote{Rom. v. 1, 2, 5; Rom. xiv. 17; xv. 13; Eph. i. 3, 4; Psa. iv. 6, 7; cxix. 32.} so far is it from inclining men to looseness.\footnote{1 John ii. 1, 2; Rom. vi. 1, 2; Titus ii. 11, 12, 14; 2 Cor. vii. 1; Rom. viii. 1, 12; 1 John iii. 2, 3; Psa. cxxx. 4; 1 John i. 6, 7.}

\emph{eam adipisci,}\footnote{1 Cor. ii. 12; 1 John iv. 13; Heb. vi. 11, 12; Eph. iii. 17–19.}\emph{ spiritu nempe quæ Deus illi gratuito donaverit cognoscendi facultatem subministrante. Proindeque tenetur quisque, quo vocationem suam sibi et electionem certmn faciat, omnem adhibere diligentiam,}\footnote{2 Pet. i. 10.}\emph{ unde cor suum habeat pace et gaudio in spiritu sancto, in Deum amore et gratitudine, in actibus observantiæ robore et alacritate dilatatum; qui certitudinus huius fructus proprii sunt ac genuini.}\footnote{Rom. v. 1, 2, 5; Rom. xiv. 17; xv. 13; Eph. i. 3, 4; Psa. iv. 6, 7; cxix. 32.}\emph{ Tantum abest ut homines inde ad omnem nequitiam discingantur.}\footnote{1 John ii. 1, 2; Rom. vi. 1, 2; Titus ii. 11, 12, 14; 2 Cor. vii. 1; Rom. viii. 1, 12; 1 John iii. 2, 3; Psa. cxxx. 4; 1 John i. 6, 7.}

IV. True believers may have the assurance of their salvation divers ways shaken, diminished, and intermitted; as, by negligence in preserving of it; by falling into some special sin, which woundeth, the conscience, and grieveth the Spirit; by some sudden or vehement temptation; by God's withdrawing the light of his countenance, and suffering even such as fear him to walk in darkness and to have no light:\footnote{Cant. v. 2, 3, 6; Psa. li. 8, 12, 14; Eph. iv. 30, 31; Psa. lxxvii. 1–10; Matt. xxvi. 69–72; Psa. xxxi. 22; lxxxviii. throughout; Isa. l. 10.} yet are they never utterly destitute of that seed of God, and life of faith, that love of Christ and the brethren, that sincerity of heart and conscience of duty, out

IV. \emph{Certitudo salutis vere fidelibus multifariam concuti potest et imminui imo et quandoque interrumpi; conservandi scilicet eam incuria; lapsu in peccatum aliquod insigne, quod conscientiam vulnerat, spiritumque contristat; tentatione aliqua vehementi ac subitanea; uti etiam Deo vultus sui lumen subducente, ac permittente ut vel illi qui ipsum timent in tenebris ambulent omni prorsus lumine viduati:}\footnote{Cant. v. 2, 3, 6; Psa. li. 8, 12, 14; Eph. iv. 30, 31; Psa. lxxvii. 1–10; Matt. xxvi. 69–72; Psa. xxxi. 22; lxxxviii. throughout; Isa. l. 10.}\emph{ nunquam tamen destituuntur penitus illo Dei semine vitaque fidei. Christi illa fratrumque dilectione, ea sinceritate cordis et pietatis officia præstandi conscientia; unde per }

of which, by the operation of the Spirit, this assurance may in due time be revived,\footnote{1 John iii. 9; Luke xxii. 32; Job xiii. 15; Psa. lxxiii. 15; li. 8, 12; Isa. 1. 10.} and by the which, in the mean time, they are supported from utter despair.\footnote{Micah vii. 7–9; Jer. lii. 40; Isa. liv. 7–10; Psa. xxii. 1; lxxxviii. throughout.}

\emph{operationem spiritus eadem illa certitudo tempestive possit reviviscere:}\footnote{1 John iii. 9; Luke xxii. 32; Job xiii. 15; Psa. lxxiii. 15; li. 8, 12; Isa. 1. 10.}\emph{ quibusque interim ne prorsus in desperationem ruant suffulciuntur.}\footnote{Micah vii. 7–9; Jer. lii. 40; Isa. liv. 7–10; Psa. xxii. 1; lxxxviii. throughout.}

\xsection{Chapter XIX. Of the Law of God}

I. God gave to Adam a law, as a covenant of works, by which he bound him and all his posterity to personal, entire, exact, and perpetual obedience; promised life upon the fulfilling, and threatened death upon the breach of it; and endued him with power and ability to keep it.\footnote{Gen. i. 26, 27, with Gen. ii. 17; Rom. ii. 14, 15; x. 5; v. 12, 19; Gal. iii. 10, 12; Eccles. vii. 29; Job xxviii. 28.}

I. \emph{Deus Adamo legem dedit ut fœdus operum, quo cum illum ipsum tum posteros ejus omnes, ad obedientiam personalem, integram, exquisitam simul et perpetuam obligavit, pollicitus vitam si observarent, violatoribus autem mortem interminatus; eundemque potentia et viribus imbuit, quibus par esset illam observando.}\footnote{Gen. i. 26, 27, with Gen. ii. 17; Rom. ii. 14, 15; x. 5; v. 12, 19; Gal. iii. 10, 12; Eccles. vii. 29; Job xxviii. 28.}

II. This law, after his fall, continued to be a perfect rule of righteousness; and, as such, was delivered by God upon mount Sinai in ten commandments, and written in two tables;\footnote{James i. 25; ii. 8, 10–12; Rom. xiii. 8, 9; Deut. v. 32; x. 4; Exod. xxxiv. 1; [Am. ed. Rom. iii. 19].} the first four commandments containing our duty towards God, and the other six our duty to man.\footnote{Matt. xxii. 37–40; [Am. ed. Exod. xx. 3–18].}

II. \emph{Lex ista post lapsum non desiit esse justitiæ regula perfectissima; quo etiam nomine a Deo est in monte Sinai tradita, tabulis duabus descripta, decem præceptis comprehensa;}\footnote{James i. 25; ii. 8, 10–12; Rom. xiii. 8, 9; Deut. v. 32; x. 4; Exod. xxxiv. 1; [Am. ed. Rom. iii. 19].}\emph{ quorum quatuor prima officium nostrum erga Deum, sex autem reliqua nostrum erga homines officium complectuntur.}\footnote{Matt. xxii. 37–40; [Am. ed. Exod. xx. 3–18].}

III. Beside this law, commonly called moral, God was pleased to give to the people of Israel, as a Church under age, ceremonial laws, containing several typical ordinances,

III. \emph{Præter autem hanc legem, quæ moralis vulgo audit, visum est Deo ut populo Israelitico tanquam Ecclesiæ minorenni leges daret ceremoniales instituta typica multifaria }

partly of worship, prefiguring Christ, his graces, actions, sufferings, and benefits;\footnote{Heb. ix.; x. 1; Gal. iv. 1–3; Col. ii. 17.} and partly holding forth divers instructions of moral duties.\footnote{1 Cor. v. 7; 2 Cor. vi. 17; Jude 23.} All which ceremonial laws are now abrogated under the New Testament.\footnote{Col. ii. 14, 16, 17; Dan. ix. 27; Eph. ii. 15, 16.}

\emph{continentes; partim de cultu, Christi gratias, actiones, perpessiones ac beneficia præfigurantia;}\footnote{Heb. ix.; x. 1; Gal. iv. 1–3; Col. ii. 17.}\emph{ partim autem de moralibus officiis institutiones varias exhibentia.}\footnote{1 Cor. v. 7; 2 Cor. vi. 17; Jude 23.}\emph{ Quæ leges ceremoniales omnes hodie sub novo instrumento sunt abrogatæ.}\footnote{Col. ii. 14, 16, 17; Dan. ix. 27; Eph. ii. 15, 16.}

IV. To them also, as a body politic, he gave sundry judicial laws, which expired together with the state of that people, not obliging any other, now, further than the general equity thereof may require.\footnote{Exod. xxi.; xxii. 1–29; Gen. xlix. 10, with 1 Pet. ii. 13, 14; Matt. v. 17, with vers. 38, 39; 1 Cor. ix. 8–10.}

IV. \emph{Iisdem etiam tanquam corpori politico leges multas dedit judiciales, quæ una cum istius populi politeia expirarunt, nullos hodie alios obligantes supra quod generalis et communis earum æquitas postularit.}\footnote{Exod. xxi.; xxii. 1–29; Gen. xlix. 10, with 1 Pet. ii. 13, 14; Matt. v. 17, with vers. 38, 39; 1 Cor. ix. 8–10.}

V. The moral law doth forever bind all, as well justified persons as others, to the obedience thereof;\footnote{Rom. xiii. 8–10; Eph. vi. 2; 1 John ii. 3, 4, 7, 8; [Am. ed. Rom. iii. 31, and vi. 15].} and that not only in regard of the matter contained in it, but also in respect of the authority of God the Creator who gave it.\footnote{James ii. 10, 11.} Neither doth Christ in the gospel any way dissolve, but much strengthen, this obligation.\footnote{Matt. v. 17–19; James ii. 8; Rom. iii. 31.}

V. \emph{Lex moralis omnes tam justificatos quam alios quosvis perpetuo ligat ad obedientiam illi exhibendam;}\footnote{Rom. xiii. 8–10; Eph. vi. 2; 1 John ii. 3, 4, 7, 8; [Am. ed. Rom. iii. 31, and vi. 15].}\emph{ neque id quidem solummodo vi materiæ quæ in illa continetur, verum etiam virtute authoritatis eandem constituentis creatoris Dei;}\footnote{James ii. 10, 11.}\emph{ neque sane hoc ejus vinculum in evangelio ulla ratione dissolvit Christus, verum idem plurimum confirmavit.}\footnote{Matt. v. 17–19; James ii. 8; Rom. iii. 31.}

VI. Although true believers be not under the law as a covenant of works, to be thereby justified or condemned;\footnote{Rom. vi. 14; Gal. ii. 16; iii. 13; iv. 4, 5; Acts xiii. 39; Rom. viii. 1.} yet is it of great use to them, as well as to others; in that, as a rule of life, informing them of the will of God and their

VI. \emph{Quamvis vere fideles non sint sub lege tanquam sub operum fœdere, unde aut justificari possint aut condemnari:}\footnote{Rom. vi. 14; Gal. ii. 16; iii. 13; iv. 4, 5; Acts xiii. 39; Rom. viii. 1.}\emph{ est tamen ea illis non minus quam aliis vehementer utilis, ut quæ quum sit vitæ norma, illos voluntatem divinam suumque officium }

duty, it directs and binds them to walk accordingly;\footnote{Rom. vii. 12, 22, 25; Psa. cxix. 4–6; 1 Cor. vii. 19; Gal. v. 14, 16, 18–23.} discovering also the sinful pollutions of their nature, hearts, and lives;\footnote{Rom. vii. 7; iii. 20.} so as, examining themselves thereby, they may come to further conviction of, humiliation for, and hatred against sin;\footnote{James i. 23–25; Rom. vii. 9, 14, 24.} together with a clearer sight of the need they have of Christ, and the perfection of his obedience.\footnote{Gal. iii. 24; Rom. vii. 24, 25; viii. 3, 4.} It is likewise of use to the regenerate, to restrain their corruptions, in that it forbids sin;\footnote{James ii. 11; Psa. cxix. 101, 104, 128.} and the threatenings of it serve to show what even their sins deserve, and what afflictions in this life they may expect for them, although freed from the curse thereof threatened in the law.\footnote{Ezra ix. 13, 14; Psa. lxxxix. 30–34.} The promises of it, in like manner, show them God's approbation of obedience, and what blessings they may expect upon the performance thereof;\footnote{Lev. xxvi. 1, 10, 14, with 2 Cor. vi. 16, Eph. vi. 2, 3; Psa. xxxvii. 11 with Matt. v. 5; Psa. xix. 11.} although not as due to them by the law as a covenant of words:\footnote{Gal. ii. 16; Luke xvii. 10.} so as a man's doing good, and refraining from evil, because the law encourageth to the one, and deterreth from the other, is no evidence of his being under the law, and not under grace.\footnote{Rom. vi. 12, 14; 1 Pet. iii. 8–12 with Psa.xxxiv. 12–16; Heb. xii. 28, 29.}

\emph{edocendo dirigit simul et obligat ad consentanee ambulandum;}\footnote{Rom. vii. 12, 22, 25; Psa. cxix. 4–6; 1 Cor. vii. 19; Gal. v. 14, 16, 18–23.}\emph{ ipsisque patere facit naturæ, cordis, vitæque suæ nefaria inquinamenta:}\footnote{Rom. vii. 7; iii. 20.}\emph{ adeo ut ad illam semet exigentes, cum peccati ulterius convinci, pro eodem humiliari, ac ejusdem odio inflammari possint;}\footnote{James i. 23–25; Rom. vii. 9, 14, 24.}\emph{ tum vero etiam ut perspicere possint evidentius quam plane necessarius eis Christus, quamque perfecta sit ejusdem obedientia.}\footnote{Gal. iii. 24; Rom. vii. 24, 25; viii. 3, 4.}\emph{ Verum ulterius etiam regenitis ea utilis esse possit, in quantum nempe corruptiones eorum peccata prohibendo coërcet,}\footnote{James ii. 11; Psa. cxix. 101, 104, 128.}\emph{ graviter autem interminando indicat tum quid vel eorum peccata commeruerint, tum etiam quas ea propter in hac vita afflictiones expectare possint, utcunque ab earum maledictione, quam lex minatur, liberentur.}\footnote{Ezra ix. 13, 14; Psa. lxxxix. 30–34.}\emph{ Quinetiam promissiones ejus demonstrant iis obedientia Deo quam accepta sit et approbata; quasque illa præstita benedictiones}\footnote{Lev. xxvi. 1, 10, 14, with 2 Cor. vi. 16, Eph. vi. 2, 3; Psa. xxxvii. 11 with Matt. v. 5; Psa. xix. 11.} (\emph{licet non tanquam lege debitas ex operum fœdere})\footnote{Gal. ii. 16; Luke xvii. 10.}\emph{ possint illi expectare. Adeo ut quod quis bonum præstet invitante lege, a malo autem abhorreat lege deterritus, nullo prorsus argumento sit, eum sub lege esse, non vero sub gratia constitutum.}\footnote{Rom. vi. 12, 14; 1 Pet. iii. 8–12 with Psa.xxxiv. 12–16; Heb. xii. 28, 29.}

VII. Neither are the forementioned uses of the law contrary to the grace of the gospel, but do sweetly comply with it:\footnote{Gal. iii. 21; [Am. ed. Titus ii. 11–14].} the Spirit of Christ subduing and enabling the will of man to do that freely and cheerfully which the will of God, revealed in the law, requireth to be done.\footnote{Ezek. xxxvi. 27; Heb. viii. 10 with Jer. xxxi. 33.}

VII. Neque interim Legis usus isti iam memorati, Evangelii gratiæ adversantur, sed cum eadem conspirant suaviter,\footnote{Gal. iii. 21; [Am. ed. Titus ii. 11–14].} voluntatem humanam ita subjugante ac imbuente Christi Spiritu, ut idem illud præstare valeat spontanee ac alacriter, quod ab illa exigit voluntas Dei in lege sua revelata.\footnote{Ezek. xxxvi. 27; Heb. viii. 10 with Jer. xxxi. 33.}

\xsection{Chapter XX. Of Christian Liberty, and Liberty of Conscience}

I. The liberty which Christ hath purchased for believers under the gospel consists in their freedom from the guilt of sin, the condemning wrath of God, the curse of the moral law;\footnote{Titus ii. 14; 1 Thess. i. 10; Gal. iii. 13.} and in their being delivered from this present evil world, bondage to Satan, and dominion of sin,\footnote{Gal. i. 4; Col. i. 13; Acts xxvi. 18; Rom. vi. 14.} from the evil of afflictions, the sting of death, the victory of the grave, and everlasting damnation;\footnote{Rom. viii. 28; Psa. cxix. 71; 1 Cor. xv. 54–57; Rom. viii. 1.} as also in their free access to God,\footnote{Rom. v. 1, 2.} and their yielding obedience unto him, not out of slavish fear, but a childlike love and\footnote{[Am. ed. inserts \emph{a} after \emph{and}.]} willing mind.\footnote{Rom. viii. 14, 15; 1 John iv. 18.} All which were common also to believers under the law;\footnote{Gal. iii. 9, 14.} but under the New Testament the liberty of Christians is further enlarged in

I. \emph{Libertas quam Christus acquisivit fidelibus sub Evangelio in eo sita est, quod a reatu peccati, ab ira Dei condemnante, a legis Moralis maledictione immunes fiant,}\footnote{Titus ii. 14; 1 Thess. i. 10; Gal. iii. 13.}\emph{ quod a præsenti malo seculo, a dura Satanæ servitute, dominioque peccati:}\footnote{Gal. i. 4; Col. i. 13; Acts xxvi. 18; Rom. vi. 14.}\emph{ ab afflictionum malo, ab aculeo mortis, a sepulchri victoria ab æterna denique damnation}\footnote{Rom. viii. 28; Psa. cxix. 71; 1 Cor. xv. 54–57; Rom. viii. 1.}\emph{ liberentur; Quodque libere eis liceat ad Deum accedere:}\footnote{Rom. v. 1, 2.}\emph{ eique non e metu servile, verum e filiali dilectione, promtoque animo præbere valeant obedientiam.}\footnote{Rom. viii. 14, 15; 1 John iv. 18.}\emph{ Atque hæc quidem omnia cum fidelibus sub lege habent communia.}\footnote{Gal. iii. 9, 14.}\emph{ Verum sub Novo Testamento ulterius adhuc se extendit libertas Christiana; in quantum }

their freedom from the yoke of the ceremonial law, to which the Jewish Church was subjected;\footnote{Gal. iv. 1–3, 6, 7; v. 1; Acts xv. 10, 11.} and in greater boldness of access to the throne of grace,\footnote{Heb. iv. 14, 16; x. 19–22.} and in fuller communications of the free Spirit of God, than believers under the law did ordinarily partake of.\footnote{John vii. 38, 39; 2 Cor. iii. 13, 17, 18.}

\emph{nempe Legis ceremonialis jugo, cui subjecta erat Ecclesia Judaica, eximuntur;}\footnote{Gal. iv. 1–3, 6, 7; v. 1; Acts xv. 10, 11.}\emph{ majoremque confidentiam ad thronum gratiæ accedendi,}\footnote{Heb. iv. 14, 16; x. 19–22.}\emph{ sed et effusiorem gratuiti Spiritus Dei communicationem sunt consecuti, quam ordinarie sub Lege fideles participarunt.}\footnote{John vii. 38, 39; 2 Cor. iii. 13, 17, 18.}

II. God alone is Lord of the conscience,\footnote{James iv. 12; Rom. xiv. 4.} and hath left it free from the doctrines and commandments of men which are in any thing contrary to his Word, or beside it in matters of faith or worship.\footnote{Acts iv. 19; v. 29; 1 Cor. vii. 23; Matt. xxiii. 8–10; 2 Cor. i. 24; Matt. xv. 9.} So that to believe such doctrines, or to obey such commands\footnote{[Am. ed. \emph{commandments}.]} out of conscience, is to betray true liberty of conscience;\footnote{Col. ii. 20–23; Gal. i. 10; v. 1; ii. 4, 5; Psa. v. l.} and the requiring of\footnote{[Am. ed. omits \emph{of}.]} an implicit faith, and an absolute and blind obedience, is to destroy liberty of conscience, and reason also.\footnote{Rom. x. 17; xiv. 23; Isa. viii. 20; Acts xvii. 11; John iv. 22; Hos. v. 11; Rev. xiii. 12, 16, 17; Jer. viii. 9.}

II. \emph{Deus solus Dominus est conscientiæ,}\footnote{James iv. 12; Rom. xiv. 4.}\emph{ quam certe exemit doctrinis et mandatis hominum, ubi aut verbo ejus adversantur, aut in rebus fidei et cultus quicquam ei superaddunt.}\footnote{Acts iv. 19; v. 29; 1 Cor. vii. 23; Matt. xxiii. 8–10; 2 Cor. i. 24; Matt. xv. 9.}\emph{ Unde qui ejusmodi aut doctrinas credunt, aut mandatis obtemperant, quasi ad id ex conscientia teneantur, veram ii conscientiæ libertatem produnt.}\footnote{Col. ii. 20–23; Gal. i. 10; v. 1; ii. 4, 5; Psa. v. l.}\emph{ Qui autem vel fidem implicitam, vel obedientiam absolutam cæcamque exigunt, næ illi id agunt, ut cum conscientiæ, tum rationis etiam destruant libertatem.}\footnote{Rom. x. 17; xiv. 23; Isa. viii. 20; Acts xvii. 11; John iv. 22; Hos. v. 11; Rev. xiii. 12, 16, 17; Jer. viii. 9.}

III. They who, upon pretense of Christian liberty, do practice any sin, or cherish any lust, do thereby destroy the end of Christian liberty; which is, that, being delivered out of the hands of our enemies, we might serve the Lord without fear, in holiness and righteousness before him, all the days of our life.\footnote{Gal. v. 13; 1 Pet. ii. 16; 2 Pet. ii. 19; John viii. 34; Luke i. 74, 75.}

III. \emph{Qui sub prætextu Christianæ libertatis, cuivis aut cupiditati indulgent aut peccato assuescunt, eo ipso libertatis Christianæ finem corrumpunt; nempe ut e manibus inimicorum nostrorum liberati, Domino in sanctimonia et justitia coram ipso omnibus diebus vitæ nostræ absque metu serviamus.}\footnote{Gal. v. 13; 1 Pet. ii. 16; 2 Pet. ii. 19; John viii. 34; Luke i. 74, 75.}

IV. And because the power\footnote{[Am. ed. \emph{powers}.]} which God hath ordained, and the liberty which Christ hath purchased, are not intended by God to destroy, but mutually to uphold and preserve one another; they who, upon pretense of Christian liberty, shall oppose any lawful power, or the lawful exercise of it, whether it be civil or ecclesiastical, resist the ordinance of God.\footnote{Matt. xii. 25; 1 Pet. ii. 13, 14, 16; Rom. xiii. 1–8; Heb. xiii. 17.} And for their publishing of such opinions, or maintaining of such practices, as are contrary to the light of nature, or to the known principles of Christianity, whether concerning faith, worship, or conversation; or to the power of godliness; or such erroneous opinions or practices, as, either in their own nature, or in the manner of publishing or maintaining them, are destructive to the external peace and order which Christ hath established in the Church; they may lawfully be called to account, and proceeded against by the censures of the Church,\footnote{Rom. i. 32 with 1 Cor. v. 1, 5, 11, 13; 2 John v. 10, 11; and 2 Thess. iii. 14, and 1 Tim. vi. 3–5, and Titus i. 10, 11, 13 and iii. 10, with Matt. xviii. 15–17; 1 Tim. i. 19, 20; Rev. ii. 2, 14, 15, 20; iii. 9.} and by the power of the Civil Magistrate.\footnote{Deut. xiii. 6–12; Rom. xiii. 3, 4, with 2 John v. 10, 11; Ezra vii. 23–28; Rev. xvii. 12, 16, 17; Neh. xiii. 15, 17, 21, 22, 25, 30; 2 Kings xxiii. 5, 6, 9, 20, 21; 2 Chron. xxxiv. 33; xv. 12, 13, 16; Dan. iii. 29; 1 Tim. ii. 2; Isa. xlix. 23; Zech. xiii. 2, 3.}\footnote{[Am. ed. omits \emph{and by the power of the Civil Magistrate}, also the prooftexts.]}

IV. \emph{Quoniam vero potestates quas Deus ordinavit, et libertas quam acquisivit Christus non in eum finem a Deo destinatæ sunt ut se mutuo perimant, verum ut se sustentent ac conservent invicem; Qui itaque sub libertatis Christianæ prætextu potestati cuivis legitimæ }(\emph{civilis sit sive Ecclesiastica}) \emph{aut legitimo ejusdem exercitio contraiverint, ordinationi divinæ resistere censendi sunt,}\footnote{Matt. xii. 25; 1 Pet. ii. 13, 14, 16; Rom. xiii. 1–8; Heb. xiii. 17.}\emph{ Quique vel ejusmodi opiniones publicaverint, praxesve defenderint, quæ lumini naturæ, aut religionis Christianæ de fide, de cultu, aut moribus principiis notis, aut pietatis denique vi ac efficaciæ adversantur; vel ejusmodi opiniones praxesve erroneas, quæ aut sua natura aut publicationis defensionisve modo, externæ paci ac eutaxiæ, quas in Ecclesia sua stabilivit Christus, perniciem minitantur; omnino licitum est tum ab iis facti rationem reposcere, tum in eos qua censuris Ecclesiasticis,}\footnote{Rom. i. 32 with 1 Cor. v. 1, 5, 11, 13; 2 John v. 10, 11; and 2 Thess. iii. 14, and 1 Tim. vi. 3–5, and Titus i. 10, 11, 13 and iii. 10, with Matt. xviii. 15–17; 1 Tim. i. 19, 20; Rev. ii. 2, 14, 15, 20; iii. 9.}\emph{ qua civilis magistratus potestate animadvertere.}\footnote{Deut. xiii. 6–12; Rom. xiii. 3, 4, with 2 John v. 10, 11; Ezra vii. 23–28; Rev. xvii. 12, 16, 17; Neh. xiii. 15, 17, 21, 22, 25, 30; 2 Kings xxiii. 5, 6, 9, 20, 21; 2 Chron. xxxiv. 33; xv. 12, 13, 16; Dan. iii. 29; 1 Tim. ii. 2; Isa. xlix. 23; Zech. xiii. 2, 3.}

\xsection{Chapter XXI. Of Religious Worship and the Sabbath-day}

I. The light of nature showeth that there is a God, who hath lordship and sovereignty over all; is good, and doeth good unto all; and is therefore to be feared, loved, praised, called upon, trusted in, and served with all the heart, and with all the soul, and with all the might.\footnote{Rom. i. 20; Acts xvii. 24; Psa. cxix. 68; Jer. x. 7; Psa. xxxi. 23; xviii. 3; Rom. x. 12; Psa. lxii. 8; Josh. xxiv. 14; Mark xii. 33.} But the acceptable way of worshiping the true God is instituted by himself, and so limited to\footnote{[Am. ed. \emph{by}.]} his own revealed will, that he may not be worshiped according to the imaginations and devices of men, or the suggestions of Satan, under any visible representations\footnote{[Am. ed. \emph{representation}.]} or any other way not prescribed in the Holy Scripture.\footnote{Deut. xii. 32; Matt. xv. 9; Acts xvii. 25; Matt. iv. 9, 10; Deut. iv. 15–20; Exod. xx. 4–6; Col. ii. 23.}

I. \emph{Constat quidem naturæ lumine esse Deum qui in universa Primatum obtinet ac absolutum Dominium, eundemque bonum esse ac omnibus beneficum, proindeque toto corde, tota anima, totisque viribus timendum esse et diligendum, laudandum ac invocandum, eique fidendum esse ac serviendum.}\footnote{Rom. i. 20; Acts xvii. 24; Psa. cxix. 68; Jer. x. 7; Psa. xxxi. 23; xviii. 3; Rom. x. 12; Psa. lxii. 8; Josh. xxiv. 14; Mark xii. 33.}\emph{ At rationem verum Deum colendi acceptabilem ipse instituit, itaque voluntate sua revelata definivit, ut coli non debeat secundum imaginationes ac inventa hominum, aut suggestiones Satanæ, sub specie quavis visibili, aut alia via quaviscunque quam scriptura sacra non præscripsit.}\footnote{Deut. xii. 32; Matt. xv. 9; Acts xvii. 25; Matt. iv. 9, 10; Deut. iv. 15–20; Exod. xx. 4–6; Col. ii. 23.}

II. Religious worship is to be given to God, the Father, Son, and Holy Ghost; and to him alone:\footnote{Matt. iv. 10 with John v. 23 and 2 Cor. xiii. 14; [Am. ed. Rev. v. 11–13].} not to angels, saints, or any other creature:\footnote{Col. ii. 18; Rev. xix. 10; Rom. i. 25.} and since the fall, not without a Mediator; nor in the mediation of any other but of Christ alone.\footnote{John xiv. 6; 1 Tim. ii. 5; Eph. ii. 18; Col. iii. 17.}

II. \emph{Cultus religiosus Deo Patri Filio et Spiritui sancto, eique soli est exhibendus,}\footnote{Matt. iv. 10 with John v. 23 and 2 Cor. xiii. 14; [Am. ed. Rev. v. 11–13].}\emph{ non angelis, non sanctis, neque alii cuivis creaturæ,}\footnote{Col. ii. 18; Rev. xix. 10; Rom. i. 25.}\emph{ nec ipsi Deo quidem post lapsum citra Mediatorem, aut quidem per Mediatorem alium quam Jesum Christum.}\footnote{John xiv. 6; 1 Tim. ii. 5; Eph. ii. 18; Col. iii. 17.}

III. Prayer with thanksgiving, being one special part of religious worship,\footnote{Phil. iv. 6.} is by God required of all

III. \emph{Supplicationem cum gratiarum actione, quæ est inter partes præcipuas divini cultus,}\footnote{Phil. iv. 6.}\emph{ Deus fieri }

men;\footnote{Psa. lxv. 2.} and that it may be accepted, it is to be made in the name of the Son,\footnote{John xiv. 13, 14; 1 Pet. ii. 5.} by the help of his Spirit,\footnote{Rom. viii. 26.} according to his will,\footnote{1 John v. 14.} with understanding, reverence, humility, fervency, faith, love, and perseverance;\footnote{Psa. xlvii. 7; Eccles. v. 1, 2; Heb. xii. 28; Gen. xviii. 27; James v. 16; i. 6, 7; Mark xi. 24; Matt. vi. 12, 14, 15; Col. iv. 2; Eph. vi. 18.} and, if vocal, in a known tongue.\footnote{1 Cor. xiv. 14.}

\emph{iubet ab hominibus universis;}\footnote{Psa. lxv. 2.}\emph{ quæ, quo Deo grata sit et accepta, est in nomine Filii,}\footnote{John xiv. 13, 14; 1 Pet. ii. 5.}\emph{ subsidio spiritus ejus,}\footnote{Rom. viii. 26.}\emph{ et secundum ipsius voluntatem,}\footnote{1 John v. 14.}\emph{ cum intellectu, reverentia, humilitate, fervore, fide, amore, ac perseverantia offerenda;}\footnote{Psa. xlvii. 7; Eccles. v. 1, 2; Heb. xii. 28; Gen. xviii. 27; James v. 16; i. 6, 7; Mark xi. 24; Matt. vi. 12, 14, 15; Col. iv. 2; Eph. vi. 18.}\emph{ et quidem, si vocalis sit, in lingua nota est efferenda.}\footnote{1 Cor. xiv. 14.}

IV. Prayer is to be made for things lawful,\footnote{1 John v. 14.} and for all sorts of men living, or that shall live hereafter;\footnote{1 Tim. ii. 1, 2; John xvii. 20; 2 Sam. vii. 29; Ruth iv. 12.} but not for the dead,\footnote{2 Sam. xii. 21–23 with Luke xvi. 25, 26; Rev. xiv. 13.} nor for those of whom it may be known that they have sinned the sin unto death.\footnote{1 John v. 16.}

IV. \emph{Preces pro rebus non nisi licitis sunt faciendæ,}\footnote{1 John v. 14.}\emph{ pro hominibus autem cuiuscunque generis, vivis scilicet, aut etiam victuris aliquando;}\footnote{1 Tim. ii. 1, 2; John xvii. 20; 2 Sam. vii. 29; Ruth iv. 12.}\emph{ pro mortuis autem neutiquam;}\footnote{2 Sam. xii. 21–23 with Luke xvi. 25, 26; Rev. xiv. 13.}\emph{ sed neque pro iis, de quibus constare possit eos peccatum ad mortem perpetrasse.}\footnote{1 John v. 16.}

V. The reading of the Scriptures with godly fear;\footnote{Acts xv. 21; Rev. i. 3.} the sound preaching;\footnote{2 Tim. iv. 2.} and conscionable hearing of the Word, in obedience unto God with understanding, faith, and reverence;\footnote{James i. 22; Acts x. 33; Matt. xiii. 19; Heb. iv. 2; Isa. lxvi. 2.} singing of psalms with grace in the heart;\footnote{Col. iii. 16; Eph. v. 19; James v. 13.} as, also, the dne administration and worthy receiving of the sacraments instituted by Christ; are all parts of the ordinary religious worship of God:\footnote{Matt. xxviii. 19; 1 Cor. xi. 23–29; Acts ii. 42.} besides religious oaths,\footnote{Deut. vi. 13 with Neh. x. 29.} vows,\footnote{Isa. xix. 21 with Eccles. v. 4, 5; [Am. ed. Acts xviii. 18.—Am. ed. reads \emph{and vows}].} solemn

V. \emph{Scripturarum lectio cum timore pio;}\footnote{Acts xv. 21; Rev. i. 3.}\emph{ verbi prædicatio solida,}\footnote{2 Tim. iv. 2.}\emph{ ejusdemque auditio religiosa ex obedientia erga Deum, cum intellectu, fide et reverentia;}\footnote{James i. 22; Acts x. 33; Matt. xiii. 19; Heb. iv. 2; Isa. lxvi. 2.}\emph{ Psalmorum cum gratia in corde cantatio,}\footnote{Col. iii. 16; Eph. v. 19; James v. 13.}\emph{ prout etiam Sacramentorum, quæ Christus instituit, debita administratio, et participatio digna, sunt divini cultus reiigiosi partes, et quidem ordinarii}.\footnote{Matt. xxviii. 19; 1 Cor. xi. 23–29; Acts ii. 42.}\emph{ Religiosa insuper juramenta,}\footnote{Deut. vi. 13 with Neh. x. 29.}\emph{ votaque;}\footnote{Isa. xix. 21 with Eccles. v. 4, 5; [Am. ed. Acts xviii. 18.—Am. ed. reads \emph{and vows}].}\emph{ solennia jejunia,}\footnote{Joel ii. 12; Esth. iv. 16; Matt. ix. 15; 1 Cor. vii. 5.}

fastings,\footnote{Joel ii. 12; Esth. iv. 16; Matt. ix. 15; 1 Cor. vii. 5.} and thanksgivings upon several\footnote{[Amer. ed. has \emph{special}]} occasions;\footnote{Psalm cvii. throughout; Esth. ix. 22.} which are, in their several times and seasons, to be used in an holy and religious manner.\footnote{Heb. xii. 28.}

\emph{solennesque gratiarum actiones, pro varietate eventuum}\footnote{Psalm cvii. throughout; Esth. ix. 22.}\emph{; suo quæque tempore ac opportunitate sancte quidem ac religiose sunt adhibenda.}\footnote{Heb. xii. 28.}

VI. Neither prayer, nor any other part of religious worship, is now, under the gospel, either tied unto, or made more acceptable by any place in which it is performed, or towards which it is directed:\footnote{John iv. 21.} but God is to be worshiped every where\footnote{Mal. i. 11; 1 Tim. ii. 8.} in spirit and\footnote{[Am. ed. inserts \emph{in}.]} truth;\footnote{John iv. 23, 24.} as in private families\footnote{Jer. x. 25; Deut. vi. 6, 7; Job i. 5; 2 Sam. vi. 18, 20; 1 Pet. iii. 7; Acts x. 2.} daily,\footnote{Matt. vi. 11; [Am. ed. Josh. xxiv. 15].} and in secret each one by himself,\footnote{Matt. vi. 6; Eph. vi. 18.} so more solemnly in the public assemblies, which are not carelessly or willfully to be neglected or forsaken, when God, by his Word or providence, calleth thereunto.\footnote{Isa. lvi. 7; Heb. x. 25; Prov. i. 20, 21, 24; viii. 34; Acts xiii. 42; Luke iv. 16; Acts ii. 42.}

VI.\emph{ Hodie sub evangelio neque preces, nec ulla pars alia religiosi cultus ita cuivis alligatur loco in quo præstetur aut versus quem dirigatur,}\footnote{John iv. 21.}\emph{ ut inde gratior evadat et acceptior; verum ubique Deus colendus est}\footnote{Mal. i. 11; 1 Tim. ii. 8.}\emph{ in spiritu ac veritate;}\footnote{John iv. 23, 24.}\emph{ quotidie}\footnote{Jer. x. 25; Deut. vi. 6, 7; Job i. 5; 2 Sam. vi. 18, 20; 1 Pet. iii. 7; Acts x. 2.}\emph{ quidem inter privatos parietes a quavis familia,}\footnote{Matt. vi. 11; [Am. ed. Josh. xxiv. 15].}\emph{ ut etiam a quolibet seorsim in secreto;}\footnote{Matt. vi. 6; Eph. vi. 18.}\emph{ at solenniter magis in conventibus publicis, qui certe quoties eo nos Deus vocat, seu verbo suo seu providentia, non sunt vel ex incuria vel obstinatione animi aut negligendi aut deserendi.}\footnote{Isa. lvi. 7; Heb. x. 25; Prov. i. 20, 21, 24; viii. 34; Acts xiii. 42; Luke iv. 16; Acts ii. 42.}

VII. As it is of the law of nature, that, in general, a due proportion of time be set apart for the worship of God; so, in his Word, by a positive, moral, and perpetual commandment, binding all men in all ages, he hath particularly appointed one day in seven for a Sabbath, to be kept holy unto him:\footnote{Exod. xx. 8, 10, 11; Isa. lvi. 2, 4, 6, 7; [Am. ed. Isa. lvi. 6].} which, from the

VII. \emph{Quemadmodum est de lege naturæ ut indefinite portio quædam temporis idonea divino cultui celebrando sejuncta sit ac assignata; ita in verbo suo Deus }(\emph{præcepto morali, positivo ac perpetuo, homines omnes cujuscunque fuerint seculi obligante})\emph{ speciatim e septenis quibusque diebus diem unum in Sabbatum designavit, sancte sibi observandum.}\footnote{Exod. xx. 8, 10, 11; Isa. lvi. 2, 4, 6, 7; [Am. ed. Isa. lvi. 6].}\emph{ Quod }

beginning of the world to the resurrection of Christ, was the last day of the week; and, from the resurrection of Christ, was changed into the first day of the week,\footnote{Gen. ii. 2, 3; 1 Cor. xvi. 1, 2; Acts xx. 7.} which in Scripture is called the Lord's day,\footnote{Rev. i. 10.} and is to be continued to the end of the world, as the Christian Sabbath.\footnote{Exod. xx. 8, 10, with Matt. v. 17, 18.}

\emph{quidem ab orbe condito ad resurrectionem usque Christi dies ultimus erat in septimana; deinde autem a Christi resurrectione in septimanæ diem primum transferebatur;}\footnote{Gen. ii. 2, 3; 1 Cor. xvi. 1, 2; Acts xx. 7.}\emph{ qui quidem in Scriptura }Dies Dominicus\footnote{Rev. i. 10.}\emph{ nuncupatur, estque perpetuo ad finem mundi tanquam Sabbatum Christianum celebrandus.}\footnote{Exod. xx. 8, 10, with Matt. v. 17, 18.}

VIII. This Sabbath is then kept holy unto the Lord, when men, after a due preparing of their hearts, and ordering of their common affairs beforehand, do not only observe an holy rest all the day from their own works, words, and thoughts, about their worldly employments and recreations;\footnote{Exod. xx. 8; xvi. 23, 25, 26, 29, 30; xxxi. 15–17; Isa. lviii. 13; Neh. xiii. 15–22.} but also are taken up the whole time in the public and private exercises of his worship, and in the duties of necessity and mercy.\footnote{Isa. lviii. 13; Matt. xii. 1–13.}

VIII. \emph{Tunc autem hoc Sabbatum Deo sancte celebratur, quum post corda rite præparata, et compositas suas res mundanas, homines non solum a suis ipsorum operibus, dictis, cogitatis; }(\emph{quæ circa illas exerceri solent})\emph{ a recreationibus etiam ludicris quietem sanctam toto observant die;}\footnote{Exod. xx. 8; xvi. 23, 25, 26, 29, 30; xxxi. 15–17; Isa. lviii. 13; Neh. xiii. 15–22.}\emph{ verum etiam in exercitiis divini cultus publicis privatisque, ac in officiis necessitatis et misericordiæ toto illo tempore occupantur.}\footnote{Isa. lviii. 13; Matt. xii. 1–13.}

\xsection{Chapter XXII. Of Lawful Oaths and Vows}

I. A lawful oath is a part of religious worship,\footnote{Deut. x. 20.} wherein, upon just occasion, the person swearing solemnly calleth God to witness what he asserteth or promiseth; and to judge him according to the truth or falsehood of what he sweareth.\footnote{Exod. xx. 7; Lev. xix. 12; 2 Cor. i. 23; 2 Chron. vi. 22, 23.}

I. \emph{Juramentum licitum est pars cultus religiosi,}\footnote{Deut. x. 20.}\emph{ qua }(\emph{occasione justa oblata})\emph{ qui jurat, Deum, de eo quod asserit aut promittit, solenni modo testatur; eundemque appellat se secundum illius quod jurat veritatem aut falsitatem judicaturum.}\footnote{Exod. xx. 7; Lev. xix. 12; 2 Cor. i. 23; 2 Chron. vi. 22, 23.}

II. The name of God only is that

II. \emph{Per solum Dei nomen jurare }

by which, men ought to swear, and therein it is to be used with all holy fear and reverence;\footnote{Deut. vi. 13.} therefore to swear vainly or rashly by that glorious and dreadful name, or to swear at all by any other thing, is sinful, and to be abhorred.\footnote{Exod. xx. 7; Jer. v. 7; Matt. v. 34, 37; James v. 12.} Yet as, in matters of weight and moment, an oath is warranted by the Word of God, under the New Testament, as well as under the Old,\footnote{Heb. vi. 16; 2 Cor. i. 23; Isa. lxv. 16.} so a lawful oath, being imposed by lawful authority, in such matters ought to be taken.\footnote{1 Kings viii. 31; Neh. xiii. 25; Ezra x. 25.}

\emph{debent homines, quod quidem cum omni timore sancto ac reverentia est inibi usurpandum.}\footnote{Deut. vi. 13.}\emph{ Proindeque per nomen illud gloriosum ac tremendum jurare leviter, aut temere, vel etiam omnino jurare per rem aliam quamviscunque, sceleratum est et quam maxime perhorrescendum.}\footnote{Exod. xx. 7; Jer. v. 7; Matt. v. 34, 37; James v. 12.}\emph{ Veruntamen sicut in rebus majoris ponderis et momenti secundum verbum Dei licitum est jusjurandum non minus quidem sub Novo quam sub Vetere Testamento:}\footnote{Heb. vi. 16; 2 Cor. i. 23; Isa. lxv. 16.}\emph{ ita sane jusjurandum licitum, authoritate legitima si exigatur, non est in rebus ejusmodi declinandum.}\footnote{1 Kings viii. 31; Neh. xiii. 25; Ezra x. 25.}

III. Whosoever taketh an oath ought duly to consider the weightiness of so solemn an act, and therein to avouch nothing but what he is fully persuaded is the truth.\footnote{Exod. xx. 7; Jer. iv. 2.} Neither may any man bind himself by oath to any thing but what is good and just, and what he believeth so to be, and what he is able and resolved to perform.\footnote{Gen. xxiv. 2, 3, 5, 6, 8, 9.} Yet it is a sin to refuse an oath touching any thing that is good and just, being imposed by lawful authority.\footnote{Numb. v. 19, 21; Neh. v. 12; Exod. xxii. 7–11.}

III. \emph{Quicunque juramentum præstat eum pondus actionis tam solennis rite secum perpendere oportet, atque juratum de nullo asseverare quod verum esse non habeat sibi persuasissimum.}\footnote{Exod. xx. 7; Jer. iv. 2.}\emph{ Neque licet cuivis ad agendum quicquam obstringere semet jurejurando, nisi quod revera bonum justumque est, quod ille ejusmodi esse credit, quodque ipse præstare potest statuitque.}\footnote{Gen. xxiv. 2, 3, 5, 6, 8, 9.}\emph{ Veruntamen de re bona justaque jusjurandum, legitima authoritate si exigatur, peccat ille qui detrectat.}\footnote{Numb. v. 19, 21; Neh. v. 12; Exod. xxii. 7–11.}

IV. An oath is to be taken in the plain and common sense of

IV. \emph{Juramentum præstandum est sensu verborum vulgari quidem ac }

the words, without equivocation or mental reservation.\footnote{Jer. iv. 2; Psa. xxiv. 4.} It can not oblige to sin; but in any thing not sinful, being taken, it binds to performance, although to a man's own hurt:\footnote{1 Sam. xxv. 22, 32–34; Psa. xv. 4.} nor is it to be violated, although made to heretics or infidels.\footnote{Ezek. xvii. 16, 18, 19; Josh. ix. 18, 19, with 2 Sam. xxi. 1.}

\emph{manifesto, sine æquivocatione aut reservatione mentali quaviscunque.}\footnote{Jer. iv. 2; Psa. xxiv. 4.}\emph{ Ad peccandum quenquam obligare nequit, verum in re qualibet cui abest peccatum, qui semel illud præstitit, adimplere tenetur, vel etiam cum damno suo;}\footnote{1 Sam. xxv. 22, 32–34; Psa. xv. 4.}\emph{ neque sane licet, quamvis hæreticis datum aut infidelibus, violare.}\footnote{Ezek. xvii. 16, 18, 19; Josh. ix. 18, 19, with 2 Sam. xxi. 1.}

V. A vow is of the like nature with a promissory oath, and ought to be made with the like religious care, and to be performed with the like faithfulness.\footnote{Isa. xix. 21; Eccles. v. 4–6; Psa. lxi. 8; lxvi. 13, 14.}

V. \emph{Votum, naturæ consimilis est cum juramento promissorio, parique debet tum religione nuncupari tum fide persolvi.}\footnote{Isa. xix. 21; Eccles. v. 4–6; Psa. lxi. 8; lxvi. 13, 14.}

VI. It is not to be made to any creature, but to God alone:\footnote{Psa. lxxvi. 11; Jer. xliv. 25, 26.} and that it may be accepted, it is to be made voluntarily, out of faith and conscience of duty, in way of thankfulness for mercy received, or for the\footnote{[Am. ed. omits \emph{the.}]} obtaining of what we want; whereby we more strictly bind ourselves to necessary duties, or to other things, so far and so long as they may fitly conduce thereunto.\footnote{Deut. xxiii. 21, 23; Psa. 1. 14; Gen. xxviii. 20–22; 1 Sam. i. 11; Psa. lxvi. 13, 14; cxxxii. 2—5.}

VI. \emph{Non est ulli creaturæ, sed Deo soli nuncupandum,}\footnote{Psa. lxxvi. 11; Jer. xliv. 25, 26.}\emph{ et quo gratum illi esse possit acceptumque, est quidem lubenter, e fide, officiique nostri conscientia suscipiendum, vel gratitudinis nostræ ob accepta beneficia testandæ causa, vel boni alicujus, quo indigemus, consequendi; per hoc autem nosmet ad officia necessaria arctius obligamus; vel etiam ad res alias quatenus quidem et quamdiu istis subserviunt.}\footnote{Deut. xxiii. 21, 23; Psa. 1. 14; Gen. xxviii. 20–22; 1 Sam. i. 11; Psa. lxvi. 13, 14; cxxxii. 2—5.}

VII. No man may vow to do any thing forbidden in the Word of God, or what would hinder any duty therein commanded, or which is not in his own power, and for the performance whereof he hath no promise

VII. \emph{Nemini quicquam vovere licet se acturum, quod aut verbo Dei prohibetur; aut officium aliquod inibi præceptum impediret, quodve non est in voventis potestate, et cui præstando vires illi Deus non est pollicitus.}\footnote{Acts xxiii. 12, 14; Mark vi. 26; Numb. xxx. 5, 8, 12, 13.}

or ability from God.\footnote{Acts xxiii. 12, 14; Mark vi. 26; Numb. xxx. 5, 8, 12, 13.} In which respect,\footnote{[Am. ed. has \emph{respects.}]} popish monastical vows of perpetual single life, professed poverty, and regular obedience, are so far from being degrees of higher perfection, that they are superstitions and sinful snares, in which no Christian may entangle himself.\footnote{Matt. xix. 11, 12; 1 Cor. vii. 2, 9; Eph. iv. 28; 1 Pet. iv. 2; 1 Cor. vii. 23.}

\emph{Unde Pontificiorum illa de perpetuo cœlibatu, de paupertate, deque obedientia regulari vota Monastica, tantum abest ut perfectionis gradus sint sublimiores, ut superstitionis plane sint ac peccati laquei, quibus nulli unquam Christiano semetipsum licet implicare.}\footnote{Matt. xix. 11, 12; 1 Cor. vii. 2, 9; Eph. iv. 28; 1 Pet. iv. 2; 1 Cor. vii. 23.}

\xsection{Chapter XXIII. Of the Civil Magistrate}

I. God, the Supreme Lord and King of all the world, hath ordained civil magistrates to be under him, over the people, for his own glory and the public good, and to this end hath armed them with the power of the sword, for the defense and encouragement of them that are good, and for the punishment of evil-doers.\footnote{Rom. xiii. 1–4; 1 Pet. ii. 13, 14.}

I. \emph{Supremus totius Mundi Rex ac Dominus Deus, Magistratus Civiles ordinavit qui vices ejus gerant supra populum ad suam ipsius gloriam, ac bonum publicum; in quem finem eosdem armavit potestate gladii, propter bonorum quidem animationem ac tutamen, animadversionem autem in maleficos.}\footnote{Rom. xiii. 1–4; 1 Pet. ii. 13, 14.}

II. It is lawful for Christians to accept and execute the office of a magistrate when called thereunto;\footnote{Prov. viii. 15, 16; Rom. xiii. 1, 2, 4.} in the managing whereof, as they ought especially to maintain piety, justice, and peace, according to the wholesome laws of each commonwealth,\footnote{Psa. ii. 10–12; 1 Tim. ii. 2; Psa. lxxxii. 3, 4; 2 Sam. xxiii. 3; 1 Pet. ii. 13.} so, for that end, they may lawfully, now under the New Testament, wage war upon just and necessary occasion.\footnote{Luke iii. 14; Rom. xiii. 4; Matt. viii. 9, 10; Acts x. 1, 2; Rev. xvii. 14, 16.}\textsuperscript{\&}\footnote{[Am. ed. has \emph{occasions}.]}

II. \emph{Christianis, quoties ad id vocantur, Magistratus munus et suscipere licet et exequi;}\footnote{Prov. viii. 15, 16; Rom. xiii. 1, 2, 4.}\emph{ in quo quidem gerendo, ut pietatem præcipue, justitiam, ac pacem secundum salubres cujusque Reipublicæ leges tueri debent,}\footnote{Psa. ii. 10–12; 1 Tim. ii. 2; Psa. lxxxii. 3, 4; 2 Sam. xxiii. 3; 1 Pet. ii. 13.}\emph{ ita quo illum finem consequantur, licitum est iis vel hodie sub Novo Testamento in causis justis ac necessariis bellum gerere.}\footnote{Luke iii. 14; Rom. xiii. 4; Matt. viii. 9, 10; Acts x. 1, 2; Rev. xvii. 14, 16.}

III. The civil magistrate may not assume to himself the administration of the Word and Sacraments, or the power of the keys of the kingdom of heaven:\footnote{2 Chron. xxvi. 18 with Matt. xviii. 17 and xvi. 19; 1 Cor. xii. 28, 29; Eph. iv. 11, 12; 1 Cor. iv. 1, 2; Rom. x. 15; Heb. v. 4.} yet he hath authority, and it is his duty to take order, that unity and peace be preserved in the Church, that the truth of God be kept pure and entire, that all blasphemies and heresies be suppressed, all corruptions and abuses in worship and discipline prevented or reformed, and all the ordinances of God duly settled, administered, and observed.\footnote{Isa. xlix. 23; Psa. cxxii. 9; Ezra vii. 23–28; Lev. xxiv. 16; Deut. xiii. 5, 6, 12; 2 Kings xviii. 4; 1 Chron. xiii. 1–9; 2 Kings xxiii. 1–26; 2 Chron. xxxiv. 33; 2 Chron. xv. 12, 13.} For the better effecting whereof he hath power to call synods, to be present at them, and to provide that whatsoever is transacted in them be according to the mind of God.\footnote{2 Chron. xix. 8–11; chaps. xxix. and xxx.; Matt. ii. 4, 5.}

III. \emph{Magistratui Civili verbi et sacramentorum administrationem, aut clavium regni cœlorum potestatem assumere sibi non est licitum:}\footnote{2 Chron. xxvi. 18 with Matt. xviii. 17 and xvi. 19; 1 Cor. xii. 28, 29; Eph. iv. 11, 12; 1 Cor. iv. 1, 2; Rom. x. 15; Heb. v. 4.}\emph{ nihilo tamen minus et jure potest ille, eique incumbit providere ut Ecclesiæ unitas ac tranquillitas conservetur, ut veritas Dei pura et integra custodiatur, ut supprimantur blasphemiæ omnes, hæresesque, ut in cultu ac disciplina omnes corruptelæ ac abusus aut præcaveantur aut reformentur, omnia denique instituta divina, ut rite statuminentur, administrentur, observentur.}\footnote{Isa. xlix. 23; Psa. cxxii. 9; Ezra vii. 23–28; Lev. xxiv. 16; Deut. xiii. 5, 6, 12; 2 Kings xviii. 4; 1 Chron. xiii. 1–9; 2 Kings xxiii. 1–26; 2 Chron. xxxiv. 33; 2 Chron. xv. 12, 13.}\emph{ Quæ omnia quo melius præstare possit, potestatem habet tum Synodos convocandi, tum ut ipsis intersit, prospiciatque, ut quicquid in iis transigatur sit menti divinæ consentaneum.}\footnote{2 Chron. xix. 8–11; chaps. xxix. and xxx.; Matt. ii. 4, 5.}

The above section is changed in the American revision, and adapted to the separation of Church and State, as follows:

[III. \emph{Civil magistrates may not assume to themselves the administration of the Word and Sacraments }(2 Chron. xxvi. 18)\emph{; or the power of the keys of the kingdom of heaven }(Matt. xvi. 19; 1 Cor. iv. 1, 2)\emph{; or, in the least, interfere in matters of faith }(John xviii. 36; Mal. ii. 7; Acts v. 29)\emph{. Yet as nursing fathers, it is the duty of civil magistrates to protect the Church of our common Lord, without giving the preference to any denomination of Christians above the rest, in such a manner that all ecclesiastical persons whatever } \emph{shall enjoy the full, free, and unquestioned liberty of discharging every part of their sacred functions, without violence or danger }(Isa. xlix. 23)\emph{. And, as Jesus Christ hath appointed a regular government and discipline in his Church, no law of any commonwealth should interfere with, let, or hinder, the due exercise thereof, among the voluntary members of any denomination of Christians, according to their own profession and belief }(Psa. cv. 15; Acts xviii. 14–16)\emph{. It is the duty of civil magistrates to protect the person and good name of all their people, in such an effectual manner as that no person be suffered, either upon pretence of religion or infidelity, to offer any indignity, violence, abuse, or injury to any other person whatsoever: and to take order, that all religious and ecclesiastical assemblies be held without molestation or disturbance }(2 Sam. xxiii. 3; 1 Tim. ii. 1; Rom. xiii. 4).]

IV. It is the duty of people\footnote{[Am. ed. reads \emph{of the people}.]} to pray for magistrates,\footnote{1 Tim. ii. 1, 2.} to honor their persons,\footnote{1 Pet. ii. 17.} to pay them tribute and other dues,\footnote{Rom. xiii. 6, 7.} to obey their lawful commands, and to be subject to their authority, for conscience' sake.\footnote{Rom. xiii. 5; Tit. i. 3.} Infidelity or difference in religion doth not make void the magistrate's just and legal authority, nor free the people from their due obedience to him:\footnote{1 Pet. ii. 13, 14, 16.} from which ecclesiastical persons are not exempted;\footnote{Rom. xiii. 1; 1 Kings ii. 35; Acts xxv. 9–11; 2 Pet. ii. 1, 10, 11; Jude 8–11.} much less hath the Pope any power or jurisdiction over them in their dominions, or over any of their people; and least of all to deprive them of their dominions or lives, if he shall judge them to be

IV. \emph{Debet populus pro Magistratibus preces fundere,}\footnote{1 Tim. ii. 1, 2.}\emph{ personas eorum honore prosequi,}\footnote{1 Pet. ii. 17.}\emph{ tributa aliaque eis debita persolvere,}\footnote{Rom. xiii. 6, 7.}\emph{ obtemperare licitis eorum mandatis, ac propter conscientiam subjici illorum authoritati;}\footnote{Rom. xiii. 5; Tit. i. 3.}\emph{ quæ si justa sit ac legitima, non eam illorum infidelitas, non religio diversa cassam reddit, neque populum liberat a debitæ, illis obedientiæ præstatione,}\footnote{1 Pet. ii. 13, 14, 16.}\emph{ qua viri quidem Ecclesiastici non eximuntur,}\footnote{Rom. xiii. 1; 1 Kings ii. 35; Acts xxv. 9–11; 2 Pet. ii. 1, 10, 11; Jude 8–11.}\emph{ multo minus in ipsos magistratus, intra ditionem suam, ant ex eorum populo quemvis potestatem ullam habet aut jurisdictionem Papa Romanus, minime vero omnium vita illos aut principatu exuendi, si ipse }

heretics, or upon any other pretense whatsoever.\footnote{2 Thess. ii. 4; Rev. xiii. 15–17.}

\emph{scilicet eos hæreticos esse judicaverit, vel etiam alio prætextu quoviscunque.}\footnote{2 Thess. ii. 4; Rev. xiii. 15–17.}

\xsection{Chapter XXIV. Of Marriage and Divorce}

I. Marriage is to be between one man and one woman: neither is it lawful for any man to have more than one wife, nor for any woman to have more than one husband at the same time.\footnote{Gen. ii. 24; Matt. xix. 5, 6; Prov. ii. 17; [Am. ed. 1 Cor. vii. 2; Mark x. 6–9].}

I. \emph{Conjugium inter unum virum ac fœminam unam contrahi debet; neque viro ulli uxores plures, nec ulli fœminæ ultra unum maritum eodem tempore habere licet.}\footnote{Gen. ii. 24; Matt. xix. 5, 6; Prov. ii. 17; [Am. ed. 1 Cor. vii. 2; Mark x. 6–9].}

II. Marriage was ordained for the mutual help of husband and wife;\footnote{Gen. ii. 18.} for the increase of mankind with a legitimate issue, and of the Church with an holy seed;\footnote{Mal. ii. 15.} and for preventing of uncleanness.\footnote{1 Cor. vii. 2, 9.}

II. \emph{Conjugium erat institutum, cum propter mariti uxorisque auxilium mutuum,}\footnote{Gen. ii. 18.}\emph{ tum propter humani generis prole legitima, Ecclesiæqeu sancto semine incrementum,}\footnote{Mal. ii. 15.}\emph{ tum vero etiam ad impudicitiam declinandam.}\footnote{1 Cor. vii. 2, 9.}

III. It is lawful for all sorts of people to marry who are able with judgment to give their consent.\footnote{Heb. xiii. 4; 1 Tim. iv. 3; 1 Cor. vii. 36–38; Gen. xxiv. 57, 58.} Yet it is the duty of Christians to marry only in the Lord.\footnote{1 Cor. vii. 39.} And, therefore, such as profess the true reformed religion should not marry with infidels, Papists, or other idolaters: neither should such as are godly be unequally yoked, by marrying with such as are notoriously wicked in their life, or maintain damnable heresies.\footnote{Gen. xxxiv. 14; Exod. xxxiv. 16; Deut. vii. 3, 4; 1 Kings xi. 4; Neh. xiii. 25–27; Mal. ii. 11, 12; 2 Cor. vi. 14.}

III. \emph{Matrimonio jungi cuivis hominum generi licitum est, qui consensum suum præbere valent cum judicio;}\footnote{Heb. xiii. 4; 1 Tim. iv. 3; 1 Cor. vii. 36–38; Gen. xxiv. 57, 58.}\emph{ Veruntamen solum in Domino connubia inire debent Christiani;}\footnote{1 Cor. vii. 39.}\emph{ proindeque quotquot religionem veram reformatamque profitentur, non debent Infidelibus, Papistis, aut aliis quibuscunque idololatris connubio sociari; neque sane debent qui pii sunt impari jugo copulari, conjugium cum illis contrahendo qui aut improbitate vitæ sunt notabiles, aut damnabiles tuentur hæreses.}\footnote{Gen. xxxiv. 14; Exod. xxxiv. 16; Deut. vii. 3, 4; 1 Kings xi. 4; Neh. xiii. 25–27; Mal. ii. 11, 12; 2 Cor. vi. 14.};

IV. Marriage ought not to be within the degrees of consanguinity or affinity forbidden in the Word;\footnote{Lev. chap. xviii.; 1 Cor. v. 1; Amos ii. 7.} nor can such incestuous marriages ever be made lawful by any law of man, or consent of parties, so as those persons may live together, as man and wife.\footnote{Mark vi. 18; Lev. xviii. 24–28.} The man may not marry any of his wife's kindred nearer in blood than he may of his own, nor the woman of her husband's kindred nearer in blood than of her own.\footnote{Lev. xx. 19–21.}

IV. \emph{Connubia intra consanguinitatis affinitatisque gradus in verbo Dei vetitos iniri non est licitum;}\footnote{Lev. chap. xviii.; 1 Cor. v. 1; Amos ii. 7.}\emph{ neque possunt ejusmodi incesta conjugia quavis aut humana lege, aut consensione partium fieri legitima, adeo ut personis illis ad instar mariti et uxoris liceat unquam cohabitare.}\footnote{Mark vi. 18; Lev. xviii. 24–28.}\emph{ Non licet viro e cognatione uxoris suæ ducere, quam si æque seipsum attingeret sanguine, ducere non liceret; sicuti nec fœminæ licet viro nubere a mariti sui sanguine minus, quam a suo liceret, alieno.}\footnote{Lev. xx. 19–21.}

V. Adultery or fornication, committed after a contract, being detected before marriage, giveth just occasion to the innocent party to dissolve that contract.\footnote{Matt. i. 18–20.} In the case of adultery after marriage, it is lawful for the innocent party to sue out a divorce,\footnote{Matt. v. 31, 32.} and after the divorce to marry another, as if the offending party were dead.\footnote{Matt. xix. 9; Rom. vii. 2, 3.}

V. \emph{Adulterium aut scortatio si admittatur post sponsalia, ac ante conjugium detegatur, personæ innocenti justam præbet occasionem contractum illum dissolvendi;}\footnote{Matt. i. 18–20.}\emph{ quod si adulterium post conjugium admittatur, licebit parti innocenti divortium lege postulare ac obtinere;}\footnote{Matt. v. 31, 32.}\emph{ atque quidem post factum divortium conjugio alteri sociari, perinde acsi mortua esset persona illa quæ conjugii fidem violabat.}\footnote{Matt. xix. 9; Rom. vii. 2, 3.}

VI. Although the corruption of man be such as is apt to study arguments, unduly to put asunder those whom God hath joined together in marriage; yet nothing but adultery, or such willful desertion as can no way be remedied by

VI. \emph{Quamvis ea sit hominis corruptio ut proclivis sit ad excogitandum argumenta, indebite illos quos Deus connubio junxit dissociandi; nihilominus tamen extra adulterium ac desertionem ita obstinatam, ut cui nullo remedio, nec ab Ecclesia nec a }

the Church or civil magistrate, is cause sufficient of dissolving the bond of marriage;\footnote{Matt. xix. 8, 9; 1 Cor. vii. 15; Matt. xix. 6.} wherein a public and orderly course of proceeding is to be observed; and the persons concerned in it, not left to their own wills and discretion in their own case.\footnote{Deut. xxiv. 1–4; [Am. ed. Ezra x. 3].}

\emph{Magistratu civili subveniri possit, sufficiens causa nulla esse potest conjugii vinculum dissolvendi.}\footnote{Matt. xix. 8, 9; 1 Cor. vii. 15; Matt. xix. 6.}\emph{ Atque hac quidem in re procedendi ordo publicus et regularis est observandus, nec personæ illæ, quarum jus agitur, sunt suo arbitrio judiciove in causa propria permittendæ.}\footnote{Deut. xxiv. 1–4; [Am. ed. Ezra x. 3].}

\xsection{Chapter XXV. Of the Church}

I. The catholic or universal Church, which is invisible, consists of the whole number of the elect, that have been, are, or shall be gathered into one, under Christ the head thereof; and is the spouse, the body, the fullness of him that filleth all in all.\footnote{Eph. i. 10, 22, 23; v. 23, 27, 32; Col. i. 18.}

I. \emph{Catholica sive Universalis Ecclesia ea quæ est invisibilis constat e toto electorum numero, quotquot fuerunt, sunt, aut erunt unquam in unum collecti, sub Christo ejusdem Capite; estque sponsa, corpus ac plenitudo ejus qui implet omnia in omnibus.}\footnote{Eph. i. 10, 22, 23; v. 23, 27, 32; Col. i. 18.}

II. The visible Church, which is also catholic or universal under the gospel (not confined to one nation as before under the law) consists of all those, throughout the world, that profess the true religion,\footnote{1 Cor. i. 2; xii. 12, 13; Psa. ii. 8; Rev. vii. 9; Rom. xv. 9–12.} and of\footnote{[Am. ed. \emph{together with}, instead of \emph{and of}.]} their children;\footnote{1 Cor. vii. 14; Acts ii. 39; Ezek. xvi. 20, 21; Rom. xi. 16; Gen. iii. 15; xvii. 7; [Am. ed. Gal. iii. 7, 9, 14; Rom. iv. throughout].} and is the kingdom of the Lord Jesus Christ,\footnote{Matt. xiii. 47; Isa. ix. 7.} the house and family of God,\footnote{Eph. ii. 19; iii. 15; [Am. ed. Prov. xxix. 18].} out of which there is no ordinary possibility of salvation.\footnote{Acts ii. 47.}

II.\emph{ Ecclesia visibilis }(\emph{quæ etiam sub Evangelio, Catholica est et universalis, non autem unius gentis finibus, ut pridem sub lege, circumscripta})\emph{ ex iis omnibus constat, undecunque terrarum sint, qui veram religionem profitentur,}\footnote{1 Cor. i. 2; xii. 12, 13; Psa. ii. 8; Rev. vii. 9; Rom. xv. 9–12.}\emph{ una cum eorundem liberis;}\footnote{1 Cor. vii. 14; Acts ii. 39; Ezek. xvi. 20, 21; Rom. xi. 16; Gen. iii. 15; xvii. 7; [Am. ed. Gal. iii. 7, 9, 14; Rom. iv. throughout].}\emph{ estque Regnum Domini Jesu Christi,}\footnote{Matt. xiii. 47; Isa. ix. 7.}\emph{ Domus et familla Dei,}\footnote{Eph. ii. 19; iii. 15; [Am. ed. Prov. xxix. 18].}\emph{ extra quam quidem ordinarie fieri nequit ut quivis salutem consequatur.}\footnote{Acts ii. 47.}

III. Unto this catholic visible Church Christ hath given the ministry, oracles, and ordinances of God, for the gathering and perfecting of the saints, in this life, to the end of the world: and doth by his own presence and Spirit, according to his promise, make them effectual thereunto.\footnote{1 Cor. xii. 23; Eph. iv. 11–13; Matt. xxviii. 19, 20; Isa. lix. 21.}

III. \emph{Catholicæ huic Ecclesiæ visibili dedit Christus ministrorum ordinem, oracula, ac instituta Dei ad sanctos usque ad finem mundi in hac vita colligendos simul et perficiendos; in quem finem præsentia sua, spirituque secundum ipsius promissionem, eadem reddit efficacia.}\footnote{1 Cor. xii. 23; Eph. iv. 11–13; Matt. xxviii. 19, 20; Isa. lix. 21.}

IV. This catholic Church hath been sometimes more, sometimes less visible.\footnote{Rom. xi. 3, 4; Rev. xii. 6, 14; [Am. ed. Acts ix. 31].} And particular churches, which are members thereof, are more or less pure, according as the doctrine of the gospel is taught and embraced, ordinances administered, and public worship performed more or less purely in them.\footnote{Rev. chaps. ii. and iii.; 1 Cor. v. 6, 7.}

IV. \emph{Ecclesia hæc Catholica extitit quandoque magis quandoque minus visibilis.}\footnote{Rom. xi. 3, 4; Rev. xii. 6, 14; [Am. ed. Acts ix. 31].}\emph{ Ecclesiæ autem particulares }(\emph{quæ sunt illius membra})\emph{ eo magis minusve puræ sunt, qui majori aut minori cum puritate in iis docetur excipiturque Evangelii doctrina, administrantur divina instituta, cultusque publicus celebratur.}\footnote{Rev. chaps. ii. and iii.; 1 Cor. v. 6, 7.}

V. The purest churches under heaven are subject both to mixture and error;\footnote{1 Cor. xiii. 12; Rev. chaps. ii. and iii.; Matt. xiii. 24–30, 47.} and some have so degenerated as to become no churches of Christ, but synagogues of Satan.\footnote{Rev. xviii. 2; Rom. xi. 18–22.} Nevertheless, there shall be always a Church on earth to worship God according to his will.\footnote{Matt. xvi. 18; Psa. lxxii. 17; cii. 28; Matt. xxviii. 19, 20.}

V. \emph{Purissimæ omnium quæ in terris sunt Ecclesiæ, cum mixturæ tum etiam errori sunt obnoxiæ,}\footnote{1 Cor. xiii. 12; Rev. chaps. ii. and iii.; Matt. xiii. 24–30, 47.}\emph{ eousque autem nonnullæ degenerarunt, ut ex Ecclesiis Christi factæ demum sint ipsius Satanæ Synagogæ;}\footnote{Rev. xviii. 2; Rom. xi. 18–22.}\emph{ nihilominus tamen nunquam deerit in terris Ecclesiæ, quæ Deum colat secundum ipsius voluntatem.}\footnote{Matt. xvi. 18; Psa. lxxii. 17; cii. 28; Matt. xxviii. 19, 20.}

VI. There is no other head of the Church but the Lord Jesus Christ:\footnote{Col. i. 18; Eph. i. 22.} nor can the Pope of Rome, in any

VI. \emph{Ecclesiæ caput extra unum Dominum Jesum Christum nullum est;}\footnote{Col. i. 18; Eph. i. 22.}\emph{ nec ullo sensu caput ejus esse }

sense be head thereof; but is that Antichrist, that man of sin and son of perdition, that exalteth himself in the Church against Christ, and all that is called God.\footnote{Matt. xxiii. 8–10; 2 Thess. ii. 3, 4, 8, 9; Rev. xiii. 6.}

\emph{potest Papa Romanus, qui est insignis ille Antichristus, homo ille peccati et perditionis filius; in Ecclesia semet efferens adversus Christum, et supra quicquid dicitur Deus.}\footnote{Matt. xxiii. 8–10; 2 Thess. ii. 3, 4, 8, 9; Rev. xiii. 6.}

\xsection{Chapter XXVI. Of the Communion of Saints}

I. All saints that are united to Jesus Christ their head, by his Spirit and by faith, have fellowship with him in his graces, sufferings, death, resurrection, and glory:\footnote{1 John i. 3; Eph. iii. 16–19; John i. 16; Eph. ii. 5, 6; Phil. iii. 10; Rom. vi. 5, 6; 2 Tim. ii. 12.} and being united to one another in love, they have communion in each other's gifts and graces,\footnote{Eph. iv. 15, 16; 1 Cor. xii. 7; iii. 21–23; Col. ii. 19.} and are obliged to the performance of such duties, public and private, as do conduce to their mutual good, both in the inward and outward man.\footnote{1 Thess. v. 11, 14; Rom. i. 11, 12, 14; 1 John iii. 16–18; Gal. vi. 10.}

I. \emph{Sancti omnes, qui capiti suo Jesu Christo per Spiritum ejus ac per fidem uniuntur, gratiarum ejus, perpessionum, mortis, resurrectionis ac gloriæ ejus habent communionem;}\footnote{1 John i. 3; Eph. iii. 16–19; John i. 16; Eph. ii. 5, 6; Phil. iii. 10; Rom. vi. 5, 6; 2 Tim. ii. 12.}\emph{ atque inde etiam amore conjuncti sibimet invicem mutuam donorum suorum gratiarumque societatem quandam ineunt,}\footnote{Eph. iv. 15, 16; 1 Cor. xii. 7; iii. 21–23; Col. ii. 19.}\emph{ ac ad ejusmodi officia præstanda publica et privata obligantur, quæ ad mutuum eorum bonum conducant, cum quoad internum tum etiam quoad externum hominem.}\footnote{1 Thess. v. 11, 14; Rom. i. 11, 12, 14; 1 John iii. 16–18; Gal. vi. 10.}

II. Saints, by profession, are bound to maintain an holy fellowship and communion in the worship of God, and in performing such other spiritual services as tend to their mutual edification;\footnote{Heb. x. 24, 25; Acts ii. 42, 46; Isa. ii. 3; 1 Cor. xi. 20.} as also in relieving each other in outward things, according to their several abilities and necessities.

II. \emph{Qui sanctos sese profitentur, sanctam illi societatem et communionem inire tenentur et conservare, cum in divino cultu, tum alia officia spiritualia præstando, quæ ad mutuam eorum ædificationem conferre possint;}\footnote{Heb. x. 24, 25; Acts ii. 42, 46; Isa. ii. 3; 1 Cor. xi. 20.}\emph{ Quin etiam porro sublevando se mutuo in rebus externis, pro ratione cujusque vel facultatum }

Which communion, as God offereth opportunity, is to be extended unto all those who, in every place, call upon the name of the Lord Jesus.\footnote{Acts ii. 44, 45; 1 John iii. 17; 2 Cor. chaps. viii. and ix.; Acts xi. 29, 30.}

\emph{vel indigentiæ. Quæ quidem communio, prout opportunitatem Deus obtulerit, est ad eos omnes, qui ubivis locorum Domini Jesu nomen invocant, extendenda.}\footnote{Acts ii. 44, 45; 1 John iii. 17; 2 Cor. chaps. viii. and ix.; Acts xi. 29, 30.}

III. This communion which the saints have with Christ, doth not make them in anywise partakers of the substance of his Godhead, or to be equal with Christ in any respect: either of which to affirm is impious and blasphemous.\footnote{Col. i. 18, 19; 1 Cor. viii. 6; Isa. xlii. 8; 1 Tim. vi. 15, 16; Psa. xlv. 7 with Heb. i. 8, 9.} Nor doth their communion one with another, as saints, take away or infringe the title or propriety\footnote{[Am. ed. \emph{property}.]} which each man hath in his goods and possessions.\footnote{Exod. xx. 15; Eph. iv. 28; Acts v. 4.}

III. \emph{Hæc autem communio qua sancti cum Christo potiuntur, eos substantiæ Deitatis ejus neutiquam reddit participes, nec ullo respectu æquales Christo: Quorum utrumvis affirmare impium est ac blasphemum;}\footnote{Col. i. 18, 19; 1 Cor. viii. 6; Isa. xlii. 8; 1 Tim. vi. 15, 16; Psa. xlv. 7 with Heb. i. 8, 9.}\emph{ neque sane communio illa, quæ iis secum mutuo quatenus sanctis intercedit; cujusquam ad bona et possessiones suas jus privatum vel tollit vel imminuit.}\footnote{Exod. xx. 15; Eph. iv. 28; Acts v. 4.}

\xsection{Chapter XXVII. Of the Sacraments}

I. Sacraments are holy signs and seals of the covenant of grace,\footnote{Rom. iv. 11; Gen. xvii. 7, 10.} immediately instituted by God,\footnote{Matt. xxviii. 19; 1 Cor. xi. 23.} to represent Christ and his benefits, and to confirm our interest in him:\footnote{1 Cor. x. 16; xi. 25, 26; Gal. iii. 27.} as also to put a visible difference between those that belong unto the Church and the rest of the world;\footnote{Rom. xv. 8; Exod. xii. 48; Gen. xxxiv. 14; [Am. ed. 1 Cor. x. 21].} and solemnly to engage them to the service of God in Christ, according to his Word.\footnote{Rom. vi. 3, 4; 1 Cor. x. 16, 21.}

I. \emph{Sacramenta sunt fœderis gratiæ signa sacra et sigilla,}\footnote{Rom. iv. 11; Gen. xvii. 7, 10.}\emph{ immediate a Deo instituta,}\footnote{Matt. xxviii. 19; 1 Cor. xi. 23.}\emph{ ad Christum ejusque beneficia repræsentandum, ad jus nostrum in illo confirmandum,}\footnote{1 Cor. x. 16; xi. 25, 26; Gal. iii. 27.}\emph{ prout etiam ad illos qui Ecclesiam spectant a reliquis illis qui sunt e mundo, visibili discrimine separandum,}\footnote{Rom. xv. 8; Exod. xii. 48; Gen. xxxiv. 14; [Am. ed. 1 Cor. x. 21].}\emph{ utque ii solenniter devinciantur ad obedientiam et cultum Deo in Christo juxta verbum ejus exhibendum.}\footnote{Rom. vi. 3, 4; 1 Cor. x. 16, 21.}

II. There is in every sacrament

II. \emph{In Sacramento quolibet est inter }

a spiritual relation or sacramental union, between the sign and the thing signified; whence it comes to pass that the names and the\footnote{Am. ed. omits \emph{the.}} effects of the one are attributed to the other.\footnote{Gen. xvii. 10; Matt. xxvi. 27, 28; Tit. iii. 5.}

\emph{signum et rem significatam relatio quædam spiritualis, sive Sacramentalis unio; unde fit ut alterius nomina et effectus alteri quandoque tribuantur.}\footnote{Gen. xvii. 10; Matt. xxvi. 27, 28; Tit. iii. 5.}

III. The grace which is exhibited in or by the sacraments, rightly used, is not conferred by any power in them; neither doth the efficacy of a sacrament depend upon the piety or intention of him that doth administer it,\footnote{Rom. ii. 28, 29; 1 Pet. iii. 21.} but upon the work of the Spirit,\footnote{Matt. iii. 11; 1 Cor. xii. 13.} and the word of institution, which contains, together with a precept authorizing the use thereof, a promise of benefit to worthy receivers.\footnote{Matt. xxvi. 27, 28; xxviii. 19, 20.}

III. \emph{Quæ in Sacramentis sive per ea rite adhibita exhibetur gratia, per vim aliquam iis intrinsecam non confertur, neque ex intentione vel pietate adininistrantis pendent Sacramenti vis ac efficacia;}\footnote{Rom. ii. 28, 29; 1 Pet. iii. 21.}\emph{ verum ex operatione Spiritus,}\footnote{Matt. iii. 11; 1 Cor. xii. 13.}\emph{ ac verbo institutionis, quod complectitur cum præceptum, unde celebrandi Sacramenti potestas fit, tum etiam promissionem de beneficiis digne percipientibus exhibendis.}\footnote{Matt. xxvi. 27, 28; xxviii. 19, 20.}

IV. There be only two sacraments ordained by Christ our Lord in the gospel, that is to say, Baptism and the Supper of the Lord: neither of which may be dispensed by any but by a minister of the Word lawfully ordained.\footnote{Matt. xxviii. 19; 1 Cor. xi. 20, 23; iv. 1; Heb. v. 4.}

IV. \emph{Sacramenta duo tantum sunt a Christo Domino nostro in Evangelio instituta, Baptismus scilicet, et cœna Domini; quorum neutrum de debet nisi a ministro verbi legitime ordinato dispensari.}\footnote{Matt. xxviii. 19; 1 Cor. xi. 20, 23; iv. 1; Heb. v. 4.}

V. The sacraments of the Old Testament, in regard of the spiritual things thereby signified and exhibited, were, for substance, the same with those of the New.\footnote{1 Cor. x. 1–4; [Am. ed. 1 Cor. v. 7, 8].}

V. \emph{Sacramenta Veteris Testamenti si res spirituales per ea significatas exhibitasque respiciamus, quoad substantiam eadem fuere cum his sub Novo.}\footnote{1 Cor. x. 1–4; [Am. ed. 1 Cor. v. 7, 8].}

\xsection{Chapter XXVIII. Of Baptism}

I. Baptism is a sacrament of the

I. \emph{Baptismus est sacramentum }

New Testament, ordained by Jesus Christ,\footnote{Matt. xxviii. 19; [Am. ed. Mark xvi. 16].} not only for the solemn admission of the party baptized into the visible Church,\footnote{1 Cor. xii. 13; [Am. ed. Gal. iii. 27, 28].} but also to be unto him a sign and seal of the covenant of grace\footnote{Rom. iv. 11 with Col. ii. 11, 12.}, of his ingrafting into Christ,\footnote{Gal. iii. 27; Rom. vi. 5.} of regeneration,\footnote{Tit. iii. 5.} of remission of sins,\footnote{Mark i. 4; [Am. ed. Acts ii. 38; xxii. 16].} and of his giving up unto God, through Jesus Christ, to walk in newness of life:\footnote{Rom. vi. 3, 4.} which sacrament is, by Christ's own appointment, to be continued in his Church until the end of the world.\footnote{Matt. xxviii. 19, 20.}

\emph{Novi Testamenti, a Jesu Christo institutum,}\footnote{Matt. xxviii. 19; [Am. ed. Mark xvi. 16].}\emph{ non solum propter solennem personæ baptizatæ in Ecclesiam visibilem admissionem,}\footnote{1 Cor. xii. 13; [Am. ed. Gal. iii. 27, 28].}\emph{ verum etiam ut signum eidem sit, et sigillum cum fœderis gratiæ,}\footnote{Rom. iv. 11 with Col. ii. 11, 12.}\emph{ tum suæ in Christum insitionis,}\footnote{Gal. iii. 27; Rom. vi. 5.}\emph{ regenerationis,}\footnote{Tit. iii. 5.}\emph{ remissionis peccatorum,}\footnote{Mark i. 4; [Am. ed. Acts ii. 38; xxii. 16].}\emph{ ac sui ipsius Deo per Christum dedicationis, ad ambulandum in vitæ novitate.}\footnote{Rom. vi. 3, 4.}\emph{ Quod quidem Sacramentum e Christi ipsius mandato est in Ecclesia ejus ad finem usque mundi retinendum.}\footnote{Matt. xxviii. 19, 20.}

II. The outward element to be used in this sacrament is water, wherewith the party is to be baptized in the name of the Father, and of the Son, and of the Holy Ghost, by a minister of the gospel lawfully called thereunto.\footnote{Matt. iii. 11; John i. 33; Matt. xxviii. 19, 20; [Am. ed. Acts x. 47; viii. 36, 38].}

II. \emph{Elementum externum in hoc Sacramento adhibendum est Aqua; qua baptizari debet admittendus, a ministro Evangelii legitime ad hoc vocato, in nomine Patris et filii et Spiritus Sancti.}\footnote{Matt. iii. 11; John i. 33; Matt. xxviii. 19, 20; [Am. ed. Acts x. 47; viii. 36, 38].}

III. Dipping of the person into the water is not necessary; but baptism is rightly administered by pouring or sprinkling water upon the person.\footnote{Heb. ix. 10, 19–22; Acts ii. 41; xvi. 33; Mark vii. 4.}

III. \emph{Baptizandi in aquam immersio necessaria non est; verum baptismus rite administratur aqua superfusa vel etiam inspersa baptizando.}\footnote{Heb. ix. 10, 19–22; Acts ii. 41; xvi. 33; Mark vii. 4.}

IV. Not only those that do actually profess faith in and obedience unto Christ,\footnote{Mark xvi. 15, 16; Acts viii. 37, 38.} but also the infants

IV. \emph{Non illi solum qui fidem in Christum eique se obedientes fore actu quidem profitentur,}\footnote{Mark xvi. 15, 16; Acts viii. 37, 38.}\emph{ verum }

of one or both believing parents are to be baptized.\footnote{Gen. xvii. 7, 9, with Gal. iii. 9, 14, and Col. ii. 11, 12, and Acts ii. 38, 39, and Rom. iv. 11, 12; 1 Cor. vii. 14; Matt. xxviii. 19; Mark x. 13–16; Luke xviii. 15; [Am. ed. Acts xvi. 14, 15, 33].}

\emph{etiam infantes qui a Parente vel altero vel utroque fideli procreantur, sunt baptizandi.}\footnote{Gen. xvii. 7, 9, with Gal. iii. 9, 14, and Col. ii. 11, 12, and Acts ii. 38, 39, and Rom. iv. 11, 12; 1 Cor. vii. 14; Matt. xxviii. 19; Mark x. 13–16; Luke xviii. 15; [Am. ed. Acts xvi. 14, 15, 33].}

V. Although it be a great sin to contemn or neglect this ordinance,\footnote{Luke vii. 30 with Exod. iv. 24–26.} yet grace and salvation are not so inseparably annexed unto it, as that no person can be regenerated or saved without it,\footnote{Rom. iv. 11; Acts x. 2, 4, 22, 31, 45, 47.} or that all that are baptized are undoubtedly regenerated.\footnote{Acts viii. 13, 23.}

V. \emph{Quamvis grave peccatum sit institutum hoc despicatui habere vel negligere;}\footnote{Luke vii. 30 with Exod. iv. 24–26.}\emph{ non tamen ei salus et gratia ita individue annectuntur, ut absque illo nemo unquam regenerari aut salvari possit,}\footnote{Rom. iv. 11; Acts x. 2, 4, 22, 31, 45, 47.}\emph{ aut quasi indubium omnino sit regenerari omnes qui baptizantur.}\footnote{Acts viii. 13, 23.}

VI. The efficacy of baptism is not tied to that moment of time wherein it is administered;\footnote{John iii. 5, 8.} yet, notwithstanding, by the right use of this ordinance the grace promised is not only offered, but really exhibited and conferred by the Holy Ghost, to such (whether of age or infants) as that grace belongeth unto, according to the counsel of God's own will, in his appointed time.\footnote{Gal. iii. 27; Tit. iii. 5; Eph. v. 25, 26; Acts ii. 38, 41.}

VI. \emph{Baptismi efficacia ei temporis momento quo administratur non adstringitur.}\footnote{John iii. 5, 8.}\emph{ Nihilominus tamen, usu debito hujus instituti non offertur solum promissa gratia, verum etiam omnibus }(\emph{tam infantibus quam adultis})\emph{ ad quos gratia illa e consilio Divinæ voluntatis pertinet, per Spiritum Sanctum in tempore suo constituto realiter confertur et exhibetur.}\footnote{Gal. iii. 27; Tit. iii. 5; Eph. v. 25, 26; Acts ii. 38, 41.}

VII. The sacrament of baptism is but once to be administered to any person.\footnote{Tit. iii. 5.}

VII. \emph{Sacramentum Baptismi eidem personæ non est nisi semel administrandum.}\footnote{Tit. iii. 5.}

\xsection{Chapter XXIX. Of the Lord's Supper}

I. Our Lord Jesus, in the night wherein he was betrayed, instituted

I. \emph{Dominus noster Jesus eadem qua prodebatur nocte instituit corporis }

the sacrament of his body and blood, called the Lord's Supper, to be observed in his Church, unto the end of the world, for the perpetual remembrance of the sacrifice of himself in his death, the sealing all benefits thereof unto true believers, their spiritual nourishment and growth in him, their further engagement in, and to all duties which they owe unto him; and to be a bond and pledge of their communion with him, and with each other, as members of his mystical body.\footnote{1 Cor. xi. 23–26; x. 16, 17, 21; xii. 13.}

\emph{et sanguinis sui sacramentum, }Cœnam Domini\emph{ quam dicimus, in Ecclesia sua ad finem usque mundi celebrandum, in perpetuam memoriam sacrificii sui ipsius in morte sua oblati, et ad beneficia istius omnia vere fidelibus obsignandum; in eorum item alimentum ac incrementum in Christo spirituale; quoque ad officia cuncta præstanda, prius quidem illi debita, arctiori adhuc nodo tenerentur; ut vinculum denique ac pignus foret communionis illius quæ iis cum Christo et secum ipsis mutuo, tanquam corporis ipsius mystici membris, intercedit.}\footnote{1 Cor. xi. 23–26; x. 16, 17, 21; xii. 13.}

II. In this sacrament Christ is not offered up to his Father, nor any real sacrifice made at all for remission of sins of the quick or dead,\footnote{Heb. ix. 22, 25, 26, 28.} but only a commemoration of that one offering up of himself, by himself, upon the cross, once for all, and a spiritual oblation of all possible praise unto God for the same;\footnote{1 Cor. xi. 24–26; Matt. xxvi. 26, 27; [Am. ed. Luke xxii. 19, 20].} so that the Popish sacrifice of the mass, as they call it, is most abominably injurious to Christ's one only sacrifice, the alone propitiation for all the sins of the elect.\footnote{Heb. vii. 23, 24, 27; x. 11, 12, 14, 18.}

II. \emph{In hoc Sacramento non Patri suo offertur Christus, sed neque inibi fit reale aliquod sacrificium ad peccatorum remissionem vivis aut mortuis procurandam;}\footnote{Heb. ix. 22, 25, 26, 28.}\emph{ verum unicæ istius oblationis, qua Christus semet ipsum ipse in cruce, et quidem omnino semel obtulit, commemoratio solum; una cum spirituali propterea laudis omnimodæ Deo redditæ oblatione.}\footnote{1 Cor. xi. 24–26; Matt. xxvi. 26, 27; [Am. ed. Luke xxii. 19, 20].}\emph{ Unde Pontificiorum istud sacrificium Missæ }(uti loqui amant)\emph{ plane detestandum sit oportet, utpote maxime injuriam uni illi unicoque Christi sacrificio, quod quidem unica est pro peccatis electorum universus propitiatio.}\footnote{Heb. vii. 23, 24, 27; x. 11, 12, 14, 18.}

III. The Lord Jesus hath, in this ordinance, appointed his ministers to declare his word of institution to the people, to pray, and bless the elements of bread and wine, and thereby to set them apart from a common to an holy use; and to take and break the bread, to take the cup, and (they communicating also themselves) to give both to the communicants;\footnote{Matt. xxvi. 26–28, and Mark xiv. 22–24, and Luke xxii. 19, 20, with 1 Cor. xi. 23–27.} but to none who are not then present in the congregation.\footnote{Acts xx. 7; 1 Cor. xi. 20.}

III. \emph{In hoc suo instituto præcepit Dominus Jesus Ministris suis, verbum institutionis populo declarare, orare, ac elementis pani scilicet ac vino benedicere, eaque hac ratione a communi ad sacrum usum separare, quinetiam panem accipere et frangere; poculum item in manus accipere; atque }\{\emph{communicantibus una ipsis})\emph{ utrumque communicantibus exhibere,}\footnote{Matt. xxvi. 26–28, and Mark xiv. 22–24, and Luke xxii. 19, 20, with 1 Cor. xi. 23–27.}\emph{ nemini autem a congregatione tunc absenti.}\footnote{Acts xx. 7; 1 Cor. xi. 20.}

IV. Private masses, or receiving this sacrament by a priest, or any other, alone;\footnote{1 Cor. x. 6.} as likewise the denial of the cup to the people;\footnote{Mark iv. 23; 1 Cor. xi. 25–29.} worshiping the elements, the lifting them up, or carrying them about for adoration, and the reserving them for any pretended religious use, are all contrary to the nature of this sacrament, and to the institution of Christ.\footnote{Matt. xv. 9.}

IV. \emph{Missæ privatæ, sive perceptio hujusce Sacramenti a solo vel Sacerdote vel alio quovis;}\footnote{1 Cor. x. 6.}\emph{ prout etiam poculi a populo detensio,}\footnote{Mark iv. 23; 1 Cor. xi. 25–29.}\emph{ elementorum adoratio, quoque adorentur elevatio aut circumgestatio, ut et prætextu religiosi usus cujuscunque asservatio, sunt quidem omnia tum hujusce Sacramenti naturæ tum Christi institutioni plane contraria.}\footnote{Matt. xv. 9.}

V. The outward elements in this sacrament, duly set apart to the uses ordained by Christ, have such relation to him crucified, as that truly, yet sacramentally only, they are sometimes called by the name of the things they represent, to wit, the body and blood of Christ;\footnote{Matt. xxvi. 26–28.} albeit, in substance and nature, they still

V. \emph{In hoc Sacramento externa elementa ad usus a Christo institutos rite separata, ita ad eum crucifixum referuntur ut rerum quas repræsentat nominibus }(\emph{corporis nempe ac sanguinis Christi})\emph{ vere quidem, at Sacramentaliter tantum, sint nuncupata,}\footnote{Matt. xxvi. 26–28.}\emph{ manent siquidem adhuc quoad substantiam et naturam vere solumque }

remain truly, and only, bread and wine, as they were before.\footnote{1 Cor. xi. 26–28; Matt. xxvi. 29.}

\emph{panis ac vinum nihilo minus quam antea fuerant.}\footnote{1 Cor. xi. 26–28; Matt. xxvi. 29.}

VI. That doctrine which maintains a change of the substance of bread and wine, into the substance of Christ's body and blood (commonly called transubstantiation) by consecration of a priest, or by any other way, is repugnant, not to Scripture alone, but even to common-sense and reason; overthroweth the nature of the sacrament; and hath been, and is the cause of manifold superstitions, yea, of gross idolatries.\footnote{Acts iii. 21 with 1 Cor. xi. 24–26; Luke xxiv. 6, 39.}

VI. \emph{Doctrina illa quæ substantiæ panis ac vini in substantiam corporis et sanguinis Christi conversionem }(\emph{transubstantiatio vulgo dicitur})\emph{ sive illam per Sacerdotis consecrationem, sive quomodocunque demum fieri statuit, non scripturæ solum, verum etiam communi omnium sensui ac rationi adversatur, sacramenti naturam evertit, superstitionis multifariæ causa extitit atque etiamnum existit, imo vero et crassissimæ idololatriæ.}\footnote{Acts iii. 21 with 1 Cor. xi. 24–26; Luke xxiv. 6, 39.}

VII. Worthy receivers, outwardly partaking of the visible elements in this sacrament,\footnote{1 Cor. xi. 28; [Am. ed. 1 Cor. v. 7, 8].} do then also inwardly by faith, really and indeed, yet not carnally and corporally, but spiritually, receive and feed upon Christ crucified, and all benefits of his death: the body and blood of Christ being then not corporally or carnally in, with, or under the bread and wine; yet as really, but spiritually, present to the faith of believers in that ordinance, as the elements themselves are, to their outward senses.\footnote{1 Cor. x. 16; [Am. ed. 1 Cor. x. 3, 4].}

VII. \emph{Digne communicantes, Elementa in hoc Sacramento visibilia dum participant,}\footnote{1 Cor. xi. 28; [Am. ed. 1 Cor. v. 7, 8].}\emph{ una cum iis interne Christum crucifixum et beneficia mortis ejus universa revera et realiter }(\emph{modo, non carnali quidem aut corporeo, sed spirituali})\emph{ per fidem recipiunt eisque vescuntur. Corpus siquidem et sanguis Christi non corporeo aut carnali modo in, cum, vel sub pane ac vino; realiter tamen, ac spiritualiter credentium fidei in hoc instituto, non minus quam externis sensibus elementa ipsa, sunt præsentia.}\footnote{1 Cor. x. 16; [Am. ed. 1 Cor. x. 3, 4].}

VIII. Although ignorant and wicked men receive the outward elements in this sacrament, yet

VIII. \emph{Homines improbi et ignari externa licet in hoc sacramento percipere possint elementa, rem tamen }

they receive not the thing signified thereby; but by their unworthy coming thereunto are guilty of the body andj blood of the Lord, to their own damnation. Wherefore all ignorant and ungodly persons, as they are unfit to enjoy communion with him, so are they unworthy of the Lord's table, and can not, without great sin against Christ, while they remain such, partake of these holy mysteries,\footnote{1 Cor. xi. 27–29; 2 Cor. vi. 14–16; [Am. ed. 1 Cor. x. 21].} or be admitted thereunto.\footnote{1 Cor. v. 6, 7, 13; 2 Thess. iii. 6, 14, 15; Matt. vii. 6.}

\emph{per ea significatam non recipiunt; verum indigne illuc accedendo, rei fiunt corporis ac sanguinis Dominici ad sui ipsorum condemnationem. Quapropter homines impii et ignari prout communioni cum Deo potiundæ nullatenus sunt idonei, ita prorsus indigni sunt qui accedant ad mensam Domini; neque sine gravi in Christum peccato, possunt }(\emph{quamdiu tales esse non destiterint Sacra hæc mysteria participare;}\footnote{1 Cor. xi. 27–29; 2 Cor. vi. 14–16; [Am. ed. 1 Cor. x. 21].}\emph{ vel ad ea participandum admitti.}\footnote{1 Cor. v. 6, 7, 13; 2 Thess. iii. 6, 14, 15; Matt. vii. 6.}

\xsection{Chapter XXX. Of Church Censures}

I. The Lord Jesus, as king and head of his Church, hath therein appointed a government in the hand of Church officers, distinct from the civil magistrate.\footnote{Isa. ix. 6, 7; 1 Tim. v. 17; 1 Thess. v. 12; Acts xx. 17, 28; Heb. xiii. 7, 17, 24; 1 Cor. xii. 28; Matt. xxviii. 18–20; [Am. ed. Psa. ii. 6–9; John xviii. 36].}

I. \emph{Dominus Jesus quatenus Rex et caput Ecclesiæ suæ constituit in ea regimen, quod in officiariorum Ecclesiasticorum manu foret, distinctum a civili Magistratu.}\footnote{Isa. ix. 6, 7; 1 Tim. v. 17; 1 Thess. v. 12; Acts xx. 17, 28; Heb. xiii. 7, 17, 24; 1 Cor. xii. 28; Matt. xxviii. 18–20; [Am. ed. Psa. ii. 6–9; John xviii. 36].}

II. To these officers the keys of the kingdom of heaven are committed, by virtue whereof they have power respectively to retain and remit sins, to shut that kingdom against the impenitent, both by the Word and censures; and to open it unto penitent sinners, by the ministry of the gospel, and by absolution from censures, as occasion shall require.\footnote{Matt. xvi. 19; xviii. 17, 18; John xx. 21–23; 2 Cor. ii. 6–8.}

II. \emph{Officiariis hisce claves regni cœlorum sunt commissæ, quarum virtute obtinent potestatem peccata vel retinendi vel remittendi pro varia peccantium conditione; impœnitentibus quidem regnum illud tam per verbum quam per censuras occludendi, peccatoribus vero pœnitentibus tam evangelii ministerio quam absolutione a censuris idem aperiendi, prout occasio postulaverit.}\footnote{Matt. xvi. 19; xviii. 17, 18; John xx. 21–23; 2 Cor. ii. 6–8.}

III. Church censures are necessary for the reclaiming and gaining of offending brethren; for deterring of others from the\footnote{[Am. ed. omits \emph{the}.]} like offenses; for purging out of that leaven which might infect the whole lump; for vindicating the honor of Christ, and the holy profession of the gospel; and for preventing the wrath of God, which might justly fall upon the Church, if they should suffer his covenant, and the seals thereof, to be profaned by notorious and obstinate offenders.\footnote{1 Cor. chap. v.; 1 Tim. v. 20; Matt. vii. 6; 1 Tim. i. 20; 1 Cor. xi. 27 to the end, with Jude 23.}

III. \emph{Omnino necessariæ sunt censuræ Ecclesiasticæ, lucrandis fratribus delinquentibus eisque in viam reducendis, reliquis autem a similibus delictis deterrendis, fermento illi malo, ne totam massam inficiat, expurgando; ad honorem Christi et Sanctam Evangelii professionem vindicandum, ut prævertatur denique ira Dei, quæ merito in Ecclesiam accendi posset, si ipsius fœdus, hujusque sigilla ab insigniter ac pertinaciter delinquentibus impune profanari pateretur.}\footnote{1 Cor. chap. v.; 1 Tim. v. 20; Matt. vii. 6; 1 Tim. i. 20; 1 Cor. xi. 27 to the end, with Jude 23.}

IV. For the better attaining of these ends, the officers of the Church are to proceed by admonition, suspension from the Sacrament of the Lord's Supper for a season, and by excommunication from the Church, according to the nature of the crime and demerit of the person.\footnote{1 Thess. v. 12; 2 Thess. iii. 6, 14, 15; 1 Cor. v. 4, 5, 13; Matt. xviii. 17; Tit. iii. 10.}

IV. \emph{Quo melius autem hosce fines consequantur, procedere debent Ecclesiæ officiarii, admonendo, a Sacramento cœnæ Dominicæ ad tempus aliquod suspendendo, excommunicando denique ab Ecclesia, pro ratione criminis, atque personæ delinquentis merito.}\footnote{1 Thess. v. 12; 2 Thess. iii. 6, 14, 15; 1 Cor. v. 4, 5, 13; Matt. xviii. 17; Tit. iii. 10.}

\xsection{Chapter XXXI. Of Synods and Councils}

I. For the better government and further edification of the Church, there ought to be such assemblies as are commonly called synods or councils.\footnote{Acts xv. 2, 4, 6.}

I. \emph{Quo melius gubernari, ac ulterius ædificari possit Ecclesia, conventus ejusmodi fieri debent, quales vulgo Synodi et Concilia nuncupantur.}\footnote{Acts xv. 2, 4, 6.}

The American edition here adds the following:

[\emph{And it belongeth to the overseers and other rulers of the particular churches, by virtue of their office, and the power which Christ hath given them for edification, and not for destruction, to appoint such assemblies }(Acts xv.)\emph{; and to convene together in them, as often as they shall judge it expedient for the good of the Church }(Acts xv. 22, 23, 25).]

II. As magistrates may lawfully call a synod of ministers and other fit persons to consult and advise with about matters of religion;\footnote{Isa. xlix. 23; 1 Tim. ii. 1, 2; 2 Chron. xix. 8–11; chaps. xxix., xxx.; Matt. ii. 4, 5; Prov. xi. 14.} so, if magistrates be open enemies to the Church, the ministers of Christ, of themselves, by virtue of their office, or they, with other fit persons, upon delegation from their churches, may meet together in such assemblies.\footnote{Acts xv. 2, 4, 22, 23, 25.}\textsuperscript{\&}\footnote{[Am. ed. omits this whole section.]}

II. \emph{Quemadmodum licitum est Magistratibus Synodum Ministrorum aliorumque qui sunt idonei convocare, quibuscum de religionis rebus deliberent ac consultent:}\footnote{Isa. xlix. 23; 1 Tim. ii. 1, 2; 2 Chron. xix. 8–11; chaps. xxix., xxx.; Matt. ii. 4, 5; Prov. xi. 14.}\emph{ Ita si Magistratus fuerint Ecclesiæ hostes aperti, licebit Christi ministris a seipsis virtute officii, eisve cum aliis idoneis, accepta prius ab Ecclesiis suis delegatione, in istiusmodi conventibus congregari.}\footnote{Acts xv. 2, 4, 22, 23, 25.}

III. [II.] It belongeth to synods and councils, ministerially, to determine controversies of faith, and cases of conscience; to set down rules and directions for the better ordering of the public worship of God, and government of his Church; to receive complaints in cases of maladministration, and authoritatively to determine the same: which decrees and determinations, if consonant to the Word of God, are to be received with reverence and submission, not only for their agreement with the Word, but also for

III. \emph{Synodorum et Conciliorum est controversias fidei et conscientiæ casus, ministerialiter quidem, determinare; regulas ac præscripta quo melius publicus Dei cultus ejusque Ecclesiæ regimen ordinentur constituere; Querelas de mala administratione delatas admittere, deque iis authoritative decernere. Quæ quidem decreta et decisiones, modo verbo Dei consenserint, cum reverentia sunt ac summissione excipienda; Non quidem solum quod verbo Dei sint consentanea, verum etiam gratia potestatis ea constituentis, ut quæ }

the power whereby they are made, as being an ordinance of God, appointed thereunto in his Word.\footnote{Acts xv. 15, 19, 24, 27–31; xvi. 4; Matt. xviii. 17–20.}

\emph{sit ordinatio Dei id ad in verbo suo designata.}\footnote{Acts xv. 15, 19, 24, 27–31; xvi. 4; Matt. xviii. 17–20.}

IV. [III.] All synods or councils since the apostles' times, whether general or particular, may err, and many have erred; therefore they are not to be made the rule of faith or practice, but to be used as a help in both.\footnote{Eph. ii. 20; Acts xvii. 11; 1 Cor. ii. 5; 2 Cor. i. 24.}

IV. \emph{Synodi omnes sive concilia post Apostolorum tempora, seu generales sive particulares, errori sunt obnoxiæ, quin neque paucæ erraverunt. Proindeque fidei aut praxeos norma constituendæ non sunt, verum in utrisque auxilii loco adhibendæ.}\footnote{Eph. ii. 20; Acts xvii. 11; 1 Cor. ii. 5; 2 Cor. i. 24.}

V. [IV.] Synods and councils are to handle or conclude nothing but that which is ecclesiastical: and are not to intermeddle with civil affairs which concern the commonwealth, unless by way of humble petition in cases extraordinary; or by way of advice for satisfaction of conscience, if they be thereunto required by the civil magistrate.\footnote{Luke xii. 13, 14; John xviii. 36.}

V. \emph{Synodi et Concilia id solum quod Ecclesiam spectat tractare debent et concludere; neque civilibus negotiis, quæ rem publicam spectant ingerere se debent, nisi humiliter supplicando in casibus, si qui acciderint, extraordinariis; aut consulendo, quoties id ab eis postulat Magistratus civilis, nempe quo conscientiæ illius satisfiat.}\footnote{Luke xii. 13, 14; John xviii. 36.}

\xsection{Chapter XXXII. Of the State of Men after Death, and of the Resurrection of the Dead}\footnote{[Am. ed. has \emph{Man}.]}

I. The bodies of men, after death, return to dust, and see corruption;\footnote{Gen. iii. 19; Acts xiii. 36.} but their souls (which neither die nor sleep), having an immortal subsistence, immediately return to God who gave them.\footnote{Luke xxiii. 43; Eccles. xii. 7.} The souls of the righteous, being then made perfect in holiness, are received into the highest heavens, where they behold

I. \emph{Hominum corpora post mortem ad pulverem rediguntur, et corruptionem vident:}\footnote{Gen. iii. 19; Acts xiii. 36.}\emph{ At animæ illorum }(\emph{quæ quidem nec morientur nec obdormiunt})\emph{ ut quæ subsistentiam habent immortalem, ad Deum continuo earum datorem revertuntur.}\footnote{Luke xxiii. 43; Eccles. xii. 7.}\emph{ Animæ quidem Justorum iam tum perfecte sanctificatæ, cœlis supremis accipiuntur, }

the face of God in light and glory, waiting for the full redemption of their bodies:\footnote{Heb. xii. 23; 2 Cor. v. 1, 6, 8; Phil. i. 23, with Acts iii. 21 and Eph. iv. 10; [Am. ed. 1 John iii. 2].} and the souls of the wicked are cast into hell, where they remain in torments and utter darkness, reserved to the judgment of the great day.\footnote{Luke xvi. 23, 24; Acts i. 25; Jude 6, 7; 1 Pet. iii. 19.} Besides these two places for souls separated from their bodies, the Scripture acknowledgeth none.

\emph{ubi Dei faciem in lumine ac gloria intuentur, corporum suorum plenum redemtionem expectantes:}\footnote{Heb. xii. 23; 2 Cor. v. 1, 6, 8; Phil. i. 23, with Acts iii. 21 and Eph. iv. 10; [Am. ed. 1 John iii. 2].}\emph{ Animæ vero improborum conjiciuntur in Gehennam, ubi inter diros cruciatus in tenebris exterioribus conclusæ manent, ad judicium magni illius diei asservatæ.}\footnote{Luke xvi. 23, 24; Acts i. 25; Jude 6, 7; 1 Pet. iii. 19.}\emph{ Locum autem animabus a corpore solutis extra hosce duos Scriptura Sacra non agnoscit ullum.}

II. At the last day, such as are found alive shall not die, but be changed;\footnote{1 Thess. iv. 17; 1 Cor. xv. 51, 52.} and all the dead shall be raised up with the self-same bodies, and none other, although with different qualities, which shall be united again to their souls forever.\footnote{Job xix. 26, 27; 1 Cor. xv. 42–44.}

II. \emph{Novissimo illo die, qui comperientur in vivis non morientur quidem sed mutabuntur;}\footnote{1 Thess. iv. 17; 1 Cor. xv. 51, 52.}\emph{ qui mortui fuerint resuscitabuntur omnes, ipsissimis iis corporibus quibus viventes aliquando fungebantur, ac non aliis, utut qualitate differentibus; quæ denuo animabus quæque suis æterno conjugio unientur.}\footnote{Job xix. 26, 27; 1 Cor. xv. 42–44.}

III. The bodies of the unjust shall, by the power of Christ, be raised to dishonor; the bodies of the just, by his Spirit, unto honor, and be made conformable to his own glorious body.\footnote{Acts xxiv. 15; John v. 28, 29; 1 Cor. xv. 42; Phil. iii. 21.}

III. \emph{Injustorum corpora ad dedecus per potentiam Christi suscitabuntur; justorum autem corpora per spiritum ejus ad honorem, fientque hæc conformia corpori ipsius glorioso.}\footnote{Acts xxiv. 15; John v. 28, 29; 1 Cor. xv. 42; Phil. iii. 21.}

\xsection{Chapter XXXIII. Of the Last Judgment}

I. God hath appointed a day wherein he will judge the world in righteousness by Jesus Christ,\footnote{Acts xvii. 31.} to

I. \emph{Diem Deus designavit quo mundum in justitia judicabit per Jesum Christum;}\footnote{Acts xvii. 31.}\emph{ cui a Patre data est }

whom all power and judgment is given of the Father.\footnote{John v. 22, 27.} In which day, not only the apostate angels shall be judged,\footnote{1 Cor. vi. 3; Jude 6; 2 Pet. ii. 4.} but likewise all persons, that have lived upon earth, shall appear before the tribunal of Christ, to give an account of their thoughts, words, and deeds; and to receive according to what they have done in the body, whether good or evil.\footnote{2 Cor. v. 10; Eccles. xii. 14; Rom. ii. 16; xiv. 10, 12; Matt. xii. 36, 37.}

\emph{omnis potestas et judicium.}\footnote{John v. 22, 27.}\emph{ Quo quidem die non solum judicabuntur Angeli apostatici,}\footnote{1 Cor. vi. 3; Jude 6; 2 Pet. ii. 4.}\emph{ verum etiam omnes homines, quotquot uspiam in orbe terrarum aliquando vixerint, coram Christi tribunali comparebunt, ut cogitationum, dictorum, factorumque suorum rationem reddant, recipiantque simul juxta id quod in corpore quisque fecerit, seu bonum fuerit sive malum.}\footnote{2 Cor. v. 10; Eccles. xii. 14; Rom. ii. 16; xiv. 10, 12; Matt. xii. 36, 37.}

II. The end of God's appointing this day, is for the manifestation of the glory of his mercy in the eternal salvation of the elect;\footnote{Rom. ix. 23; Matt. xxv. 21.} and of his justice in the damnation of the reprobate, who are wicked and disobedient.\footnote{Rom. ii. 5, 6; 2 Thess. i. 7, 8; Rom. ix. 22.} For then shall the righteous go into everlasting life, and receive that fullness of joy and refreshing which shall come from the presence of the Lord:\footnote{Matt. xxv. 31–34; Acts iii. 19; 2 Thess. i. 7.} but the wicked, who know not God, and obey not the gospel of Jesus Christ, shall be cast into eternal torments, and be punished with everlasting destruction from the presence of the Lord, and from the glory of his power.\footnote{Matt. xxv. 41, 46; 2 Thess. i. 9; [Am. ed. Isa. lxvi. 24].}

II. \emph{Eo autem consilio Diem hum præstituit Deus, quo nempe misericordiæ suæ constaret gloria ex æterna salute electorum, justitiæ autem e damnatione reproborum, qui improbi sunt et contumaces. Tunc enim justi introibunt in vitam æternam, recipientque plenitudinem illam gaudii ac refrigerii, quæ a præsentia Domini ventura sunt: Impii autem, qui Deum ignorant, quique Evangelio Jesu Christi non morem gerunt, in æternos cruciatus detrudentur, æternaque perditione punientur a præsentia Domini et a potentiæ ipsius gloria profligati.}\footnote{Matt. xxv. 41, 46; 2 Thess. i. 9; [Am. ed. Isa. lxvi. 24].}

III. As Christ would have us to be certainly persuaded that there shall be a day of judgment, both to deter all men from sin, and for the

III. \emph{Quemadmodum Christus nobis, futurum esse aliquando diem judicii, esse velit persuasissimum; tum quo omnes a peccato absterreantur, }

greater consolation of the godly in their adversity:\footnote{2 Pet. iii. 11, 14; 2 Cor. v. 10, 11; 2 Thess. i. 5–7; Luke xxi. 27, 28; Rom. viii. 23–25.} so will he have that day unknown to men, that they may shake off all carnal security, and be always watchful, because they know not at what hour the Lord will come; and may be ever prepared to say, Come, Lord Jesus, come quickly.\footnote{Matt. xxiv. 36, 42–44; Mark xiii. 35–37; Luke xii. 35, 36; Rev. xxii. 20.} Amen.

\emph{tum ob majus piorum solatium in rebus adversis:}\footnote{2 Pet. iii. 11, 14; 2 Cor. v. 10, 11; 2 Thess. i. 5–7; Luke xxi. 27, 28; Rom. viii. 23–25.}\emph{ ita sane diem ipsum vult ab hominibus ignorari, quo securitatem omnem carnalem excutiant, et nunquam non sint vigilantes }(\emph{quum qua hora venturus sit Dominus ignorant})\emph{ utque semper sint parati ad dicendum Veni Domine Jesu, etiam cito veni.}\footnote{Matt. xxiv. 36, 42–44; Mark xiii. 35–37; Luke xii. 35, 36; Rev. xxii. 20.} Amen.

C\textsc{harles} H\textsc{erle}, \emph{Prolocutor}.

C\textsc{ornelius} B\textsc{urges}, \emph{Assessor}.

H\textsc{erbert} P\textsc{almer}, \emph{Assessor.}

H\textsc{enry} R\textsc{obroughe}, \emph{Scriba}.

A\textsc{doniram} B\textsc{yfield}, \emph{Scriba}.

\xchapter{The Westminster Shorter Catechism, A.D. 1647}

[This Catechism was prepared by the Westminster Assembly in 1647, and adopted by the General Assembly of the Church of Scotland, 1648; by the Presbyterian Synod of New York and Philadelphia, May, 1788: and by nearly all the Calvinistic Presbyterian and Congregational Churches of the English tongue. It was translated into Greek, Hebrew, Arabic, and many other languages, and appeared in innumerable editions. Although little known on the continent of Europe, it is more extensively used than any other Protestant catechism, except perhaps the Small Catechism of Luther and the Heidelberg Catechism. Want of space compels us to omit the Assembly's Larger Catechism, which is easy of access. For the same reason we have omitted the Scripture proofs.

The English original is conformed to the edition of the Presbyterian Board, compared with the London edition of 1658 and other older English and Scotch editions, which present no variations of any account. The Latin translation is from the Cambridge and Edinburgh editions, containing the Confession and both Catechisms, and reprinted in Niemeyer's Appendix.]

T\textsc{he} S\textsc{horter} C\textsc{atechism.}

C\textsc{atechismus} M\textsc{inor}.

Question. 1. What is the chief end of man?

Quæstio. \emph{Quis hominis finis est præcipuus?}

\emph{Answer.} Man's chief end is to glorify God, and to enjoy him forever.

Responsio. \emph{Præcipuus hominis finis est, Deum glorificare, eodemque frui in æternum.}

Ques. 2. What rule hath God given to direct us how we may glorify and enjoy him?

Quæs. \emph{Quam nobis regulam dedit Deus, qua nos ad ejus glorificationem ac fruitionem dirigamur?}

\emph{Ans.} The Word of God, which is contained in the Scriptures of the Old and New Testaments,\footnote{The London edition of 1658, Dunlop's Collection of 1719, and other editions read \emph{Testament.}} is the only rule to direct us how we may glorify and enjoy him.

Resp. \emph{Verbum Dei }(\emph{quod Scripturis Veteris ac Novi instrumenti comprehenditur}) \emph{est unica regula, qua nos ad ejus glorificationem ac fruitionem dirigamur.}

Ques. 3. What do the Scriptures principally teach?

Quæs. \emph{Quid est quod Scripturæ præcipue docent?}

\emph{Ans.} The Scriptures principally teach what man is to believe concerning God, and what duty God requires of man.

Resp. \emph{Duo imprimis sunt quæ Scripturæ docent, quid homini de Deo sit credendum, quidque officii exigat ab homine Deus.}

Ques.4. What is GOD?

Quæs. \emph{Quid est Deus?}

\emph{Ans.} God is a Spirit, infinite, eternal, and unchangeable, in his

Resp. \emph{Deus est Spiritus essentia, sapientia, potentia, sanctitate, justitia,}

being, wisdom, power, holiness, justice, goodness, and truth.

\emph{bonitate ac veritate infinitus, æternus, ac immutabilis.}

\emph{Ques.} 5. \emph{Are there more Gods than one?}

Quæs. \emph{Suntne plures uno Deo?}

\emph{Ans.} There is but one only, the living and true God.

Resp. \emph{Unus est unicusque, vivens ille verusque Deus.}

\emph{Ques.} 6. \emph{How many persons are there in the Godhead?}

Quæs. \emph{Quot sunt personæ in Deitate?}

\emph{Ans.} There are three persons in the Godhead: the Father, the Son, and the Holy Ghost; and these three are one God, the same in substance, equal in power and glory.

Resp. \emph{In Deitate personæ tres sunt, Pater, Filius, ac Spiritus Sanctus; suntque hæ tres personæ Deus unus, substantia eædem, potentia ac gloria coæquales.}

\emph{Ques.} 7. \emph{What are the decrees of God?}

Quæs. \emph{Quid sunt decreta Dei?}

\emph{Ans.} The decrees of God are his eternal purpose according to the counsel of his will, whereby, for his own glory, he hath fore-ordained whatsoever comes to pass.

Resp. \emph{Decreta Dei sunt æternum ejus propositum secundum voluntatis suæ consilium, quo quicquid unquam evenit, propter suam ipsius gloriam præordinavit.}

\emph{Ques.} 8. \emph{How doth God execute his decrees?}

Quæs. \emph{Quomodo decreta sua exequitur Deus?}

\emph{Ans.} God executeth his decrees in the works of creation and providence.

Resp. \emph{Deus exequitur decreta sua creationis operibus ac providentiæ.}

\emph{Ques.} 9. \emph{What is the work of creation?}

Quæs. \emph{Quid est opus creationis?}

\emph{Ans.} The work of creation is God's making all things of nothing, by the word of his power, in the space of six days, and all very good.

Resp. \emph{Opus creationis est quo Deus per verbum potentiæ suæ omnia sex dierum spatio ex nihilo condidit, atque omnia quidem valde bona.}

\emph{Ques.} 10. \emph{How did God create man?}

Quæs. \emph{Qualem creavit Deus hominem?}

\emph{Ans.} God created man, male and female, after his own image, in knowledge, righteousness, and holiness,

Resp. \emph{Deus hominem creavit marem ac fœminam, juxta suam ipsius imaginem, in cognitione, justitia ac }

with dominion over the creatures.

\emph{sanctitate, dominium habentem in creaturas.}

\emph{Ques.} 11. \emph{What are God's works of providence?}

Quæs. \emph{Quænam sunt opera Divinæ providentiæ?}

\emph{Ans.} God's works of providence are his most holy, wise, and powerful preserving and governing all his creatures, and all their actions.

Resp. \emph{Providentiæ Divinæ opera sunt sanctissima Dei, sapientissima potentissimaque creaturarum suarum omnium, earumque actionum conservatio et gubernatio.}

\emph{Ques.} 12. \emph{What special act of providence did God exercise towards man, in the estate wherein he was created?}

Quæs. \emph{Quem peculiarem providentiæ suæ actum exercebat Deus circa hominem in statu creationis suæ existentem?}

\emph{Ans.} When God had created man, he entered into a covenant of life with him, upon condition of perfect obedience; forbidding him to eat of the tree of knowledge of good and evil, upon pain of death.

Resp. \emph{Postquam Deus hominem condidisset, inibat cum illo fœdus vitæ, sub conditione perfectæ, obedientiæ; esu de arbore scientiæ boni malique sub pœna mortis eidem interdicens.}

\emph{Ques.} 13. \emph{Did our first parents continue in the estate wherein they were created?}

Quæs. \emph{An vero Primi nostri Parentes in quo creati fuerant statu perstitere?}

\emph{Ans.} Our first parents, being left to the freedom of their own will, fell from the estate wherein they were created, by sinning against God.

Resp. \emph{Primi Parentes voluntatis suæ libertati permissi peccando in Deum statu in quo creati fuerant exciderunt.}

\emph{Ques.} 14. \emph{What is sin?}

Quæs. \emph{Quid est peccatum?}

\emph{Ans.} Sin is any want of conformity unto, or transgression of, the law of God.

Resp. \emph{Peccatum est defectus quilibet conformitatis cum lege Divina, seu quævis ejusdem transgressio.}

\emph{Ques.} 15. What was the sin whereby our first parents fell from the estate wherein they were created?

Quæs. \emph{Quodnam erat peccatum istud quo primi parentes statu in quo creati fuerant exciderunt?}

\emph{Ans.} The sin whereby our first parents fell from the estate wherein

Resp. \emph{Peccatum istud quo primi parentes statu in quo creati fuerant}

they were created, was their eating the forbidden fruit.

\emph{exciderunt, erat comestio fructus interdicti.}

\emph{Ques.} 16. \emph{Did all mankind fall in Adam's first transgression?}

Quæs. \emph{Totumne genus humanum cecidit in prima Adami transgressione?}

\emph{Ans.} The covenant being made with Adam, not only for himself, but for his posterity, all mankind descending from him by ordinary generation, sinned in him, and fell with him, in his first transgression.

Resp. \emph{Quandoquidem fœdus cum Adamo ictum fuerat non suo tantum sed et posterorum suorum nomine; exinde factum est ut totum genus humanum ab illo generatione ordinaria procreatum, in eo peccaverit, cumque eo ceciderit, in prima ejus transgressione.}

\emph{Ques.} 17. \emph{Into what estate did the fall bring mankind?}

Quæs. \emph{In quem vero statum præcipitavit lapsus iste humanum genus?}

\emph{Ans.} The fall brought mankind into an estate of sin and misery.

Resp. \emph{Lapsus iste humanum genus in statum peccati ac miseriæ præcipitavit.}

\emph{Ques.} 18. \emph{Wherein consists the sinfulness of that estate whereinto man fell?}

Quæs. \emph{In quo consistit status illius in quem lapsus est homo peccaminositas?}

\emph{Ans.} The sinfulness of that estate whereinto man fell, consists in the guilt of Adam's first sin, the want of original righteousness, and the corruption of his whole nature, which is commonly called original sin; together with all actual transgressions which proceed from it.

Resp. \emph{Status in quem lapsus est homo peccaminositas consistit in reatu primi illius peccati quod Adamus admisit, in defectu originalis justitiæ, totiusque naturæ corruptione, quod Peccatum originale vulgo dicitur; una cum omnibus peccatis actualibus exinde profluentibus.}

\emph{Ques.} 19. \emph{What is the misery of that estate whereinto man fell?}

Quæs. \emph{Quæ miseria est illius status in quem homo lapsus est?}

\emph{Ans.} All mankind by their fall lost communion with God, are under his wrath and curse, and so made liable to all the\footnote{Older editions omit \emph{the}.} miseries in

Resp. \emph{Universum genus humanum lapsu suo communionem cum Deo perdidit, sub ira ejus et maledictione est constitutum, adeoque cunctis hujus }

this life, to death itself, and to the pains of hell forever.

\emph{vitæ miseriis, ipsi morti, infernique cruciatibus in æternum est obnoxium.}

\emph{Ques.} 20. \emph{Did God leave all mankind, to perish in the estate of sin and misery?}

Quæs. \emph{An vero Deus humanum genus universum in statu peccati ac miseriæ periturum dereliquit?}

\emph{Ans.} God, having out of his mere good pleasure, from all eternity, elected some to everlasting life, did enter into a covenant of grace, to deliver them out of the estate of sin and misery, and to bring them into an estate of salvation by a Redeemer.

Resp. \emph{Deus cum ex mero suo beneplacito nonnullos ad vitam æternam ab omni retro æternitate elegisset, fœdus gratiæ cum eis iniit; se nempe liberaturum eos e statu peccati ac miseriæ, atque in statum salutis per redemptorem translaturum.}

\emph{Ques.} 21. \emph{Who is the Redeemer of God's elect?}

Quæs. \emph{Quis est Redemptor electorum Dei?}

\emph{Ans.} The only Redeemer of God's elect is the Lord Jesus Christ, who being the eternal Son of God became man, and so was, and continueth to be, God and man, in two distinct natures, and one person forever.

Resp. \emph{Dominus Jesus Christus est electorum Dei Redemptor unicus, qui æternus Dei Filius cum esset, factus est homo; adeoque fuit, est, eritque } \textgreek{Θεάνθρωπος}, \emph{e naturis duabus distinctis persona unica in sempiternum.}

\emph{Ques.} 22. \emph{How did Christ, being the Son of God, become man?}

Quæs. \emph{Qui autem Christus, Filius Dei cum esset, factus est homo?}

\emph{Ans.} Christ, the Son of God, became man, by taking to himself a true body, and a reasonable soul, being conceived by the power of the Holy Ghost, in the womb of the Virgin Mary, and born of her, yet without sin.

Resp. \emph{Christus Filius Dei factus est homo, dum corpus verum, animamque rationalem assumeret sibi vi Spiritus Sancti in utero eque substantia Virginis Mariæ conceptus, et ex eadem natus, immunis tamen a peccato.}

\emph{Ques.} 23. \emph{What offices doth Christ execute as our Redeemer?}

Quæs. \emph{Quæ munera Christus ut Redemptor noster obit?}

\emph{Ans.} Christ, as our Redeemer, executeth the offices of a Prophet,

Resp. \emph{Christus quatenus Redemptor noster obit munera Prophetæ, }

of a Priest, and of a King, both in his estate of humiliation and exaltation.

\emph{Sacerdotis ac Regis, cum in humiliationis tum in exaltationis suæ statu.}

\emph{Ques.} 24. \emph{How doth Christ execute the office of a Prophet?}

Quæs. \emph{Quomodo Prophetæ munere defungitur Christus?}

\emph{Ans.} Christ executeth the office of a Prophet, in revealing to us by his Word and Spirit, the will of God for our salvation.

Resp. \emph{Christus defungitur Prophetæ munere, voluntatem Dei in salutem nostram nobis per verbum suum spiritumque revelando.}

\emph{Ques.} 25. \emph{How doth Christ execute the office of a Priest?}

Quæs. \emph{Qua ratione exequitur Christus munus Sacerdotale?}

\emph{Ans.} Christ executeth the office of a Priest, in his once offering up of himself a sacrifice to satisfy divine justice, and reconcile us to God, and in\footnote{Older editions omit in.} making continual intercession for us.

Resp. \emph{Christus exequitur Sacerdotale munus, semetipsum semel in sacrificium offerendo, quo justitiæ divinæ satisfaceret, nosque Deo conciliaret; prout etiam perpetuo pro nobis intercedendo.}

\emph{Ques.} 26. \emph{How doth Christ execute the office of a King?}

Quæs. \emph{Qui exequitur Christus munus Regium?}

\emph{Ans.} Christ executeth the office of a King, in subduing us to himself, in ruling and defending us, and in\footnote{Older editions omit in.} restraining and conquering all his and our enemies.

Resp. \emph{Christus exequitur munus Regium nos sibi subjugando, nos gubernando, tuendoque, ut etiam hostes suos nostrosque coërcendo ac debellando.}

\emph{Ques.} 27. Wherein did Christ's humiliation consist?

Quæs. \emph{In quo constitit Christi humiliatio?}

\emph{Ans.} Christ's humiliation consisted in his being born, and that in a low condition, made under the law, undergoing the miseries of this life, the wrath of God, and the cursed death of the cross; in being buried, and continuing under the power of death for a time.

Resp. \emph{Humiliatio Christi in eo constitit quod fuerit natus, et quidem humili conditione, factus sub lege, quodque vitæ hujus miserias, iram Dei mortemque crucis execrabilem subierit; quod sepultus fuerit, et sub potestate mortis aliquandiu commoratus.}

\emph{Ques.} 28. \emph{Wherein consisteth Christ's exaltation?}

Quæs. \emph{In quo consistit Christi exaltatio?}

\emph{Ans.} Christ's exaltation consisteth in his rising again from the dead on the third day, in ascending up into heaven, in sitting at the right hand of God the Father, and in coming to judge the world at the last day.

Resp. \emph{Exaltatio Christi consistit in resurrectione ejus a mortuis tertio die, ascensu in cœlum, sessione ad dextram Dei Patris, adventu ejus ad mundum judicandum die novissimo.}

\emph{Ques.} 29. \emph{How are we made partakers of the redemption purchased by Christ?}

Quæs. \emph{Qua ratione participes efficimur redemptionis per Christum acquisitæ?}

\emph{Ans.} We are made partakers of the redemption purchased by Christ, by the effectual application of it to us by his Holy Spirit.

Resp. \emph{Redemptionis per Christum acquisitæ participes efficimur ejusdem nobis efficaci per Spiritum ejus Sanctum, applicatione.}

\emph{Ques.} 30. \emph{How doth the Spirit apply to us the redemption purchased by Christ?}

Quæs. \emph{Quomodo nobis applicat Spiritus redemptionem per Christum acquisitam?}

\emph{Ans.} The Spirit applieth to us the redemption purchased by Christ, by working faith in us, and thereby uniting us to Christ in our effectual calling.

Resp. \emph{Spiritus nobis applicat redemptionem per Christum acquisitam fidem in nobis efficiendo, ac per eandem nos Christo in vocatione nostra efficaci uniendo.}

\emph{Ques.} 31. \emph{What is effectual calling?}

Quæs. \emph{Quid est vocatio efficax?}

\emph{Ans.} Effectual calling is the work of God's Spirit, whereby, convincing us of our sin and misery, enlightening our minds in the knowledge of Christ, and renewing our wills, he doth persuade and enable us to embrace Jesus Christ, freely offered to us in the gospel.

Resp. \emph{Vocatio efficax est Spiritus Dei opus, quo nos peccati ac miseriæ nostræ arguens, mentes nostras cognitione Christi illuminans, voluntates nostras renovans, prorsus nobis persuadet, et vires sufficit, ut Jesum Christum amplectamur, gratuito nobis oblatum in Evangelio.}

\emph{Ques.} 32. \emph{What benefits do they }

Quæs. \emph{Quænam beneficia in hac }

that are effectually called partake of in this life?

\emph{vita consequuntur ii qui sunt vocati efficaciter?}

\emph{Ans.} They that are effectually called do in this life partake of justification, adoption, sanctification, and the several benefits which, in this life, do either accompany or flow from them.

Resp. \emph{Qui vocati sunt efficaciter, justificationem, adoptionem, et sanctificationem in hac vita consequuntur, una cum omnibus iis beneficiis quæcunque solent in hac vita comitari illas, aut ab iisdem promanare. }

\emph{Ques.} 33. \emph{What is justification?}

Quæs. \emph{Quid est justificatio?}

\emph{Ans.} Justification is an act of God's free grace, wherein he pardoneth all our sins, and accepteth us as righteous in his sight, only for the righteousness of Christ imputed to us, and received by faith alone.

Resp. \emph{Justificatio est actus gratiæ Dei gratuitæ, quo peccata nobis condonat omnia, nosque tanquam justos in conspectu suo acceptat, propter solam Christi justitiam nobis imputatam, per fidem tantum apprehensam.}

\emph{Ques.} 34. \emph{What is adoption?}

Quæs.\emph{ Quid est adoptio?}

\emph{Ans.} Adoption is an act of God's free grace, whereby we are received into the number, and have a right to all the privileges, of the sons of God.

Resp. \emph{Adoptio est actus gratiæ Dei gratuitæ, quo in numerum recipimur ac jus obtinemus ad omnia privilegia filiorum Dei.}

\emph{Ques.} 35. \emph{What is sanctification?}

Quæs. \emph{Quid est sanctificatio?}

\emph{Ans.} Sanctification is the work of God's free grace, whereby we are renewed in the whole man after the image of God, and are enabled more and more to die unto sin, and live unto righteousness.

Resp. \emph{Sanctificatio est opus gratiæ Dei gratuitæ, quo in toto homine secundum imaginem Dei renovamur, et potentes efficimur, qui magis in dies magisque peccato quidem moriamur, justitiæ autem vivamus.}

\emph{Ques.} 36. \emph{What are the benefits which in this life do accompany or flow from justification, adoption, and sanctification?}

Quæs. \emph{Quænam sunt illa beneficia quæ justificationem, adoptionem et sanctificationem in hac vita vel comitantur, vel ab eis promanant?}

\emph{Ans.} The benefits which in this life do accompany or flow from justification, adoption, and sanctification, are, assurance of God's love,

Resp. \emph{Quæ justificationem, adoptionem et sanctificationem in hac vita vel comitantur vel ab eis promanant beneficia, sunt certitudo amoris Dei, }

peace of conscience, joy in the Holy Ghost, increase of grace, and perseverance therein to the end.

\emph{pax conscientiæ, gaudium in Spiritu Sancto, gratiæ incrementum, in eaque ad finem usque perseverantia.}

\emph{Ques.} 37. \emph{What benefits do believers receive from Christ at death?}

Quæs. \emph{Quænam a Christo beneficia in morte percipiunt fideles?}

\emph{Ans.} The souls of believers are, at their death, made perfect in holiness, and do immediately pass into glory; and their bodies, being still united to Christ, do rest in their graves till the resurrection.

Resp. \emph{Animæ fidelium in morte fiunt perfecte sanctæ, ac protinus in gloriam transferuntur; corpora vero usque Christo unita in sepulchris ad resurrectionem usque quiescunt.}

\emph{Ques.} 38. \emph{What benefits do believers receive from Christ at the resurrection?}

Quæs. \emph{Quæ tandem beneficia a Christo percipiunt fideles in resurrectione?}

\emph{Ans.} At the resurrection, believers being raised up in glory, shall be openly acknowledged and acquitted in the day of judgment, and made perfectly blessed in the full enjoying of God to all eternity.

Resp. \emph{In resurrectione fideles suscitati in gloria, palam agnoscentur et absolventur in die judicii, fientque perfecte beati plena Dei in omne æternum fruitione.}

\emph{Ques.} 39. \emph{What is the duty which God requireth of man?}

Quæs. \emph{Quid autem officii ac observantiæ ab homine exposcit Deus?}

\emph{Ans.} The duty which God requireth of man is obedience to his revealed will.

Resp. \emph{Officium quod ab homine Deus exposcit, est obedientia voluntati ejus revelatæ exhibenda.}

\emph{Ques.} 40. \emph{What did God at first reveal to man for the rule of his obedience?}

Quæs. \emph{Quid homini primum revelavit Deus, quod foret ipsi obedientiæ regula?}

\emph{Ans.} The rule which God at first revealed to man, for his obedience, was the moral law.

Resp. \emph{Obedientiæ regula, quam Deus homini primum revelavit, erat Lex moralis.}

\emph{Ques.} 41. \emph{Wherein is the moral law summarily comprehended?}

Quæs. \emph{Ubinam summatim comprehenditur lex moralis?}

\emph{Ans.} The moral law is summarily comprehended in the ten commandments.

Resp. \emph{Lex moralis summatim comprehenditur in Decalogo.}

\emph{Ques.} 42. \emph{What is the sum of the ten commandments?}

Quæs. \emph{Dic quænam sit Decalogi summa?}

\emph{Ans.} The sum of the ten commandments is, to love the Lord our God with all our heart, with all our soul, with all our strength, and with all our mind; and our neighbor as ourselves.

Resp. \emph{Summa Decalogi est ut Dominum nostrum toto corde, tota anima, tota mente, totisque viribus nostris diligamus; proximum vero nostrum sicut nosmetipsos.}

\emph{Ques.} 43. \emph{What is the preface to the ten commandments?}

Quæs. \emph{Quænam est Decalogi præfatio?}

\emph{Ans.} The preface to the ten commandments is in these words: \emph{I am the Lord thy God, which brought thee out of the land of Egypt, out of the house of bondage.}

Resp. \emph{Decalogi præfatio hisce verbis continetur} [Ego sum Dominus Deus tuus, qui te eduxi e terra Ægypti, e domo servitutis].

\emph{Ques.} 44. \emph{What doth the preface to the ten commandments teach us?}

Quæs. \emph{Quid nos edocet Decalogi præfatio?}

\emph{Ans.} The preface to the ten commandments teacheth us, that because God is the Lord, and our God and Redeemer, therefore we are bound to keep all his commandments.

Resp. \emph{Decalogi præfatio nos docet, quod quoniam Deus est Dominus, nosterque Deus ac redemptor, ea propter præcepta ejus omnia tenemur observare.}

\emph{Ques.} 45. \emph{Which is the first commandment?}

Quæs. \emph{Quodnam est mandatum primum?}

\emph{Ans.} The first commandment is, \emph{Thou shalt have no other gods before me.}

Resp. \emph{Mandatum primum est} [Non habebis Deos alios coram me].

\emph{Ques.} 46. \emph{What is required in the first commandment?}

Quæs. \emph{In mandato primo quid exigitur?}

\emph{Ans.} The first commandment requireth us to know and acknowledge God, to be the only true God, and our God; and to worship and glorify him accordingly.

Resp. \emph{In mandato primo exigitur ut Jehovam esse unicum illum verumque Deum, Deumque nostrum cognoscamus simul et agnoscamus, atque ut talem colamus, ac glorificemus.}

\emph{Ques.} 47. \emph{What is forbidden in the first commandment?}

Quæs. \emph{Quid est quod prohibetur mandato primo?}

\emph{Ans.} The first commandment forbiddeth the denying, or not worshiping and glorifying the true God, as God, and our God; and the giving that worship and glory to any other which is due to him alone.

Resp. \emph{In primo mandato prohibetur veri Dei abnegatio, neglectusque ipsum tanquam Deum, Deumque nostrum colendi ac glorificandi; prout etiam cultum ac gloriam illi soli debita alii cuivis tribuere aut exhibere.}

\emph{Ques.} 48. \emph{What are we specially taught by these words, }''before me,'' \emph{in the first commandment?}

Quæs. \emph{Quid imprimis docemur verbis istis mandati primi} [Coram me]?

\emph{Ans.} These words, ``\emph{before me}'' in the first commandment, teach us that God, who seeth all things, taketh notice of, and is much displeased with, the sin of having any other God.

Resp. \emph{Verba isthæc} [Coram me] \emph{in mandato primo nos docent, Deum qui omnia intuetur, peccatum alium habendi Deum cum imprimis advertere, tum vero eodem offendi plurimum.}

\emph{Ques.} 49. \emph{Which is the second commandment?}

Quæs. \emph{Quodnam est præceptum secundum?}

\emph{Ans.} The second commandment is, \emph{Thou shalt not make unto thee any graven image, or any likeness of any thing that is in heaven above, or that is in the earth beneath, or that is in the water under the earth; thou shalt not bow down thyself to them, nor serve them; for I the Lord thy God am a jealous God, visiting the iniquity of the fathers upon the children, unto the third and fourth generation of them that hate me, and showing mercy unto thousands of them that love me and Jceep my commandments.}

Resp. \emph{Secundum præceptum est} [Non facies tibi imaginem quamvis sculptilem, aut similitudinem rei cujusvis quæ est in cœlis superne, aut inferius in terris, aut in aquis infra terram; non incurvabis te iis, nec eis servies: siquidem ego Dominus Deus tuus Deus sum Zelotypus, visitans iniquitates patrum in filios ad tertiam usque quartamque progeniem osorum mei, exhibens vero misericordiam ad millenas usque diligentium me, ac mandata mea observantium].

\emph{Ques.} 50. \emph{What is required in the second commandment?}

Quæs. \emph{Quid exigitur in secundo præcepto?}

\emph{Ans.} The second commandment requireth the receiving, observing, and keeping pure and entire, all such religious worship and ordinances as God hath appointed in his Word.

Resp. \emph{Præceptum secundum exigit, ut cultus omnes ac instituta religionis quæcunque Deus in verbo suo constituit, excipiamus, observemus, pura denique ac integra custodiamus.}

\emph{Ques.} 51. \emph{What is forbidden in the second commandment?}

Quæs. \emph{Quid est quod in secundo præcepto prohibetur?}

\emph{Ans.} The second commandment forbiddeth the worshiping of God by images, or any other way not appointed in his Word.

Resp. \emph{Secundum præceptum interdicit nobis cultu Dei per simulacra, aut alia ratione quaviscunque quam in verbo suo Deus non præscripsit.}

\emph{Ques.} 52. \emph{What are the reasons annexed to the second commandment?}

Quæs. \emph{Quænam sunt rationes præcepto secundo annexe?}

\emph{Ans.} The reasons annexed to the second commandment are, God's sovereignty over us, his propriety\footnote{London ed. of 1658 reads \emph{property}, and \emph{his} zeal.} in us, and the\footnote{London ed. of 1658 reads \emph{property}, and \emph{his} zeal.} zeal he hath to his own worship.

Resp. \emph{Rationes secundo præcepto annexæ sunt, supremum Dei in nos dominium, illius jus in nobis peculiare, zelusque quo suum ipsius cultum prosequitur.}

\emph{Ques.} 53. \emph{Which is the third commandment?}

Quæs. \emph{Age quodnam est tertium mandatum?}

\emph{Ans.} The third commandment is, \emph{Thou shalt not take the name of the Lord thy God in vain: for the Lord will not hold him guiltless that taketh his name in vain.}

Resp. \emph{Mandatum tertium sic habetur} [Nomen Domini Dei tui inaniter non usurpabis; non enim eum pro insonte habebit Dominus qui nomen ejus inaniter adhibuerit].

\emph{Ques.} 54. \emph{What is required in the third commandment?}

Quæs. \emph{Quid exigitur in mandate tertio?}

\emph{Ans.} The third commandment requireth the holy and reverent use of God's names, titles, attributes, ordinances, word, and works.

Resp. \emph{Mandatum tertium exigit ut Dei nomina, titulos, attributa, instituta, verba, operaque sancte summaque cum reverentia adhibeamus.}

\emph{Ques.} 55. \emph{What is forbidden in the third commandment?}

Quæs. \emph{Quid prohibetur mandato tertio?}

\emph{Ans.} The third commandment forbiddeth all profaning or abusing of any thing whereby God maketh himself known.

Resp. \emph{Mandatum tertium prohibet rei cujusvis qua Deus se notum facit, profanationem omnem ac abusum.}

\emph{Ques.} 56. \emph{What is the reason annexed to the third commandment?}

Quæs. \emph{Quænam est ratio subnexa mandato tertio?}

\emph{Ans.} The reason annexed to the third commandment is, that however the breakers of this commandment may escape punishment from men, yet the Lord our God will not suffer them to escape his righteous judgment.

Resp. \emph{Ratio mandato tertio subnexa est, quod licet hujus præcepti violatores ab hominibus quandoque nil supplicii ferant, nihilominus tamen Dominus Deus noster eos justum ejus judicium neutiquam patietur subterfugere.}

\emph{Ques.} 57. \emph{Which is the fourth commandment?}

Quæs. \emph{Recita mandatum quartum?}

\emph{Ans.} The fourth commandment is, \emph{Remember the Sabbath-day to keep it holy. Six days shalt thou labor, and do all thy work: but the seventh day is the Sabbath of the Lord thy God: in it thou shalt not do any work, thou, nor thy son, nor thy daughter, thy man-servant, nor thy maid-servant, nor thy cattle, nor thy stranger that is within thy gates; for in six days the Lord made heaven and earth, the sea, and all that in them is, and rested the seventh day: wherefore the Lord blessed the Sabbath-day and hallowed it.}

Resp. \emph{Mandati quarti verba sunt isthæc} [Memineris diem Sabbati ut sanctifices eum; sex diebus operaberis et facies omne opus tuum, septimus vero dies sabbatum est Domini Dei tui, opus in eo nullum facies tu, neque filius tuus, neque filia tua, nec servus tuus, nec ancilla tua, neque jumentum tuum, nec hospes tuus quicunque intra portas tuas commoratur: Nam sex diebus perfecit Dominus cœlum terramque, mare, et quicquid in illis continetur, septimo vero die requievit, quamobrem benedixit Dominus diei sabbati, eumque sanctificavit.]

\emph{Ques.} 58. \emph{What is required in the fourth commandment?}

Quæs. \emph{Quid a nobis exigit mandatum quartum?}

\emph{Ans.} The fourth commandment

Resp. \emph{Quartum mandatum a nobis}

requireth the keeping holy to God such set times as he hath appointed in his Word; expressly one whole day in seven, to be a holy Sabbath to himself.\footnote{1 London ed. of 1658: \emph{unto the Lord.}}

\emph{exigit, ut statum illud tempus quod in verbo suo designavit Deus, sanctum ei observemus; integrum nempe Diem e septenis unum in sanctum illi sabbatum celebrandum.}

\emph{Ques.} 59. \emph{Which day of the seven hath God appointed to be the weekly Sabbath?}

Quæs. \emph{E septenis autem quem diem sabbato hebdomadario designavit Deus?}

\emph{Ans.} From the beginning of the world to the resurrection of Christ, God\footnote{London ed. of 1658 inserts \emph{hath}.} appointed the seventh day of the week to be the weekly Sabbath; and the first day of the week, ever since, to continue to the end of the world, which is the Christian Sabbath.

Resp. \emph{Deus hebdomadario sabbato designavit septimum diem hebdomadæ ah initio mundi usque ad Christi resurrectionem, exinde vero ad finem usque mundi duraturum, diem septimanæ primum, quod est sabbatum Christianum.}

\emph{Ques.} 60. \emph{How is the Sabbath to be sanctified?}

Quæs. \emph{Qui autum est sabbatum sanctificandum?}

\emph{Ans.} The Sabbath is to be sanctified by a holy resting all that day, even from such worldly employments and recreations as are lawful on other days; and spending the whole time in the public and private exercises of God's worship, except so much as is to be taken up in the works of necessity and mercy.

Resp. \emph{Sabbatum est sanctificandum diem illum integrum sancte quiescendo, etiam a negotiis et recreationibus mundanis, aliis quidem diebus haud illicitis; totumque illud temporis} (\emph{præterquam quod operibus necessitatis ac misericordiæ insumendum fuerit}) \emph{cultus Divini exercitiis publicis privatisque impendendo.}

\emph{Ques.} 61. \emph{What is}\footnote{London ed. of 1658: \emph{what are the sins}.} \emph{forbidden in the fourth commandment?}

Quæs. \emph{Quid prohibetur in mandato quarto?}

\emph{Ans.} The fourth commandment forbiddeth the omission,\footnote{London ed. of 1658: \emph{the omission of careful}.} or careless performance, of the duties required, and the profaning the day

Resp. \emph{Mandatum quartum prohibet officiorum quæ inibi requiruntur, cum omissionem tum præstationem negligentem; prout etiam ejus diei }

by idleness, or doing that which is in itself sinful, or by unnecessary thoughts, words, or works about our worldly employments and\footnote{London ed. of 1658 reads, \emph{or}.} recreations.

\emph{profanationem qualemcunque, sive illum otiose consumendo, sive quod in se peccatum est faciendo, seu denique circa mundana negotia vel recreationes cogitationibus, dictis, factisve non necessariis.}

\emph{Ques.} 62. \emph{What are the reasons annexed to the fourth commandment?}

Quæs. \emph{Quænam sunt quarto præcepto rationes annexe?}

\emph{Ans.} The reasons annexed to the fourth commandment are, God's allowing us six days of the week for our own employments,\footnote{London ed. of 1658: \emph{employment}.} his challenging a special propriety\footnote{London ed. of 1658: \emph{property}.} in the seventh, his own example, and his blessing the Sabbath-day.

Resp. \emph{Rationes quarto præcepto annexæ sunt istiusmodi; quoniam e septimana qualibet sex dies concesserit nobis Deus nostris ipsorum negotiis insumendos; quoniam in septimo jus sibi vendicat peculiare; quoniam Deus exemplo suo nobis præivit, ac diei sabbati benedixit.}

\emph{Ques.} 63. \emph{Which}\footnote{London ed. of 1658: \emph{what}.} \emph{is the fifth commandment?}

Quæs. \emph{Quodnam est præceptum quintum?}

\emph{Ans.} The fifth commandment is, \emph{Honor thy father and thy mother: that thy days may be long upon the land which the Lord thy God giveth thee.}

Resp. \emph{Quintum præceptum est hujusmodi} [Honora patrem tuum ac matrem tuam ut prolongentur dies tui in terra illa quam tibi largitur Dominus Deus tuus].

\emph{Ques.} 64. \emph{What is required in the fifth commandment?}

Quæs. \emph{Quid est quod jubemur mandato quinto?}

\emph{Ans.} The fifth commandment requireth the preserving the honor of,\footnote{London ed. of 1658 omits \emph{of}.} and performing the duties belonging to, every one in their several places and relations, as superiors, inferiors, or equals.

Resp. \emph{Mandatum quintum nos jubet honorem conservare, ac officia persolvere unicuique pro ratione ordinis ac relationis in quibus fuerit exhibenda, seu superior nobis fiet, sive inferior, sive denique æqualis.}

\emph{Ques.} 65. \emph{What is forbidden in the fifth commandment?}

Quæs. \emph{Quid est quod mandatum quintum vetat?}

\emph{Ans.} The fifth commandment forbiddeth the neglecting of, or doing any thing against, the honor and duty which belongeth to every one in their several places and relations.

Resp. \emph{Quintum mandatum vetat honorem, officiumque singulis debitum pro ratione ordinis ac relatione in quibus fuerint, aut negligere, aut adversus ea quicquam machinari.}

\emph{Ques.} 66. \emph{What is the reason annexed to the fifth commandment?}

Quæs. \emph{Quæ ratio subnectitur quinto præcepto?}

\emph{Ans.} The reason annexed to the fifth commandment is, a promise of long life and prosperity (as far as it shall serve for God's glory, and their own good) to all such as keep this commandment.

Resp. \emph{Ratio quinto præcepto subnexa est promissio longævitatis, prosperitatisque} (\emph{quatenus nempe Dei gloriæ ipsorumque conducant utilitati}) \emph{omnibus facta hoc præceptum observantibus.}

\emph{Ques.} 67. \emph{Which is the sixth commandment?}

Quæs. \emph{Cedo mandatum sextum?}

\emph{Ans.} The sixth commandment is, \emph{Thou shalt not kill.}

Resp. \emph{Mandatum sextum hisce verbis comprehenditur} [Non occides].

\emph{Ques.} 68. \emph{What is required in the sixth commandment?}

Quæs. \emph{Quid a nobis exigit mandatum sextum?}

\emph{Ans.} The sixth commandment requireth all lawful endeavors to preserve our own life, and the life of others.

Resp. \emph{Exigit a nobis mandatum sextum, ut vitam cum nostram tum aliorum honestis quibuscunque rationibus tueamur.}

\emph{Ques.} 69. \emph{What is forbidden in the sixth commandment?}

Quæs. \emph{Quid vero prohibet sextum mandatum?}

\emph{Ans.} The sixth commandment forbiddeth the taking away of our own life, or the life of our neighbor unjustly, or whatsoever tendeth thereunto.

Resp. \emph{Sextum mandatum prohibet vitam nobismetipsis, aut injuste proximo vitam adimere, aut quidvis quod eo tendat agere.}

\emph{Ques.} 70. \emph{Which is the seventh commandment?}

Quæs. \emph{Quodnam est mandatum septimum?}

\emph{Ans.} The seventh commandment

Resp. \emph{Mandatum septimum hæc }

is, \emph{Thou shalt not commit adultery.}

\emph{verba complectuntur} [Non mœchaberis].

\emph{Ques.} 71. \emph{What is required in the seventh commandment?}

Quæs. \emph{Quid exigitur mandato septimo?}

\emph{Ans.} The seventh commandment requireth the preservation of our own and our neighbor's chastity, in heart, speech, and behavior.

Resp. \emph{Mandatum septimum exigit ut tam nostram quam proximorum castitatem animo, sermone, gestuque conservemus.}

\emph{Ques.} 72. \emph{What is forbidden in the seventh commandment?}

Quæs. \emph{Quid prohibetur mandato septimo?}

\emph{Ans.} The seventh commandment forbiddeth all unchaste thoughts, words, and actions.

Resp. \emph{Septimum mandatum prohibet cogitationes, sermones, actionesque omnes impudicas.}

\emph{Ques.} 73. \emph{Which is the eighth commandment?}

Quæs. \emph{Quodnam est præceptum octavum?}

\emph{Ans.} The eighth commandment is, \emph{Thou shalt not steal.}

Resp. \emph{Præceptum octavum hoc est} [Non furaberis].

\emph{Ques.} 74 \emph{What is required in the eighth commandment?}

Quæs. \emph{Mandatum octavum quid a nobis exigit?}

\emph{Ans.} The eighth commandment requireth the lawful procuring and furthering the wealth and outward estate of ourselves and others.

Resp. \emph{Octavum mandatum a nobis exigit, facultates ac rem externam nostri aliorumque ut procuremus ac promoveamus.}

\emph{Ques.} 75. \emph{What is forbidden in the eighth commandment?}

Quæs. \emph{In octavo præcepto quid prohibetur?}

\emph{Ans.} The eighth commandment forbiddeth whatsoever doth, or may, unjustly hinder our own, or our neighbor's wealth or outward estate.

Resp. \emph{Octavum mandatum prohibet quicquid nostris aut proximorum nostrorum opibus rebusque externis injusto aut est aut esse possit impedimento.}

\emph{Ques.} 76. \emph{Which}\footnote{London ed. of 1658: \emph{what}.} \emph{is the ninth commandment?}

Quæs. \emph{Quodnam est præceptum nonum?}

\emph{Ans.} The ninth commandment is, \emph{Thou shalt not hear false witness against thy neighbor?}

Resp. \emph{Præceptum nonum sic se habet} [Non eris adversus proximum tuum testis mendax].

\emph{Ques.} 77. \emph{What is required in the ninth commandment?}

Quæs. \emph{Quid a nobis exigit præceptum nonum?}

\emph{Ans.} The ninth commandment requireth the maintaining and promoting of truth between man and man, and of our own and our neighbor's good name, especially in witness-bearing.

Resp. \emph{Præceptum nonum id a nobis exigit ut veritatem inter homines mutuo, utque bonum nomen et existimationem cum nostri tum proximorum nostrorum conservemus ac promoveamus, cum primis vero in ferendo testimonio.}

\emph{Ques.} 78. \emph{What is forbidden in the ninth commandment?}

Quæs. \emph{Quid prohibetur nono præcepto?}

\emph{Ans.} The ninth commandment forbiddeth whatsoever is prejudicial to truth, or injurious to our own or our neighbor's good name.

Resp. \emph{Nonum præceptum prohibet quicquid est aut veritati inimicum; aut existimationi nostri vel proximorum nostrorum injurium.}

\emph{Ques.} 79. \emph{Which}\footnote{London ed. of 1658: \emph{what}.} \emph{is the tenth commandment?}

Quæs. \emph{Quale est mandatum decimum?}

\emph{Ans.} The tenth commandment is, Thou shalt not covet thy neighbor's house, thou shalt not covet thy neighbor's wife, nor his man-servant, nor his maid-servant, nor his ox, nor his ass, nor any thing that is thy neighbor's.

Resp. \emph{Mandatum decimum hæc verba exhibent} [Non concupisces proximi tui domum, non concupisces proximi tui uxorem, non servum, non ancillam, non bovem, non asinum, neque aliud denique quicquam quod est proximi tui].

\emph{Ques.} 80. \emph{What is required in the tenth commandment?}

Quæs. \emph{In decimo præcepto quid exigitur?}

\emph{Ans.} The tenth commandment requireth full contentment with our own condition, with a right and charitable frame of spirit toward our neighbor, and all that is his.

Resp. \emph{Præceptum decimum exigit ut sorti nostræ plane acquiescamus, utque in proximum et quæcunque sunt ejus debite, benevoleque afficiamur.}

\emph{Ques.} 81. \emph{What is forbidden in the tenth commandment?}

Quæs. \emph{Quæ prohibentur decimo mandato?}

\emph{Ans.} The tenth commandment forbiddeth all discontentment with

Resp. \emph{Mandatum decimum prohibet rerum nostrarum displicentiam, }

our own estate, envying or grieving at the good of our neighbor, and all inordinate motions or\footnote{London ed. of 1658: \emph{and}.} affections to any thing that is his.

\emph{invidiam ac dolorem de bono proximi, una cum animi nostri motibus et affectionibus circa ea quæ proximi sunt inordinatis quibuscunque.}

\emph{Ques.} 82. \emph{Is any man able perfectly to keep the commandments of God?}

Quæs. \emph{Quisquamne potis est mandata Dei perfecte observare?}

\emph{Ans.} No mere man, since the fall, is able, in this life, perfectly to keep the commandments of God; but doth daily break them, in thought, word, and deed.

Resp. \emph{Post lapsum nemo extat humana tantum natura constans, qui mandata Dei perfecte in hac vita implere potest, quominus ea quotidie tum cogitatione, tum dictis factisque violet.}

\emph{Ques.} 83. \emph{Are all transgressions of the law equally heinous?}

Quæs. \emph{An vero sunt omnes violationes legis ex æquo graves?}

\emph{Ans.} Some sins in themselves, and by reason of several aggravations, are more heinous in the sight of God than others.

Resp. \emph{Peccata sunt nonnulla aliis cum sua natura, tum propter varias eorum aggravationes in conspectu Dei graviora.}

\emph{Ques.} 84. \emph{What doth every sin deserve?}

Quæs. \emph{Quid est quod meretur peccatum unumquodque?}

\emph{Ans.} Every sin deserveth God's wrath and curse, both in this life and that which is to come.

Resp. \emph{Unumquodque peccatum iram Dei meretur ac maledictionem cum in vita præsenti, tum in futura.}

\emph{Ques.} 85. \emph{What doth God require of us, that we may escape his wrath and curse, due to us for sin?}

Quæs. \emph{Quid autem exigit a nobis Deus, quo nobis ob peccatum debitas iram ejus ac maledictionem effugiamus?}

\emph{Ans.} To escape the wrath and curse of God, due to us for sin, God requireth of us faith in Jesus Christ, repentance unto life, with the diligent use of all the outward means whereby Christ communicateth to us the benefits of redemption.

Resp. \emph{Quo iram Dei ac maledictionem ob peccatum nobis debitas effugiamus, exigit a nobis Deus fidem in Jesum Christum, resipiscentiam ad vitam, una cum usu mediorum omnium externorum diligenti, quibus Christus nobis communicat redemptionis suæ beneficia.}

\emph{Ques.} 86. \emph{What is faith in Jesus Christ?}

Quæs. \emph{Quid est fides in Jesum Christum?}

\emph{Ans.} Faith in Jesus Christ is a saving grace, whereby we receive and rest upon him alone for salvation, as he is offered to us in the gospel.

Resp. \emph{Fides in Jesum Christum est gratia salvifica, qua illum recipimus, eoque solo nitimur, ut salvi simus, prout ille nobis offertur in evangelio.}

\emph{Ques.} 87. \emph{What is repentance unto life?}

Quæs. \emph{Quid est resipiscentia ad vitam?}

\emph{Ans.} Repentance unto life is a saving grace, whereby a sinner, out of a\footnote{London ed. of 1658 omits \emph{a}.} true sense of his sin, and apprehension of the mercy of God in Christ, doth, with grief and hatred of his sin, turn from it unto God, with full purpose of, and endeavor after, new obedience.

Resp. \emph{Resipiscentia ad vitam est gratia salvifica, qua peccator e vero peccati sui sensu, ac apprehensione divinæ in Christo misericordiæ, dolens ac perosus peccatum suum ab illo ad Deum convertitur, cum novae, obedientiæ pleno proposito et conatu.}

\emph{Ques.} 88. \emph{What are the outward and ordinary means whereby Christ communicateth to us the benefits of redemption?}

Quæs. \emph{Quænam sunt externa media quibus Christus nobis communicat redemptionis suæ beneficia?}

\emph{Ans.} The outward and ordinary means whereby Christ communicateth to us the benefits of redemption, are his ordinances, especially the word, sacraments, and prayer; all which are made effectual to the elect for salvation.

Resp. \emph{Media externa ac ordinaria quibus Christus nobis communicat redemptionis suæ beneficia sunt ejus instituta, verbum præsertim, sacramenta, et oratio; quæ quidem omnia electis redduntur efficacia ad salutem.}

\emph{Ques.} 89. \emph{How is the word made effectual to salvation?}

Quæs. \emph{Qua ratione fit verbum efficax ad salutem?}

\emph{Ans.} The Spirit of God maketh the reading, but especially the preaching of the word, an effectual means of convincing and converting

Resp. \emph{Spiritus Dei lectionem verbi præcipue vero prædicationem ejus reddit medium efficax convincendi, convertendique peccatores, eosdemque }

sinners, and of building them up in holiness and comfort through faith unto salvation.

\emph{in sanctimonia et consolatione ædificandi per fidem ad salutem.}

\emph{Ques.} 90. \emph{How is the Word to be read and heard, that it may become effectual to salvation?}

Quæs. \emph{Quomodo legi debet ac audiri verbum, ut evadat efficax ad salutem?}

\emph{Ans.} That the Word may become effectual to salvation, we must attend thereunto with diligence, preparation, and prayer; receive it with faith and love, lay it up in our hearts, and practice it in our lives.

Resp. \emph{Quo verbum evadat efficax ad salutem, debemus ei cum præparatione, ac oratione diligenter attendere; idemque fide excipere ac amore, in animis nostris recondere, ac in vita nostra exprimere.}

\emph{Ques.} 91. \emph{How do the sacraments become effectual means of salvation?}

Quæs. \emph{Qui evadunt sacramenta media efficacia ad salutem?}

\emph{Ans.} The sacraments become effectual means of salvation, not from any virtue in them, or in him that doth administer them, but only by the blessing of Christ, and the working of his Spirit in them that by faith receive them.

Resp. \emph{Sacramenta evadunt efficacia ad salutem media, non ulla in ipsis vi, nec in eo qui illa administrat; verum Christi solummodo benedictione, ac Spiritus ejus in iis qui illa per fidem recipiunt operatione.}

\emph{Ques.} 92. \emph{What is a sacrament?}

Quæs. \emph{Quid est sacramentum?}

\emph{Ans.} A sacrament is a holy ordinance instituted by Christ; wherein, by sensible signs, Christ and the benefits of the new covenant are represented, sealed, and applied to believers.

Resp. \emph{Sacramentum est ordinatio sacra a Christo instituta, in qua fidelibus per signa in sensus incurrentia Christus novique fœderis beneficia repræsentantur, obsignantur, et applicantur.}

\emph{Ques.} 93. \emph{Which are the sacraments of the New Testament?}

Quæs. \emph{Quænam sunt sacramenta Novi Testamenti?}

\emph{Ans.} The sacraments of the New Testament are Baptism and the Lord's Supper.

Resp. \emph{Sacramenta Novi Testamenti sunt Baptismus ac cœna Dei.}

\emph{Ques.} 94. \emph{What is Baptism?}

Quæs. \emph{Quid est baptismus?}

\emph{Ans.} Baptism is a sacrament,

Resp. \emph{Baptismus est Sacramentum, }

wherein the washing with water, in the name of the Father, and of the Son, and of the Holy Ghost, doth signify and seal our ingrafting into Christ and partaking of the benefits of the covenant of grace, and our engagement to be the Lord's.

\emph{in quo ablutio per aquam in nomine Patris ac Filii ac Spiritus Sancti, nostram in Christum insitionem, et beneficiorum fœderis gratiæ, participationem, pactumque nostrum, nos nempe Domini futuros esse totos, significat obsignatque.}

\emph{Ques.} 95. \emph{To whom is Baptism to be administered?}

Quæs. \emph{Quibus est Baptismus administrandus?}

\emph{Ans.} Baptism is not to be administered to any that are out of the visible Church, till they profess their faith in Christ, and obedience to him; but the infants of such as are members of the visible church, are to be baptized.

Resp. \emph{Baptismus non est administrandus quibusdam extra Ecclesiam visibilem constitutis, donec se in Christum credere, eique obedientes fore professi fuerint; verum infantes eorum qui membra sunt Ecclesiæ visibilis sunt baptizandi.}

\emph{Ques.} 96. \emph{What is the Lord's Supper?}

Quæs. \emph{Quid est cœna Domini?}

\emph{Ans.} The Lord's Supper is a sacrament, wherein, by giving and receiving bread and wine, according to Christ's appointment, his death is showed forth, and the worthy receivers are, not after a corporal and carnal manner, but by faith, made partakers of his body and blood, with all his benefits, to their spiritual nourishment and growth in grace.

Resp. \emph{Cœna Domini est Sacramentum, in quo pane ac vino secundum Christi institutum datis acceptisque, mors ejus ostenditur; quæ qui digne participant, corporis ejus et sanguinis }(\emph{non quidem corporeo et carnali modo, verum}) \emph{per fidem fiunt participes, omniumque ipsius beneficiorum ad nutritionem suam spiritualem suumque in gratia incrementum.}

\emph{Ques.} 97. \emph{What is required to the worthy receiving of the Lord's Supper?}

Quæs. \emph{Ut digne quis participet cœnam Dominicam quid requiritur?}

\emph{Ans.} It is required of them that would worthily partake of the Lord's Supper, that they examine

Resp. \emph{Qui cœnam Dominicam digne cupiunt participare, requiritur, ut semet examinent cum de cognitione }

\emph{themselves of their knowledge to discern the Lord's body, of their faith to feed upon him, of their repentance, love, and new obedience; lest coming unworthily, they eat and drink judgment to themselves.}

\emph{sua, qua corpus Domini valeant discernere, tum de fide sua, qua vescantur ipso, tum etiam de resipiscentia sua, amore ac obedientia nova; ne forte indigni si advenerint, judicium edant bibantque sibimetipsis.}

\emph{Ques.} 98. \emph{What is prayer?}

Quæs. \emph{Quid est precatio?}

\emph{Ans.} Prayer is an offering up of our desires unto God, for things agreeable to his will, in the name of Christ, with confession of our sins, and thankful acknowledgment of his mercies.

Resp. \emph{Precatio est qua petitiones nostras pro rebus divinæ voluntati congruis offerimus Deo, in nomine Christi, una cum peccatorum nostrorum confessione, et grata beneficiorum ejus agnitione.}

\emph{Ques.} 99. \emph{What rule hath God given for our direction in prayer?}

Quæs. \emph{Quam nobis regulum præscripsit Deus precibus nostris dirigendis?}

\emph{Ans.} The whole Word of God is of use to direct us in prayer, but the special rule of direction is that form of prayer which Christ taught his disciples, commonly called, \emph{The Lord's Prayer.}

Resp. \emph{Totum Dei verbum utile est nobis in oratione dirigendis; specialis vero directionis norma est illa orationis formula quam discipulos suos edocuit Christus, }oratio dominica \emph{quæ vulgo dicitur.}

\emph{Ques.} 100. \emph{What doth the preface of the Lord's Prayer teach us?}

Quæs. \emph{Quid nos docet orationis Dominicæ præfatio?}

\emph{Ans.} The preface of the Lord's Prayer, which is, '\emph{Our Father which art in heaven,}' teacheth us to draw near to God with all holy reverence and confidence, as children to a father, able\footnote{London ed. of 1658 omits \emph{able and}.} and ready to help us; and that we should pray with and for others.

Resp. \emph{Orationis Dominicæ præfatio nempe} [Pater noster, qui es in cœlis] \emph{nos docet accedere ad Deum cum omni sancta reverentia ac confidentia, tanquam filios ad patrem, qui et potis est ut paratus nobis opitulari; prout etiam cum aliis atque pro aliis orare.}

\emph{Ques.} 101. \emph{What do we pray for in the first petition?}

Quæs. \emph{Quid est quod oramus in petitione prima?}

\emph{Ans.} In the first petition, which

Resp. \emph{In petitione prima, scil. }

is, '\emph{Hallowed be thy name,}' we pray that God would enable us and others to glorify him in all that whereby he maketh himself known, and that he would dispose all things to his own glory.

[Sanctificetur nomen tuum] \emph{oramus et efficere velit Deus, ut eum nos aliique, in eis, quibuscunque se notum nobis facit, glorificare valeamus; atque ad suam ipsius gloriam omnia dirigere velit ac disponere.}

\emph{Ques.} 102. \emph{What do we pray for in the second petition?}

Quæs. \emph{Quid petimus in secunda petitione?}

\emph{Ans.} In the second petition, which is, '\emph{Thy kingdom, come},' we pray that Satan's kingdom may be destroyed, and that the kingdom of grace may\footnote{London ed. of 1658: \emph{might}.} be advanced, ourselves and others brought into it, and kept in it, and that the kingdom of glory may be hastened.

Resp. \emph{In petitione secunda, quæ hujusmodi est} [adveniat regnum tuum] \emph{petimus ut destruatur regnum Satanæ, gratiæ vero regnum ut promoveatur, ut nos aliique in eo simus cum constituti tum conservati ne excidamus, utque regnum gloriæ velit Deus adproperare.}

\emph{Ques.} 103. \emph{What do we pray for in the third petition?}

Quæs. \emph{In petitione tertia quid precamur?}

\emph{Ans.} In the third petition, which is, '\emph{Thy will be done on earth as it is in heaven},' we pray that God by his grace would make us able and willing to know, obey, and submit to his will in all things, as the angels do in heaven.

Resp. \emph{In petitione tertia, scil. hisce verbis} [fiat voluntas tua in terris sicut in cœlis] \emph{precamur efficere velit Deus, ut nos per gratiam voluntatem ejus tum cognoscere, tum ei in omnibus obtemperare, et nos submittere, id quod in cœlis faciunt Angeli, et valeamus et velimus.}

\emph{Ques.} 104. \emph{What do we pray for in the fourth petition?}

Quæs. \emph{Quid oramus in petitione quarta?}

\emph{Ans.} In the fourth petition, which is, ',' we pray that of God's free gift we may receive a competent portion of the good things of this life, and enjoy his blessing with them.

Resp. \emph{In quarta petitione quæ sic habetur} [Panem nostrum quotidianum da nobis hodie] \emph{oramus ut e donatione Dei gratuita, bonorum quæ hujus vitæ sunt portionem idoneam obtineamus, ejusque una cum iis benedictione perfruamur.}

\emph{Ques.} 105. \emph{What do we pray for in the fifth petition?}

Quæs. \emph{Quid precamur in petitions quinta?}

\emph{Ans.} In the fifth petition, which is, '\emph{And forgive us our debts as we forgive our debtors},' we pray that God, for Christ's sake, would freely pardon all our sins; which we are the rather encouraged to ask, because by his grace we are enabled from the heart to forgive others.

Resp. \emph{In petitione quinta, cujus verba sunt} [Ac remitte nobis debita nostra, sic ut remittimus debitoribus nostris] \emph{precamur ut Deus peccata nostra omnia gratis velit propter Christum condonare, quod quidem ut petamus eo magis animus nobis fit, quod aliis animitus condonare gratia ipsius auxiliante valeamus.}

\emph{Ques.} 106. \emph{What do we pray for in the sixth petition?}

Quæs. \emph{Quid petimus in sexta petitione?}

\emph{Ans.} In the sixth petition, which is, '\emph{And lead us not into temptation, but deliver us from evil},' we pray that God would either keep us from being tempted to sin, or support and deliver us when we are tempted.

Resp. \emph{In petitione sexta, quam hæc verba complectuntur} [Et ne nos inducas in tentationem, sed libera nos a malo] \emph{oramus ut velit nos Deus aut immunes a tentatione ad peccatum conservare, aut certe tentatos suffulcire ac liberare.}

\emph{Ques.} 107. \emph{What doth the conclusion of the Lord's Prayer teach us?}

Quæs. \emph{Quid nos docet orationis Dominicæ conclusio?}

\emph{Ans.} The conclusion of the Lord's Prayer, which is, '\emph{For thine is the kingdom, and the power and the glory forever, Amen,}' teacheth us to take our encouragement in prayer from God only, and in our prayers to praise him; ascribing kingdom, power, and glory to him; and in testimony of our desire and assurance to be heard, we say, \emph{Amen}.

Resp. \emph{Orationis Dominicæ conclusio} [Quia tuum est regnum, potentia et gloria, in secula, Amen] \emph{Nos docet animos ac confidentiam nobis in orando a solo Deo derivare, eumque in precibus nostris laudare, regnum ei, potentiam, ac gloriam tribuendo; quoque desiderium nostrum testemur, et exauditionis confidentiam, dicimus, Amen.}

T\textsc{he} T\textsc{en} C\textsc{ommandments}.

D\textsc{ecalogus}.

E\textsc{xodus xx}.

E\textsc{xod. xx}.

God spake all these words, saying, I am the Lord thy God, which have brought thee out of the land of Egypt, out of the house of bondage.

\emph{Locutus est Deus omnia hæc verba, dicendo; Ego sum Dominus Deus tuus, qui te eduxi e terra Ægypti, e domo servitutis.}

I. Thou shalt have no other gods before me.

I. \emph{Non habebis deos alios coram me.}

II. Thou shalt not make unto thee any graven image, or any likeness of any thing that is in heaven above, or that is in the earth beneath, or that is in the water under the earth: thou shalt not bow down thyself to them, nor serve them: for I the Lord thy God am a jealous God, visiting the iniquity of the fathers upon the children unto the third and fourth generation of them that hate me; and showing mercy unto thousands of them that love me and keep my commandments.

II. \emph{Non facies tibi imaginem quamvis sculptilem, aut similitudinem rei cujusvis quæ est in cœlis superne, aut inferius in terris, aut in aquis infra terram; non incurvabis te iis, nec eis servies: siquidem ego Dominus Deus tuus Deus sum zelotypus, visitans iniquitates patrum in filios ad tertiam usque quartamque progeniem osorum mei, exhibens vero misericordiam ad millenas usque diligentium me ac mandata mea observantium.}

III. Thou shalt not take the name of the Lord thy God in vain: for the Lord will not hold him guiltless that taketh his name in vain.

III. \emph{Nomen Domini Dei tui inaniter non usurpabis; non enim eum pro insonte habebit Dominus qui nomen ejus inaniter adhibuerit.}

IV. Remember the Sabbath-day to keep it holy. Six days shalt thou labor, and do all thy work; but the seventh day is the Sabbath of the Lord thy God; in it thou shalt not do any work, thon, nor thy son: nor thy daughter, thy man-servant, nor thy maid-servant, nor thy cattle, nor thy stranger that is

IV. \emph{Memineris diem Sabbati ut sanctifices eum; sex diebus operaberis et facies omne opus tuum, septimus vero dies sabbatum est Domini Dei tui, opus in eo nullum facies tu, neque servus tuus, nec ancilla tua, neque jumentum tuum, neque hospes tuus quicunque intra portas tuas commoratur: Nam sex diebus perfecit }

within thy gates; for in six days the Lord made heaven and earth, the sea, and all that in them is, and rested the seventh day; wherefore the Lord blessed the Sabbath-day, and hallowed it.

\emph{Dominus cœlum, terramque, mare et quicquid in illis continetur: septimo vero die requievit; quamobrem benedixit Dominus diei Sabbati, eumque sanctificavit.}

V. Honor thy father and thy mother; that thy days may be long upon the land which the Lord thy God giveth thee.

V. \emph{Honora patrem tuum ac matrem tuam, ut prolongentur dies tui in terra illa quam tibi largitur Dominus Deus tuus.}

VI. Thou shalt not kill.

VI. \emph{Non occides.}

VII. Thou shalt not commit adultery.

VII. \emph{Non mæchaberis.}

VIII. Thou shalt not steal.

VIII. \emph{Non furaberis.}

IX. Thou shalt not bear false witness against thy neighbor.

IX. \emph{Non eris adversus proximum tuum testis mendax.}

X. Thou shalt not covet thy neighbor's house, thou shalt not covet thy neighbor's wife, nor his man-servant, nor his maid-servant, nor his ox, nor his ass, nor any thing that is thy neighbor's.

X. \emph{Non concupisces proximi tui domum, non concupisces proximi tui uxorem, non servum, non ancillam, non bovem, non asinum neque aliud denique quicquam quod est proximi tui.}

T\textsc{he} L\textsc{ord's} P\textsc{rayer}.

O\textsc{ratio} D\textsc{ominica}.

M\textsc{att. vi}.

M\textsc{att. vi.}

Our Father which art in heaven, hallowed be thy name. Thy kingdom come. Thy will be done in earth as it is in heaven. Give us this day our daily bread. And forgive us our debts, as we forgive our debtors. And lead us not into temptation, but deliver us from evil. For thine is the kingdom, and the power, and the glory, forever. Amen.

\emph{Pater noster qui es in cœlis, sanctificetur nomen tuum, adveniat regnum tuum, fiat voluntas tua in terris sicut in cœlis, panem nostrum quotidianum da nobis hodie, ac remitte nobis debita nostra, sicut nos remittimus debitoribus nostris, et ne nos inducas in tentationem, sed libera nos a malo, quia tuum est regnum, potentia et gloria in secula. Amen.}

T\textsc{he} C\textsc{reed}.

S\textsc{ymbolum.}

I believe in God the Father almighty, maker of heaven and earth; and in Jesus Christ his only Son, our Lord; who\footnote{London ed. of 1658: \emph{which.}} was conceived by the Holy Ghost, born of the Virgin Mary; suffered under Pontius Pilate, was crucified, dead, and buried; he descended into hell:\footnote{\emph{i. e.}, Continued in the state of the dead, and under the power of death, until the third day.} the third day he rose again from the dead; he ascended into heaven, and sitteth on the right hand of God the Father almighty; from thence he shall come to judge the quick and the dead. I believe in the Holy Ghost; the holy catholic church; the communion of saints; the forgiveness of sins; the resurrection of the body; and the life everlasting. Amen.

\emph{Credo in Deum Patrem omnipotentem, creatorem cœli ac terræ; et in Jesum Christum filium ejus unicum, Dominum nostrum; qui conceptus est e Spiritu sancto, natus ex Maria Virgine; passus sub Pontio Pilato, crucifixus, mortuus et sepultus; descendit ad inferos:}\footnote{i. e., \emph{Permansit in statu mortuorum et sub potestate mortis usque ad diem tertium.}} \emph{ tertio die resurrexit a mortuis: ascendit in cœlum, et sedet ad dextram Dei patris omnipotentis: unde venturus est ad judicandum vivos et mortuos. Credo in Spiritum Sanctum: Sanctam ecclesiam catholicam: Sanctorum communionem: remissionem peccatorum, resurrectionem corporis et vitam æternam. Amen.}

The oldest editions of the Westminster Shorter Catechism have the following addendum:

So much of every Question both in the Larger and Shorter Catechism, is repeated in the Answer, as maketh every Answer an entire Proposition, or sentence in itself: to the end the Learner may further improve it upon all occasions, for his increase in knowledge and piety, even out of the course of catechising, as well as in it.

\emph{E quæstione qualibet utriusque catechismi repetitum dedimus in responsione quantum responsionem quamlibet reddat propositionem integram, sive sententiam absolutam. Eo nempe consilio ut discenti ulterius utilis esse possit, quoties occasio tulerit, ad cognitionis ac pietatis incrementum, vel extra catechisandi rationem.}

And albeit the substance of the doctrine comprised in that Abridgment commonly called, \emph{The Apostles' Creed}, be fully set forth in each of the Catechisms, so as there is no necessity of inserting the Creed itself, yet it is here annexed, not as though it were composed by the Apostles, or ought to be esteemed Canonical Scripture, as the Ten Commandments, and the Lord's Prayer (much less a Prayer, as ignorant people have been apt to make both it and the Decalogue), but because it is a brief sum of the Christian faith, agreeable to the Word of God, and anciently received in the Churches of Christ.

\emph{Et quamvis in alterutro Catechismo substantia doctrinæ in compendia illo \{Symbolo apostolico vulgo dicto) comprehensæ plene ac perfecte exhibeatur, adeo quidem ut nulla supersit necessitas symbolum ipsum inserendi: nihilominus tamen hic illud subnectendum esse duximus; non perinde quasi aut ab ipsis Apostolis fuerit concinnatum, aut pariter cum decalogo, ac oratione Dominica pro Scriptura canonica haberi debeat:} (\emph{nedum certe pro oratione, quo nomine ignara plebecula cum illud tum decalogum in proclivi fuit ut usurparet}), \emph{verum quod sit fidei Christianæ breve compendium, verbo Dei consentaneum, ac in Ecclesiis Christi antiquitus receptum.}

C\textsc{ornelius} B\textsc{urges}, \emph{Prolocutor pro tempore}.

H\textsc{enry} R\textsc{oborough}, \emph{Scriba.}

A\textsc{doniram} B\textsc{yfield}, \emph{Scriba.}
